% =============================================================================
% Conceptual Topology Mapping Benchmark
% Jameson Hodge, 2025-2026
% =============================================================================
\documentclass[11pt]{article}

% --- Page geometry ---
\usepackage[margin=1in]{geometry}

% --- Math ---
\usepackage{amsmath}
\usepackage{amssymb}

% --- Graphics ---
\usepackage{graphicx}

% --- Tables ---
\usepackage{booktabs}
\usepackage{longtable}
\usepackage{array}
\usepackage{calc}       % for \real{} in pandoc longtables
\usepackage{multirow}

% --- Lists ---
\usepackage{enumitem}

% --- Citations ---
\usepackage[round]{natbib}
\bibliographystyle{plainnat}

% --- Hyperlinks ---
\usepackage[colorlinks=true,linkcolor=blue,citecolor=blue,urlcolor=blue]{hyperref}
\usepackage{url}

% --- Font encoding ---
\usepackage[T1]{fontenc}

% --- Unicode character mappings (for pdflatex) ---
\usepackage{newunicodechar}
\newunicodechar{→}{$\rightarrow$}
\newunicodechar{−}{$-$}
\newunicodechar{×}{$\times$}
\newunicodechar{≈}{$\approx$}
\newunicodechar{≤}{$\leq$}
\newunicodechar{≥}{$\geq$}
\newunicodechar{∞}{$\infty$}
\newunicodechar{∪}{$\cup$}
\newunicodechar{η}{$\eta$}
\newunicodechar{α}{$\alpha$}
\newunicodechar{Σ}{$\Sigma$}
\newunicodechar{Δ}{$\Delta$}
\newunicodechar{²}{$^2$}
\newunicodechar{ᵢ}{$_i$}
\newunicodechar{₂}{$_2$}

% --- Pandoc compatibility ---
\providecommand{\tightlist}{%
  \setlength{\itemsep}{0pt}\setlength{\parskip}{0pt}}

% --- Custom commands ---
\newcommand{\eg}{e.g.\@}
\newcommand{\ie}{i.e.\@}

% =============================================================================
\title{Conceptual Topology Mapping Benchmark:\\
  Waypoint Elicitation Reveals Structured, Asymmetric,\\
  Model-Specific Navigation in LLM Conceptual Space}

\author{Jameson Hodge\\
  \texttt{jameson@jamesonhodge.com}\\
  \url{https://jamesonhodge.com}}

\date{March 2026}

% =============================================================================
\begin{document}

\maketitle

% --- Abstract ---
\begin{abstract}
Static geometric structure in LLM representations is increasingly well-characterized: linear directions encode concepts, hierarchies form polytopes, and translation symmetry in data statistics determines manifold shape. But no work has tested whether this geometry supports reliable \emph{behavioral traversal} --- whether models can consistently navigate between concepts, not merely represent them. The distinction between map and route is logically independent, and static structure alone does not entail navigational consistency.

We introduce the Conceptual Topology Mapping Benchmark, a waypoint-elicitation paradigm that probes navigational geometry across approximately 21,540 runs, 12 models from 11 independent training pipelines, and 11 experimental phases. Models are given concept pairs and asked to produce intermediate waypoints; the resulting paths are evaluated against metric axioms, compositional structure, and causal perturbation, yielding six robust empirical claims plus one qualified.

Models exhibit distinct ``conceptual gaits'' --- navigational signatures with consistency spanning a 2.9x range (0.258 to 0.747 across 12 models) --- on a shared underlying geometry. Navigation is fundamentally asymmetric (mean directional asymmetry 0.811), consistent with quasimetric rather than metric structure. Bridge concepts function as structural bottlenecks, not mere associations, with frequencies exceeding 0.97 under robustness testing. These properties generalize across all 12 models and resist protocol variation. Single-variable explanations consistently fail: seven follow-up hypotheses --- spanning mechanistic, model-deficit, methodological, and structural types --- were falsified and one resurrected with additional data, revealing a mechanism ceiling where qualitative structure is characterizable but quantitative dynamics remain opaque to current tools. Twenty-nine documented dead ends provide transparent accounting of explanatory limits. The benchmark establishes that conceptual navigation has geometric regularity that can be reliably detected but cannot yet be mechanistically decomposed.
\end{abstract}

% =============================================================================
% Main text
% =============================================================================

\section{Introduction}\label{introduction}

\subsection{The Navigation Gap}\label{the-navigation-gap}

LLM embedding spaces have rich geometric structure. Linear directions encode concepts \citep{park2024linear}. Hierarchies form polytopes \citep{park2025geometry}. Hyperbolic geometry captures tree-like relations \citep{nickel2017poincare}. Most fundamentally, \citet{karkada2025translation} proved that this geometry arises inevitably from translation symmetry in data statistics --- the manifold shape is determined by co-occurrence structure, not architectural choices.

The static geometry of LLM representations is increasingly well-characterized, and theoretically grounded. What remains uncharacterized is the \emph{dynamic} behavior on that geometry.

But geometry is not navigation.

A map is not a route.

Static analysis tells us what landmarks look like; it does not tell us whether models can reliably walk between them. Embedding probes reveal that ``warm'' sits between ``hot'' and ``cold'' in representational space, but they do not test whether a model traversing from hot to cold will actually pass through warm --- or whether it will take a different route entirely, through ``tepid'' and ``chilly'' and ``brisk.'' The distinction matters: two agents can share an identical map and prefer entirely different paths.

The gap is not merely philosophical. Static geometry answers questions about structure --- what is near what, which directions encode which features, what the curvature looks like. Navigation answers questions about behavior --- which routes models actually take, whether those routes are consistent across repeated traversals, whether they differ by direction, and whether they compose when chained through intermediate concepts. A model might have perfect geometric representations and still navigate erratically, or conversely, navigate with rigid consistency through a geometry that embedding analysis would not predict. The static and dynamic questions are logically independent, and neither entails the other.

This paper asks the dynamic question directly: \textbf{Do LLMs have consistent, measurable geometric structure in how they \emph{navigate} between concepts?}

Not what their representations look like, but how they behaviorally traverse conceptual space when asked to produce intermediate waypoints between two endpoints.

The answer, across approximately 21,540 runs, 12 models from 11 independent training pipelines, and 11 experimental phases, is yes --- with substantial qualifications about what that structure can and cannot explain. Navigation is structured (gait consistency spanning a 2.9x range across models), asymmetric (mean directional asymmetry 0.811), organized around bottleneck landmarks (bridge frequencies exceeding 0.97 under robustness testing), and consistent with quasimetric topology (triangle inequality holding at approximately 91\%). But 7 of 8 follow-up hypotheses were falsified (one later resurrected), and 29 documented dead ends attest to the gap between characterizing structure and explaining it.

\textbf{Theoretical complement.} \citet{karkada2025translation} proved \emph{why} static geometry exists: translation symmetry in co-occurrence statistics determines manifold shape. We test \emph{whether that geometry supports behavioral navigation} --- whether the geometric structure that arises inevitably from data statistics produces consistent, structured, asymmetric paths when models are asked to walk between concepts. Their contribution explains the map; ours probes the routes. The two contributions are complementary: static geometry is necessary for navigation to be structured, but not sufficient --- the dynamics could still be noisy, inconsistent, or model-independent. They are none of these.

\subsection{Why Navigation Matters}\label{why-navigation-matters}

Four independent lines of research motivate the navigation question.

\textbf{Cognitive science predicts navigable conceptual spaces.} \citet{gardenfors2000} theorized that concepts live in geometric spaces with measurable distance --- spaces that support not just representation but traversal. \citet{constantinescu2016grid} provided neural evidence: grid cells in the human entorhinal cortex, originally characterized for spatial navigation, fire during \emph{conceptual} navigation, suggesting that spatial mechanisms have been repurposed for abstract thought. If human conceptual spaces are navigable, the question of whether LLM conceptual spaces share this property is natural and testable.

Our benchmark provides a large-scale computational probe: 12 models, 100+ concept pairs, and over 21,000 runs testing the navigability prediction directly.

\textbf{Representational convergence does not imply navigational convergence.} The Platonic Representation Hypothesis \citep{huh2024platonic} claims that models trained on similar data converge toward shared representations, supported by measures like CKA alignment. But representational alignment is not navigational alignment. Two agents can agree perfectly on a map --- identical embeddings, identical similarity structures --- and still prefer different routes.

Our benchmark tests the behavioral complement: do models that share representations also share navigational behavior? The answer requires distinguishing three levels. Navigational \emph{structure} converges --- all 12 models exhibit gait consistency, directional asymmetry, and bridge bottleneck topology. Navigational \emph{content} converges among frontier-scale models --- the same bridge concepts (``spectrum,'' ``dog,'' ``warm'') appear across providers, with aggregate bridge frequency differences whose confidence intervals include zero. But navigational \emph{dynamics} diversify: the 2.9x gait range from GPT (0.258) to Mistral (0.747) does not collapse with scale.

Same map, different routes. The Platonic Representation Hypothesis receives partial confirmation with a clear boundary: structure and content converge, but gaits remain model-specific.

\textbf{The reversal curse reveals directional structure.} \citet{berglund2024reversal} demonstrated that autoregressive models cannot reverse learned associations: training on ``A is B'' does not teach ``B is A.'' This has been framed as a failure mode. We reframe it as a structural signal.

\citet{tversky1977features} established decades earlier that human similarity judgments are inherently asymmetric --- North Korea is judged more similar to China than China is to North Korea. Our benchmark finds pervasive directional asymmetry in LLM navigation (mean 0.811 in the Phase 2 cohort, all 12 tested models exceeding the 0.60 threshold), consistent with quasimetric topology. Forward and reverse paths between the same two concepts share less than 19\% of their waypoints on average. The asymmetry is not noise --- it tracks semantic properties of the concept pairs and model-specific navigational styles.

The connection across these three findings is thematic: all concern directional asymmetry in conceptual structure. We do not claim that the reversal curse mechanistically causes our observed asymmetry; rather, both phenomena reflect a shared property of directed conceptual spaces in which the path from A to B traverses fundamentally different intermediate territory than the path from B to A.

\textbf{Multi-hop reasoning failures call for richer diagnostics.} The compositionality gap \citep{press2022compositionality} and the two-hop curse~\citep{twohopcurse2025} document that models struggle with latent multi-step reasoning. These findings use binary pass/fail evaluation: the model either retrieves the composed fact or it does not.

Our benchmark generalizes the diagnostic. Rather than asking whether a model can compose two facts, we ask what \emph{shape} the composition takes --- which intermediate concepts appear, how consistently they appear across repeated runs, whether they compose properly (the triangle inequality holds at approximately 91\%), and whether they differ by direction. This produces a geometric characterization of compositional structure rather than a scalar accuracy measure.

The characterization is informative even when the model ``fails.'' Asymmetry and bridge fragility are structured phenomena, not random noise. A model that produces different paths in opposite directions is not failing at composition --- it is revealing the directional structure of its conceptual space.

\subsection{Contributions}\label{contributions}

This paper makes six contributions, spanning paradigm design, empirical findings, and methodological transparency:

\begin{enumerate}
\def\labelenumi{\arabic{enumi}.}
\item
  \textbf{A novel evaluation paradigm} --- waypoint elicitation as behavioral probe for conceptual geometry, distinct from static embedding analysis, factual QA, and free association. The paradigm is self-grounding: model outputs generate their own ground truth, enabling scalable evaluation without human annotation. This introduces a circularity concern addressed through external mathematical tests (metric axioms), cross-model validation (12 independent models), and causal intervention (pre-fill experiments).
\item
  \textbf{Six robust empirical claims} about navigational structure, replicated across multiple phases and models, plus one qualified claim (R5, control validation) that reveals a fundamental paradigm limitation (Section 9.5):

  \begin{itemize}
  \tightlist
  \item
    Models have characteristic conceptual gaits spanning a 2.9x range {[}robust{]} R1.
  \item
    Navigation is consistent with quasimetric structure: systematically asymmetric (mean 0.811) with the triangle inequality holding at approximately 91\% {[}robust{]} R2.
  \item
    Polysemy sense differentiation is genuine, with cross-pair Jaccard of 0.011-0.062 for different senses {[}robust{]} R3.
  \item
    Hierarchical paths are compositional with 4.9x higher transitivity than random controls {[}robust{]} R4.
  \item
    Control validation supports the paradigm, though stapler-monsoon fails strict criteria for all 12 models {[}robust, qualified{]} R5.
  \item
    Bridge concepts function as structural bottlenecks with frequencies exceeding 0.97 under robustness testing {[}robust{]} R6.
  \item
    Cue-strength gradient exists: 12/16 family-model combinations show monotonic decrease {[}robust{]} R7.
  \item
    R3, R4, and R7 are untested beyond the original 4-model cohort.
  \end{itemize}
\item
  \textbf{A mechanism ceiling result} --- 7 follow-up hypotheses fail across Phases 8--10 --- spanning mechanistic, model-deficit, methodological, and effect-direction types --- with one additional hypothesis initially falsified then resurrected with more data, establishing that conceptual navigation resists simple explanations (Section 7). Prediction accuracy descended from 81.3\% for characterization questions to 20--24\% for mechanistic ones --- tracking a systematic shift in prediction type, not random degradation.
\item
  \textbf{Cross-architecture generality} --- navigational structure and content generalize across 12 models from 11 independent training pipelines; scale differentiates (Section 8). Llama 3.1 8B retains gait and asymmetry but navigates through different landmarks (bridge frequency 0.200 vs 0.717--0.817 for frontier models). The within-family Llama comparison (Maverick 0.724 vs 8B 0.200) provides the sharpest evidence.
\item
  \textbf{Protocol robustness} --- bridge frequency insensitive to waypoint count and temperature (exceeding 0.97 across all conditions); gait rankings largely stable (W = 0.840); model identity dominates protocol variation (eta-squared = 0.242 vs approximately 0 for protocol factors) (Section 9). Asymmetry is resolution-dependent --- a measurement sensitivity finding that qualifies R2 without undermining it.
\item
  \textbf{29 documented dead ends} as honest accounting of the benchmark's learning curve. Failed predictions --- ``metaphor'' at frequency 0.00 for language-to-thought, gait normalization producing 0.000 improvement, three successive Gemini deficit characterizations all falsified --- contain more methodological information than several of the successful findings. The graveyard is a methodological contribution, not a failure record.
\end{enumerate}

\subsection{Paper Structure}\label{paper-structure}

The paper follows a six-act narrative structure, tracing the benchmark from initial characterization through mechanism and generality to robustness.

\textbf{Section 2: The Benchmark.} Defines the waypoint elicitation paradigm, derived metrics (gait consistency, path asymmetry, bridge frequency, navigational distance, triangle inequality), concept pair selection across 9 categories, the 12-model cohort accessed via OpenRouter, and the evidential strategy including self-grounding circularity mitigations.

\textbf{Section 3: Related Work.} Positions the benchmark against existing work on LLM geometry, conceptual spaces, compositionality, and evaluation methodology.

\textbf{Section 4: Act I -- Structure.} Establishes that models have distinct, stable conceptual gaits spanning a 2.9x range; paths are self-consistent across time (no temporal drift) and resolution (5-waypoint paths appear as subsequences of 10-waypoint paths 67.9\% of the time); both endpoints shape the trajectory via a dual-anchor effect; and polysemy sense differentiation (cross-pair Jaccard near zero for ``bank,'' ``bat,'' ``crane'') confirms that paths reflect conceptual structure rather than surface word statistics.

\textbf{Section 5: Act II -- Topology.} Characterizes navigational geometry as consistent with quasimetric structure: systematically asymmetric (mean 0.811), with the triangle inequality holding at approximately 91\% across three independent samples, compositional hierarchy at 4.9x over random controls, bottleneck bridge topology where navigational salience depends on directional information content rather than associative strength, and non-uniform navigational traffic concentrated at preferred waypoints.

\textbf{Section 6: Act III -- Mechanism.} Pre-filling the first waypoint causally displaces bridges (mean displacement 0.515, CI excluding zero), revealing associative primacy as the primary mechanism with secondary content modulation. The relation class of the pre-fill systematically affects bridge survival (Friedman p = 0.034), and the effect is bidirectional --- displacement for dominant bridges, facilitation up to 8x for marginal bridges under semantically aligned pre-fills.

\textbf{Section 7: Act IV -- Limits.} Seven mechanistic hypotheses fail across three phases; prediction accuracy descends from 81.3\% to 20\%; cross-model distance is fundamentally model-dependent (r = 0.170, gait normalization produces 0.000 improvement); the Gemini mystery survives three successive characterization attempts; and the mechanism ceiling establishes that navigation resists single-variable explanation.

\textbf{Section 8: Act V -- Generality.} Gait and asymmetry replicate universally across 12 models from 11 independent training pipelines. Bridge content generalizes among frontier-scale models but not to Llama 3.1 8B (bridge frequency 0.200 vs 0.717--0.817). The structure/content/scale hierarchy emerges as the organizing framework for cross-architecture comparison.

\textbf{Section 9: Act VI -- Robustness.} Model identity dominates protocol variation (eta-squared = 0.242 vs approximately 0 for protocol factors). Bridge frequency is the most robust property, exceeding 0.97 across all tested conditions. Asymmetry is resolution-dependent: 9-waypoint conditions clear the 0.60 threshold; 5-waypoint conditions fall just short. The control pair problem --- 0 of 24 replacement cells passing strict criteria --- is real and acknowledged directly.

\textbf{Section 10: Discussion.} Synthesizes the six-act narrative, states nine limitations directly (six benchmark design, three phenomenon), connects findings to Karkada et al.'s geometric inevitability, Gardenfors's conceptual spaces, the Platonic Representation Hypothesis, and spectral semantic attractors \citep{wyss2025spectral}, and identifies ten future directions from multi-variable mechanism models to human comparison.

\textbf{Section 11: Conclusion.} Distills the central empirical contribution: LLMs navigate conceptual space with structured, asymmetric, model-specific geometry that generalizes across architectures and resists mechanistic decomposition --- structure is detectable before it is explainable.
     % Section 1: Introduction
\input{02-benchmark}        % Section 2: The Benchmark
\section{Related Work}\label{related-work}

The benchmark sits at the intersection of seven distinct literatures: static geometry of LLM representations, the theoretical origin of that geometry, conceptual navigation and interpolation, consistency and directionality, multi-hop reasoning, cross-model comparison, and cognitive science foundations. Each has produced results that inform our design and findings, but none asks the question we ask: whether models can \emph{navigate} conceptual space consistently and whether the resulting paths have measurable geometric structure. We review each literature and state precisely what gap the benchmark fills.

\subsection{Static Geometry of LLM Representations}\label{static-geometry-of-llm-representations}

The linear representation hypothesis \citep{park2024linear} established that individual concepts correspond to linear directions in activation space. Its extension \citep{park2025geometry} showed that hierarchical relationships form polytope structures --- convex regions whose geometric relationships encode taxonomic nesting. Separately, hyperbolic embedding methods \citep{nickel2017poincare} demonstrated that tree-like hierarchies are better captured in negatively curved spaces, a finding extended to LLM representations by recent work on HELM~\citep{helm2025} and direct measurement of delta-hyperbolicity and ultrametricity~\citep{deltahyperbolicity2025}, which quantify how tree-like a given embedding space actually is.

Topological data analysis has added finer-grained structural characterization. Explainable Mapper~\citep{explainablemapper2025} identified circular manifolds for temporal concepts and linear manifolds for ordered quantities in LLM representations. Persistent Topological Features~\citep{persistenttopological2025} used persistence diagrams to characterize the homological structure of concept clusters. \citet{wyss2025spectral} proved mathematically that LLM activation manifolds partition into discrete basins of attraction via spectral analysis of semantic attractors --- a result directly relevant to the bottleneck bridge structure we observe empirically (Section 5.4). \citet{joshi2025dimensionality} characterized intrinsic dimensionality profiles across transformer layers, finding a universal early-expand-late-compress pattern that holds across architectures.

\textbf{Gap we fill.} This entire literature characterizes \emph{what shape the space has} --- its curvature, dimensionality, manifold topology, and basin structure. All of it is static. We probe whether models can \emph{navigate} that space consistently, producing paths with measurable geometric properties. The distinction is between a map and a route: static geometry gives us the map; our benchmark tests whether models can walk it.

\subsection{Theoretical Origin of Geometric Structure}\label{theoretical-origin-of-geometric-structure}

\citet{karkada2025translation} proved that translation symmetry in word co-occurrence statistics mathematically determines manifold geometry in embedding spaces. Their result is architecture-independent: the same geometric structures arise in Word2Vec, GloVe, and Gemma because all train on data distributions with shared statistical regularities. Geometry comes from data, not architecture.

This theoretical result connects directly to our empirical findings. If geometric structure is determined by data statistics that all models share, then navigational structure derived from that geometry should also generalize across architectures --- which is what we observe (R1 and R2 universal across 12 models from 11 training pipelines; Section 8). The scale effect (Llama 8B as outlier in bridge frequency) is consistent with their framework: if extracting the full geometric structure requires sufficient capacity to represent translation symmetries, smaller models may have the coarse structure but lack the resolution to navigate through the same landmarks as frontier models.

\textbf{Gap we fill.} Karkada et al.~study the static manifold shape and its mathematical origin. We test the dynamic complement: whether that geometry supports consistent behavioral navigation. Their theory is consistent with our empirical finding of cross-architecture generalization; our data provides the behavioral evidence their framework implies but does not test. The contributions are complementary --- they explain why the geometry exists; we show what it supports.

\subsection{Conceptual Navigation and Interpolation}\label{conceptual-navigation-and-interpolation}

The mathematical framework for navigation through latent spaces exists but has not been applied to LLM behavioral outputs. \citet{arvanitidis2018latent} developed Riemannian geodesics as semantically meaningful paths in generative model latent spaces --- shortest paths through curved space that respect the local geometry. The framework is directly applicable to conceptual navigation but has been used for image generation and VAE interpolation, not for probing how language models traverse between concepts.

The closest precedent for waypoint-based navigation is Relation Embedding Chains~\citep{relationchains2023}, which discovered intermediate concepts X such that A-to-B can be decomposed as A-to-X-to-B using knowledge graph paths. This is the nearest prior work to our waypoint elicitation paradigm, though it operates on static knowledge graph structure rather than behavioral generation. The Geometry of Knowledge~\citep{geometryknowledge2025} traversed semantic manifolds for generative diversity --- the closest existing work to ``conceptual traversal'' --- but used traversal as a tool for improving generation, not as an evaluation paradigm.

More recently, reasoning-as-path-search has entered the mainstream. RiemannInfer~\citep{riemanninfer2026} frames multi-step inference as trajectory optimization on a Riemannian manifold, and Geometric Reasoner~\citep{geometricreasoner2026} uses geodesic metaphors for chain-of-thought planning. These works treat the path metaphor as productive for reasoning but do not test the geometric properties of the paths themselves.

\textbf{Gap we fill.} No prior work asks models to produce intermediate waypoints between two concepts and then tests whether the resulting paths have consistent geometric structure --- satisfying metric axioms, exhibiting compositional properties, and responding to causal perturbation. The mathematical framework exists (Riemannian geodesics); the precedent for intermediate concepts exists (relation embedding chains); the path metaphor is becoming mainstream (reasoning-as-search). What is missing is the empirical test of whether behaviorally generated paths have measurable, replicable geometry. We provide that test.

\subsection{Consistency, Directionality, and the Reversal Curse}\label{consistency-directionality-and-the-reversal-curse}

\citet{berglund2024reversal} established the reversal curse: autoregressive models trained on ``A is B'' systematically fail to learn ``B is A.'' This is an architectural limitation --- the left-to-right generation order creates inherent directional bias. the directional optimization asymmetry literature~\citep{directionalasymmetry2025} deepened this finding by showing that inverse mappings incur higher excess loss even under symmetric training objectives, establishing that the A-to-B versus B-to-A gap is architectural, not merely a training-data artifact. \citet{barakat2026hysteresis} extended directionality to semantic caching, demonstrating that input ordering steers inference trajectories into different attractor basins --- directional hysteresis at the system level.

On the consistency side, recent work has documented pervasive self-consistency failures in frontier models~\citep{selfconsistency2025}, with high variability even in GPT-4 across repeated evaluations on identical prompts. Trajectory variance in agentic settings~\citep{trajectoryvariance2025} shows substantial action-sequence diversity across runs of the same agent on the same task, establishing that behavioral variability is not confined to simple generation tasks.

\textbf{Our contribution.} Existing work treats inconsistency and asymmetry as failure modes --- problems to be solved or architectural limitations to be documented. We reframe asymmetry as data. The pervasive directional asymmetry we observe (mean 0.811 in the Phase 2 cohort, exceeding 0.60 across all 12 models tested) is not a bug but a structural property of navigational space, consistent with quasimetric topology. The pattern of asymmetry reveals structure: it tracks semantic properties of concept pairs and model-specific navigational styles, producing systematic geometric signatures rather than noise.

A defensive clarification is warranted here. Reviewers will ask why this is not simply self-consistency work with a geometry gloss. The contrast is precise: self-consistency work measures the \emph{magnitude} of variance across samples and treats it as error to be minimized. We test whether the \emph{structure} of that variance satisfies formal geometric axioms --- metric properties, compositionality, causal sensitivity. We do not ask whether models are consistent; we ask whether the pattern of their inconsistency has geometric regularity. Variance is our data, not our failure mode.

\subsection{Multi-Hop Reasoning and Compositionality}\label{multi-hop-reasoning-and-compositionality}

The compositionality gap \citep{press2022compositionality} demonstrated that models can answer individual sub-questions correctly while failing the composed question --- knowing both ``A is B'' and ``B is C'' but failing ``A is C.'' the two-hop curse~\citep{twohopcurse2025} formalized this for knowledge composition: models trained on A-to-B and B-to-C facts fail A-to-C inference without explicit chain-of-thought prompting. TWOHOPFACT~\citep{twohopfact2024} studied latent multi-hop reasoning through bridge entities, identifying the intermediate step as the critical point of failure --- the model must internally activate the bridge concept to compose the two hops.

\textbf{Our contribution.} We generalize the compositionality question from binary pass/fail to geometric characterization. Multi-hop reasoning work asks whether the model gets the composed answer \emph{correct}. We ask what \emph{shape} the intermediate sequence takes --- whether waypoints compose consistently (hierarchical transitivity 4.9x over random controls), whether bridge concepts function as structural bottlenecks (frequencies of 0.60-1.00 for predicted bridges), and whether the compositional structure satisfies the triangle inequality (\textasciitilde91\% across three independent samples). Waypoints externalize latent composition, making the intermediate steps available for geometric analysis rather than accuracy scoring.

A second defensive clarification applies here. The paradigm will invite comparison to chain-of-thought elicitation. The contrast is fundamental: chain-of-thought elicits reasoning steps toward a \emph{correct answer} --- intermediate steps are evaluated by whether they lead to a right response. We elicit navigational waypoints toward a \emph{destination concept} --- there is no correct answer, only geometric structure to characterize. Chain-of-thought is evaluated by task accuracy; our paradigm is evaluated by structural consistency, metric properties, and causal sensitivity. The two paradigms share surface similarity (both produce ordered intermediate steps) but differ in what they measure and how they are validated.

\subsection{Cross-Model Comparison and the Platonic Representation Hypothesis}\label{cross-model-comparison-and-the-platonic-representation-hypothesis}

The Platonic Representation Hypothesis \citep{huh2024platonic} predicts that model representations converge with scale, measured by CKA alignment between activation spaces. Larger models from different families produce increasingly similar internal representations --- a convergence toward a shared ``platonic'' geometry. Multi-way alignment work~\citep{multiwayalignment2026} extended this to shared reference spaces across N models, and model stitching experiments~\citep{modelstitching2025} demonstrated that features transfer meaningfully but non-uniformly across architectures.

On the behavioral side, LLM behavioral fingerprinting~\citep{behavioralfingerprint2025} established that models have characteristic behavioral signatures --- stable patterns of response that distinguish one model from another across tasks. \citet{nanda2026attractors} identified conversational attractor states: model-specific behavioral loops (Claude trending toward philosophical synthesis, GPT toward system-building) that function as the conversational analog of our quantitative gaits. Behavioral failure manifold mapping~\citep{failuremanifold2025} characterized model-specific topographies with qualitatively different basin structures, showing that models fail in geometrically distinct ways.

\textbf{Our contribution.} The PRH measures static representational alignment. We test \emph{dynamic navigational alignment} --- whether shared representations imply shared routes. The finding is a three-level hierarchy: structure generalizes (all 12 models have gait and asymmetry), content mostly generalizes among frontier-scale models (aggregate bridge frequency difference CI includes zero), but gaits diverge (the 2.9x gait range does not collapse with scale). Same map, different routes. This is a behavioral refinement of the PRH: representational convergence does not imply behavioral convergence. The models share structural properties (gait and asymmetry are universal) and largely converge on the strongest navigational landmarks, but they navigate with model-specific route preferences that remain distinctive even at frontier scale --- and their navigational distances are fundamentally model-dependent (Section 7.4).

\subsection{Cognitive Science Foundations}\label{cognitive-science-foundations}

\citet{gardenfors2000} proposed Conceptual Spaces as geometric similarity spaces organized along quality dimensions, with natural concepts forming convex regions and betweenness relations defining navigable paths. This theoretical framework posits that conceptual structure is inherently geometric --- concepts have locations, distances between them are meaningful, and navigation between concepts follows principled geometric trajectories.

Cognitive neuroscience has found biological mechanisms consistent with this theory. \citet{tolman1948cognitive} proposed cognitive maps for spatial navigation. \citet{okeefe1971hippocampus} discovered place cells in the hippocampus. \citet{constantinescu2016grid} demonstrated that grid cells in the human entorhinal cortex fire during \emph{conceptual} navigation --- not just physical movement through space but traversal through abstract conceptual dimensions. The spatial navigation machinery is repurposed for abstract thought.

Two additional cognitive science results bear directly on our findings. \citet{tversky1977features} established that human similarity judgments are asymmetric: North Korea is judged more similar to China than China is to North Korea. This is the direct cognitive parallel to our quasimetric finding --- asymmetry in conceptual distance is not an artifact of autoregressive generation but a property also observed in human cognition. \citet{collins1975spreading} proposed spreading activation through semantic networks, where activation propagates directionally from a source concept through associative links --- a process model that maps onto the early-anchoring and bridge-bottleneck phenomena we observe.

More recently, LLM World of Words~\citep{llmwow2025} demonstrated that free association norms elicited from LLMs are directly comparable to human norms collected over decades, establishing that LLM conceptual organization shares at least surface-level structure with human semantic memory.

\textbf{Our contribution.} Gardenfors theorized that concepts live in navigable geometric spaces. Cognitive neuroscience found biological mechanisms (grid cells, place cells) that support navigation through those spaces. We provide a computational probe of this theoretical framework at scale --- 12 models, 100+ concept pairs, \textasciitilde21,540 runs. The partial confirmation is substantive: LLM behavior is consistent with concepts inhabiting navigable spaces with measurable, compositional distance structure (triangle inequality \textasciitilde91\%, hierarchical transitivity 4.9x over random). But the pervasive asymmetry (mean 0.811) challenges the standard metric assumption in conceptual spaces theory, suggesting that quasimetric structure may be a more appropriate formal framework. We do not claim that LLMs replicate human conceptual navigation --- only that they exhibit analogous geometric regularities testable by similar methods. The parallel to Tversky's asymmetry and Constantinescu's grid-cell navigation is structural, not mechanistic.

\begin{center}\rule{0.5\linewidth}{0.5pt}\end{center}

The seven literatures reviewed here converge on a single gap. Static geometry is well-characterized and its theoretical origin is understood. The mathematical framework for navigation exists. Consistency and directionality are documented as phenomena. Multi-hop composition is studied as a reasoning failure mode. Cross-model convergence is measured at the representational level. Cognitive science provides both the theoretical framework and the biological precedent. What is missing --- and what this benchmark provides --- is an empirical test of whether LLM conceptual navigation has consistent, measurable geometric structure: not what the space looks like, but whether models can walk through it reliably, and what the shape of their walking reveals.
     % Section 3: Related Work
\section{Act I --- Structure: Models Have Navigational Geometry}\label{act-i-structure-models-have-navigational-geometry}

The first three experimental phases established the empirical foundation for the benchmark: models navigate conceptual space with measurable, stable, model-specific structure. This section summarizes the five structural findings from the early phases; some were later replicated across expanded model cohorts (R1), while others remain well-supported observations from the original cohort (O1, O2, O9, O10).

\subsection{Conceptual Gaits Are Distinct and Stable}\label{conceptual-gaits-are-distinct-and-stable}

The most robust finding in the benchmark is that each model navigates conceptual space with a characteristic level of consistency --- a property we term its \emph{conceptual gait}. Gait is operationalized as the mean pairwise Jaccard similarity across independent runs on the same (model, pair, direction) condition: high gait indicates that a model produces similar waypoints each time it traverses a given conceptual path; low gait indicates that each traversal explores different intermediate territory.

\textbf{The 12-model gait spectrum.} Across 12 models from 11 independent training pipelines, conceptual gait spans a 2.9× range:

\begin{longtable}[]{@{}llll@{}}
\toprule\noalign{}
Model & Mean Jaccard & Cohort & Phases \\
\midrule\noalign{}
\endhead
\bottomrule\noalign{}
\endlastfoot
Mistral Large 3 & 0.747 & Phase 11A & 11 \\
Claude Sonnet 4.6 & 0.578 & Original & 1--11 \\
DeepSeek V3.2 & 0.540 & Phase 11A & 11 \\
Llama 4 Maverick & 0.539 & Phase 11A & 11 \\
Qwen 3.5 397B-A17B & 0.508 & Phase 10A & 10 \\
Cohere Command A & 0.502 & Phase 11A & 11 \\
MiniMax M2.5 & 0.419 & Phase 10A & 10 \\
Moonshot Kimi K2.5 & 0.414 & Phase 10A & 10 \\
Gemini 3 Flash & 0.372 & Original & 1--9 \\
Llama 3.1 8B Instruct & 0.298 & Phase 10A & 10 \\
Grok 4.1 Fast & 0.293 & Original & 1--9 \\
GPT-5.2 & 0.258 & Original & 1--11 \\
\end{longtable}

At the extremes, Mistral produced near-deterministic paths on several pairs (Jaccard 0.936 on music→mathematics, 0.934 on hot→cold), while GPT's most consistent experimental pair (hot→cold, Jaccard 0.911 in Phase 1) was an outlier against its otherwise exploratory profile.

\textbf{Temporal stability.} For the four core models (Claude, GPT, Grok, Gemini), gait consistency was measured across Phases 1 through 9 --- a span of multiple days and over 9,500 runs. The gait ordering (Claude \textgreater{} Gemini \textgreater{} Grok \textgreater{} GPT) was preserved throughout. Claude's mean Jaccard remained within {[}0.55, 0.60{]} across all phases in which it was tested; GPT's remained within {[}0.24, 0.28{]}. No systematic drift was detected.

\textbf{Cross-architecture generality.} Gait was initially characterized on 4 models from 4 providers (Phase 1). Phase 10A added 4 models from 4 additional providers, all falling within the previously observed range (0.298--0.508). Phase 11A added 4 more models extending the range upward (Mistral 0.747). Across 12 models from 11 independent training pipelines, every model exhibited a characteristic gait, confirming that this is a universal property of LLM conceptual navigation rather than a quirk of the original cohort.

\textbf{Protocol robustness.} Phase 11C tested whether gait was an artifact of the default elicitation protocol (7 waypoints, temperature 0.7) by varying both parameters in a 2×2 grid (5 and 9 waypoints × temperatures 0.5 and 0.9) for three models (Claude, GPT, DeepSeek). An ANOVA decomposition found that model identity accounted for the largest share of variance in Jaccard consistency (η² = 0.242, p ≈ 0.001), while waypoint count (η² = 0.008) and temperature (η² = 0.002) were negligible. The gait rank ordering across models was largely preserved (Kendall's W = 0.840), though not perfectly --- GPT and DeepSeek swapped positions in 4 of 5 conditions, with DeepSeek's baseline condition showing anomalously low Jaccard that was not replicated at other settings. These p-values are approximate (chi-squared approximation), and the robustness test covers only three models (Claude, GPT, DeepSeek); it addresses waypoint-count and temperature variation but not prompt sensitivity more broadly.

\textbf{Control validation.} The nonsense control pair (xkplm→qrvzt) produced mean Jaccard of 0.062 across all models --- effectively zero consistency. GPT's nonsense Jaccard ranged from 0.006 to 0.015, indistinguishable from chance. This supports the interpretation that gait is content-sensitive rather than a decoding artifact: when there is no conceptual structure to navigate, consistency vanishes. (Control validation is qualified by Phase 11B's finding that single-pair controls are insufficient for broader generality testing --- see Section 2.6 --- but the relative gap between experimental and control conditions remains large and consistent across the original cohort.)

\textbf{Interpretation.} Gait reflects a model's propensity to re-traverse the same conceptual territory across independent runs. A high-gait model like Mistral or Claude has a strongly determined internal map --- given two concepts, it finds essentially the same intermediate landmarks every time. A low-gait model like GPT explores broadly, sampling different landmarks on each traversal while still arriving at the destination. Neither extreme is intrinsically better; they represent different navigational strategies. Crucially, gait is pair-dependent even within a model: Claude's Jaccard on hot→cold (0.911) far exceeds its Jaccard on random control pairs (0.291--0.343), indicating that gait interacts with the structure of the conceptual terrain, not just the model's decoding behavior.

\subsection{Paths Are Self-Consistent Within Models}\label{paths-are-self-consistent-within-models}

Before attributing structure to waypoint paths, it is necessary to establish that they are stable model properties rather than transient artifacts of API versioning or server state.

\textbf{No temporal drift.} Phase 4B (the targeted bridge topology experiment) collected data across multiple days. Cross-batch Jaccard --- the similarity between runs collected on different days --- fell within 0.05 of within-batch Jaccard for all tested (model, pair) combinations. The navigational behavior measured in the benchmark was stable over the tested collection window.

\textbf{Scale consistency.} Phase 1 collected paths at both 5 and 10 waypoints for all 21 reporting pairs and 4 models. The 5-waypoint paths were not independent inventions but genuine coarse-grainings of the 10-waypoint paths: on average, 70.5\% of 5-waypoint concepts also appeared in the corresponding 10-waypoint paths, and 67.9\% of the time, the 5-waypoint sequence appeared as a proper subsequence of the 10-waypoint sequence --- same concepts in the same order, with additional waypoints interpolated between them. For example, GPT's 5-waypoint path from hot to cold (\emph{warm, tepid, lukewarm, cool, chilly}) appeared intact within its 10-waypoint path (\emph{scald, sear, warm, tepid, lukewarm, cool, chilly, brisk, frigid, icy}), with the longer path extending coverage at both endpoints and interpolating finer detail.

These two findings support the conclusion that waypoint elicitation measures a stable, resolution-consistent property of each model's navigational geometry, not a transient or scale-dependent artifact.

\subsection{The Dual-Anchor Effect}\label{the-dual-anchor-effect}

Phase 2 established that conceptual navigation is fundamentally asymmetric: forward (A→B) and reverse (B→A) paths share less than 19\% of their waypoints on average (full results in Section 5). This prompted the starting-point hypothesis --- that models construct paths by greedy forward search from the starting concept, producing near-complete divergence between forward and reverse trajectories. Phase 3A tested this hypothesis by examining where in the waypoint sequence forward and reverse paths converge.

\textbf{The U-shaped convergence profile.} To measure positional convergence, forward position \emph{i} was compared against reverse position \emph{(N+1−i)} --- the mirror position. If forward waypoint 1 matched reverse waypoint 5, forward waypoint 2 matched reverse waypoint 4, and so on, this would indicate that both paths traversed the same intermediate concepts in opposite order. The per-position mirror-match rates across all 84 pair/model combinations were:

\begin{longtable}[]{@{}llll@{}}
\toprule\noalign{}
Position & Forward & Reverse (mirror) & Match Rate \\
\midrule\noalign{}
\endhead
\bottomrule\noalign{}
\endlastfoot
1 & wp1 & wp5 & 0.102 \\
2 & wp2 & wp4 & 0.057 \\
3 & wp3 & wp3 & 0.065 \\
4 & wp4 & wp2 & 0.085 \\
5 & wp5 & wp1 & 0.129 \\
\end{longtable}

The starting-point hypothesis predicted a monotonic increase from position 1 to position 5 --- early waypoints should diverge most (each direction dominated by its starting concept's neighborhood), with convergence increasing as both paths approach the shared target. Instead, the data revealed a U-shape: positions 1 and 5 were both elevated relative to the middle. The overall convergence slope was weakly positive (0.0082, 95\% CI {[}−0.0030, 0.0198{]}) and only 50\% of pair/model combinations showed a positive slope. The starting-point hypothesis, in its simple form, was not supported.

\textbf{The dual-anchor interpretation.} Both endpoints exert ``gravitational pull'' on nearby waypoints. Position 1 (forward waypoint 1 vs.~reverse waypoint 5) shows elevated overlap because both measurements sample from the neighborhood of concept A --- the forward path's first step away from A and the reverse path's last step arriving at A draw from the same constrained vicinity. Position 5 shows even higher overlap for the same reason applied to concept B. The middle positions (2--4) form a valley because this is the region of maximum navigational freedom --- unconstrained by either endpoint, dominated by model-specific and run-specific choices.

\textbf{Category-dependent signatures.} The U-shape was not uniform across relationship types. Each category produced a distinctive convergence profile:

\begin{itemize}
\item
  \emph{Antonyms} (hot--cold) showed extreme late convergence: 0.054, 0.035, 0.215, 0.338, 0.344. The temperature axis acts as a funnel --- once either direction has traversed 2--3 waypoints, both paths enter the same tightly constrained central region. This was the one category that supported the original starting-point hypothesis (slope 0.088, CI above zero).
\item
  \emph{Identity controls} (apple--apple) showed middle domination: 0.000, 0.044, 0.319, 0.048, 0.018. With no distance to traverse, both paths circulate through the same core associations (fruit, tree, harvest), and the middle waypoint is where they are most likely to align.
\item
  \emph{Random controls} showed a pure U-shape: 0.099, 0.025, 0.010, 0.031, 0.121. Both endpoints exert local pull, but no semantic bridge constrains the middle. The path between telescope and jealousy wandered freely through intermediate territory.
\item
  \emph{Nonsense controls} were flat near zero: 0.001, 0.003, 0.008, 0.001, 0.009. No anchoring effect, because nonsense strings have no semantic neighborhoods to constrain waypoint selection.
\end{itemize}

The category signatures in the convergence profile encode the type of semantic relationship between endpoints. Tightly constrained relationships (antonyms, hierarchies) show structured convergence; unconstrained relationships (random, nonsense) show the pure dual-anchor pattern or no pattern at all. This suggests that positional convergence profiles may help discriminate relationship types, though this observation has not been validated as a predictive tool in subsequent phases.

\textbf{Model-level variation.} Claude showed the only negative convergence slope (−0.005) among the four core models, meaning its U-shape tilted slightly toward the start-side anchor. This is consistent with Claude's rigid navigational style: its paths are so determined by the pair's global topology that the target exerts less incremental pull than it does for more variable navigators. GPT (0.015) and Grok (0.016) showed positive slopes, consistent with their more exploratory styles producing stronger late-stage convergence toward the target.

\subsection{Three Models Agree, Gemini Diverges}\label{three-models-agree-gemini-diverges}

Across the 6 non-control triples analyzed in Phase 4A, Claude, GPT, and Grok aligned more closely with one another than with Gemini, which systematically diverged. The sample is small (6 triples, with two universal-zero triples that inflate pairwise correlations), but the pattern was consistent across multiple metrics.

\textbf{Bridge frequency agreement.} Phase 4A analyzed inter-model agreement on bridge concept usage across 6 non-control triples. Claude, GPT, and Grok showed 100\% binary bridge agreement --- for every triple, all three models either found the bridge or all three missed it. Gemini agreed with the other models on only 50\% of triples.

Pairwise Pearson correlations on the 6-element bridge frequency vectors:

\begin{longtable}[]{@{}lll@{}}
\toprule\noalign{}
Model Pair & Pearson \emph{r} & Binary Agreement \\
\midrule\noalign{}
\endhead
\bottomrule\noalign{}
\endlastfoot
Claude--GPT & 0.772 & 100\% \\
Claude--Grok & 0.768 & 100\% \\
GPT--Grok & 0.667 & 100\% \\
GPT--Gemini & 0.679 & 50\% \\
Claude--Gemini & 0.446 & 50\% \\
Grok--Gemini & 0.340 & 50\% \\
\end{longtable}

To quantify Gemini's divergence, the mean absolute bridge frequency difference (\textbar Δfreq\textbar) was computed for Gemini-paired model pairs versus non-Gemini pairs. Gemini-paired mean \textbar Δfreq\textbar{} was 0.336, compared to 0.200 for non-Gemini pairs --- an isolation index of 0.136. The 95\% CI {[}−0.092, 0.378{]} includes zero, so this trend may not be reliable at the aggregate level. The qualitative pattern, however, was consistent: Gemini's bridge frequency disagreements were concentrated on specific triples where the other three models agreed.

\textbf{Characteristic divergence patterns.} The divergence was not random. Gemini consistently failed to route through bridge concepts that the other three models found reliably:

\begin{itemize}
\tightlist
\item
  \emph{bank→ocean}: Claude, GPT, and Grok all routed through ``river'' (frequency 0.90--1.00); Gemini did not (0.00), following alternative non-river routes.
\item
  \emph{music→mathematics}: Claude (1.00), Grok (0.85), and GPT (0.15) routed through ``harmony''; Gemini did not (0.00).
\item
  \emph{emotion→melancholy via nostalgia}: Grok found the bridge reliably (0.90); Claude and GPT found it at low frequency (0.10); Gemini did not (0.00).
\end{itemize}

The only case where all four models agreed on strong bridge presence was the hierarchical triple animal→dog→poodle (bridge frequency 0.95--1.00). (All four models also agreed on universal zero for the hot→energy→cold and beyoncé→justice→erosion triples, where no model found the designated bridge.) Gemini's topology preserved hierarchical/taxonomic structure but fragmented on cross-domain and abstract bridges.

\textbf{Connection to later findings.} The Gemini divergence persisted through the later core-cohort phases in which Gemini was tested (through Phase 9). Three mechanistic hypotheses for Gemini's deficit --- frame-crossing threshold (Phase 5), gradient blindness (Phase 8B), and transformation-chain blindness (Phase 9B) --- were all falsified (graveyard entries G7, G21, G24). The deficit is real, stable within the tested phases, and mechanistically uncharacterized as of the end of data collection. It is discussed further in Section 7.5.

\subsection{Polysemy Sense Differentiation}\label{polysemy-sense-differentiation}

Polysemous concepts --- words with multiple distinct meanings --- provided a natural test of whether waypoint elicitation measures genuine conceptual structure or surface-level word statistics. If paths reflect word co-occurrence patterns, a polysemous source word should produce partially overlapping paths regardless of the target (since the same word token drives the associations). If paths reflect conceptual navigation, different targets should activate different senses and produce non-overlapping paths.

\textbf{The result was unambiguous.} Three polysemy groups were tested in Phase 1 (bank, bat, crane), each with two sense-steering target concepts. Phase 2 supplementary runs corrected an initial artifact (the holdout pairs had lacked pilot data, inflating the Phase 1 zero-overlap result). The corrected cross-pair Jaccard values --- measuring the overlap between paths from the same source word to different sense-targets --- were:

\begin{longtable}[]{@{}llll@{}}
\toprule\noalign{}
Group & Pair A & Pair B & Cross-Pair Jaccard \\
\midrule\noalign{}
\endhead
\bottomrule\noalign{}
\endlastfoot
bank & bank→river & bank→mortgage & 0.062 \\
bat & bat→cave & bat→baseball & 0.059 \\
crane & crane→construction & crane→wetland & 0.011 \\
\end{longtable}

All values were near zero. The crane group (0.011) was particularly clean --- essentially a single shared waypoint across all models and runs between the construction-crane and bird-crane paths. For comparison, within-pair Jaccard for these same models on meaningful pairs ranged from 0.258 (GPT) to 0.578 (Claude), and even the random control mean was 0.337.

\textbf{Interpretation.} Under the benchmark's lexical similarity measure, the target concept strongly determines which sense of the polysemous source is activated, and this sense activation redirects the navigational trajectory. The word ``bank'' paired with ``river'' activates a geographic/hydrological frame; the same word paired with ``mortgage'' activates a financial frame. The resulting paths share almost no intermediate territory. This is strong evidence that waypoint elicitation probes conceptual structure --- the semantic content of the path --- rather than surface statistics about the source word.

The polysemy result also connects to the dual-anchor finding (Section 4.3): even though the source concept (``bank'') is identical in both conditions, the target concept exerts sufficient gravitational pull to redirect the path from waypoint 1 onward. This is incompatible with a pure forward-chaining model in which only the starting concept determines the trajectory. The target's influence is detectable from the very beginning of the path, at least when the source is polysemous and the target provides strong sense-disambiguating signal.

\textbf{Extension to later phases.} Phase 5B extended the polysemy finding to forced crossings --- concept triples where a polysemous word serves as the only connection between two domains (e.g., loan→bank→shore). These triples produced the highest bridge frequencies in the benchmark (0.95--1.00 for three of four models), confirming that polysemous concepts can serve as obligatory navigational bottlenecks when they are the sole path between semantic domains. The full forced-crossing analysis appears in Section 5.4.

\begin{center}\rule{0.5\linewidth}{0.5pt}\end{center}

The five findings in this section --- distinct gaits, self-consistent paths, dual-anchor convergence, cross-model agreement structure, and polysemy sense differentiation --- collectively support the conclusion that waypoint elicitation measures real geometric structure in LLM conceptual navigation, as captured by lexical overlap under the current canonicalization pipeline (Section 2.3). The structure is stable over time, consistent across resolutions, driven by semantic content, and characteristic of each model. This foundation supports the more demanding geometric claims in the sections that follow.
        % Section 4: Act I -- Structure
\section{Act II --- Topology: Navigational Geometry Is Consistent with Quasimetric Structure}\label{act-ii-topology-navigational-geometry-is-consistent-with-quasimetric-structure}

Section 4 established that waypoint elicitation measures real navigational structure. This section characterizes the \emph{geometry} of that structure. The central finding is that LLM conceptual spaces exhibit properties consistent with a quasimetric space --- satisfying non-negativity and triangle inequality while systematically violating symmetry. We present the evidence for each property, then show that the space exhibits compositional structure (hierarchical paths compose), selective bottleneck routing (bridge concepts are not mere associations), and non-uniform navigational traffic.

\textbf{Methodological caveat.} The asymmetry result (Section 5.1) and the triangle inequality result (Section 5.2) use different operationalizations of ``distance.'' Asymmetry is measured via cross-direction Jaccard (comparing A→B paths against B→A paths), while the triangle inequality uses within-direction Jaccard (consistency of repeated runs on a single direction) to define navigational distance. The quasimetric characterization is therefore \emph{consistent with} the data rather than formally established on a single unified distance function. We present this as convergent evidence from two complementary distance measures, not as a deductive proof.

\subsection{Asymmetry Is Fundamental}\label{asymmetry-is-fundamental}

Phase 2 reversed all 21 reporting pairs across 4 models (840 reverse-direction runs; \textasciitilde960 total including polysemy supplementary) and computed the directional asymmetry index: 1 minus the mean Jaccard similarity between forward (A→B) and reverse (B→A) waypoint sets. A value of 0 indicates perfect symmetry (identical paths in both directions); a value of 1 indicates complete non-overlap.

\textbf{The symmetry axiom fails comprehensively.} The mean directional asymmetry across all 84 pair/model combinations was 0.811 (95\% CI {[}0.772, 0.848{]}). Of 84 combinations, 73 (87\%) showed statistically significant asymmetry by permutation test (p \textless{} 0.05, 1,000 resamples). The distribution was heavily right-skewed: 54 of 84 combinations (64\%) exceeded 0.8, and only 5 (6\%) fell below 0.5. Forward and reverse paths typically shared less than 19\% of their waypoints.

\textbf{Contextualizing the 0.811 mean.} The Phase 2 asymmetry was computed on the original 4-model cohort at 5 waypoints, and the 0.811 mean includes control pairs --- nonsense controls (0.986) and random controls (0.908) --- which inflate the overall average. Experimental-category means were lower but still substantial (antonym 0.596, hierarchy 0.683, cross-domain 0.822, polysemy 0.824, anchor 0.911). The 0.811 figure should be understood as characterizing the full distribution rather than as a universal constant.

Phase 11A confirmed asymmetry across 4 additional models, all exceeding the 0.60 threshold: Mistral 0.729, DeepSeek 0.722, Cohere 0.718, Llama 4 Maverick 0.673. Phase 11C revealed a resolution-dependence qualification in the 3-model robustness subset (Claude, GPT, DeepSeek): only the 9-waypoint conditions cleared the 0.60 threshold (0.669--0.684); the 5-waypoint conditions (0.593--0.594) and the 7-waypoint baseline (0.599) did not. The asymmetry property is real but requires sufficient path length to manifest reliably --- a measurement sensitivity finding, not a protocol artifact.

\textbf{Statistical caveat.} The 87\% significance rate was computed at α = 0.05 across 84 independent tests with no explicit multiple-comparison correction. With Bonferroni correction (threshold ≈ 0.0006), the percentage would be lower, and the permutation test's resolution (1,000 resamples, minimum achievable p = 0.001) limits precision at corrected thresholds. The qualitative conclusion --- that the vast majority of pair/model combinations show genuine asymmetry --- is robust, but the specific 87\% figure should be interpreted conservatively.

\textbf{Theoretical connection.} \citet{tversky1977features} established that human similarity judgments are asymmetric: North Korea is judged more similar to China than China is to North Korea. \citet{berglund2024reversal} showed that autoregressive models exhibit a reversal curse --- training on ``A is B'' does not teach ``B is A.'' Our finding connects these: the reversal curse is not merely a failure mode but produces structured asymmetry consistent with quasimetric topology. The asymmetry patterns track semantic properties of the concept pairs (Section 4.3) and model-specific navigational styles (Section 4.1), indicating systematic structure rather than noise.

\subsection{Triangle Inequality Shows Stable Replication}\label{triangle-inequality-shows-stable-replication}

The triangle inequality --- d(A,C) ≤ d(A,B) + d(B,C) --- was tested across three independent samples using navigational distance defined as 1 minus mean within-direction Jaccard similarity.

\begin{longtable}[]{@{}
  >{\raggedright\arraybackslash}p{(\columnwidth - 8\tabcolsep) * \real{0.1373}}
  >{\raggedright\arraybackslash}p{(\columnwidth - 8\tabcolsep) * \real{0.2549}}
  >{\raggedright\arraybackslash}p{(\columnwidth - 8\tabcolsep) * \real{0.1569}}
  >{\raggedright\arraybackslash}p{(\columnwidth - 8\tabcolsep) * \real{0.2157}}
  >{\raggedright\arraybackslash}p{(\columnwidth - 8\tabcolsep) * \real{0.2353}}@{}}
\toprule\noalign{}
\begin{minipage}[b]{\linewidth}\raggedright
Phase
\end{minipage} & \begin{minipage}[b]{\linewidth}\raggedright
N Triangles
\end{minipage} & \begin{minipage}[b]{\linewidth}\raggedright
Models
\end{minipage} & \begin{minipage}[b]{\linewidth}\raggedright
\% Holding
\end{minipage} & \begin{minipage}[b]{\linewidth}\raggedright
Violations
\end{minipage} \\
\midrule\noalign{}
\endhead
\bottomrule\noalign{}
\endlastfoot
3B & 8 & 4 & 90.6\% (29/32) & Gemini ×2 (music-harmony-math, bank-river-ocean), Claude ×1 marginal (hot-energy-cold) \\
4B & 8 & 4 & 93.8\% (30/32) & Claude ×1 (tree-forest-ecosystem), Grok ×1 (light-chandelier-color, a random-control triple) \\
7B & 8 & 4 & 90.6\% (29/32) & Claude ×1 (deposit-bank-shore), Gemini ×2 (seed-germination-garden, music-harmony-math) \\
\end{longtable}

The satisfaction rate shows a stable replication pattern around 91\% across three independent triple sets. This is remarkably high given the magnitude of the asymmetry violations: a space where forward and reverse distances differ by 0.811 on average might be expected to show widespread triangle inequality failures, yet the directional distances compose properly in over 90\% of cases.

The violations in Phase 3B were concentrated in Gemini: the music→harmony→mathematics triple showed a slack of −0.319 (the direct path was much longer than the sum of the two legs), and bank→river→ocean showed −0.061. Claude's hot→energy→cold violation was marginal (−0.023). GPT and Grok each satisfied the triangle inequality in all 8 of their respective Phase 3B cases.

\textbf{Quasimetric characterization.} Combined with the asymmetry data, the metric-axiom profile of LLM conceptual space is: - \emph{Identity:} not directly testable with Jaccard-based distance (identity controls show low directional divergence for some models --- Claude 0.017 --- but mean cross-direction asymmetry was 0.456, and within-direction distance was not measured for identity pairs) - \emph{Non-negativity:} holds trivially - \emph{Symmetry:} \textbf{fails comprehensively} (mean violation 0.811) - \emph{Triangle inequality:} holds in \textasciitilde91\% of cases

This profile is consistent with a quasimetric space --- a space satisfying all metric axioms except symmetry. Quasimetric spaces are standard mathematical objects arising in directed graphs, asymmetric norms, and computational complexity theory. The characterization provides a useful framework for interpreting the navigational geometry observed throughout the benchmark, though it rests on convergent evidence from two complementary distance operationalizations rather than a single unified metric (see methodological caveat above).

\subsection{Hierarchical Paths Are Compositional}\label{hierarchical-paths-are-compositional}

If conceptual navigation merely produces plausible-sounding word chains, there is no reason to expect that the path from A to C should systematically include the concept B that lies ``between'' them in a taxonomic hierarchy. Phase 3B tested this directly with 8 concept triples across 4 models.

\textbf{Hierarchical triples showed 4.9× higher waypoint transitivity than random controls.} Mean transitivity for hierarchical triples was 0.175 (95\% CI {[}0.112, 0.238{]}); for random controls, 0.036 (95\% CI {[}0.011, 0.070{]}). The confidence intervals do not overlap. Bridge concepts appeared systematically for taxonomic triples --- ``dog'' at 15--100\% frequency on the animal→poodle path --- and never for random controls (``stapler'' at 0\% on music→mathematics, ``flamingo'' at 0\% on hot→cold).

\textbf{Polysemy-extend triples} (bank→river→ocean) showed even higher bridge frequency (0.725) than hierarchical triples (0.456), reflecting the bottleneck effect of sense-mediated routing: once ``bank'' resolves to its geographic sense, ``river'' is the only natural waypoint en route to ``ocean.'' Semantic-chain triples (music→harmony→mathematics, hot→energy→cold) showed intermediate transitivity (0.073) with strong model dependence --- ``harmony'' appeared on 100\% of Claude's music→mathematics paths but 0\% of Gemini's.

The key negative result was hot→energy→cold, where ``energy'' never appeared as a bridge despite being associated with both endpoints. The hot→cold paths traversed the experiential temperature gradient (warm, tepid, cool, chilly); the hot→energy and energy→cold paths traversed physics terminology (thermodynamics, work, heat). Association with both endpoints is necessary but not sufficient for bridge function --- the concept must lie \emph{on the navigational path} between the endpoints, not merely in their shared associative neighborhood.

\subsection{Bridge Concepts Are Bottlenecks, Not Associations}\label{bridge-concepts-are-bottlenecks-not-associations}

Phases 4 and 5 converged on a characterization of what makes a concept function as a navigational bridge. Associative strength alone is insufficient; structural role --- whether the concept names the primary axis of connection between the endpoints --- can override it (though cue strength remains a real contributing factor; see Section 5.5).

\textbf{Bottleneck bridges succeed universally.} ``Spectrum'' (the mechanism by which light produces color) appeared at 1.00 frequency across all four models on the light→color path. ``Deposit'' (the action connecting bank to savings) appeared at 0.95--1.00. ``Sentence'' (the compositional unit between word and paragraph) appeared at 1.00. These bridges name the \emph{mechanism} or \emph{transformation} that connects the endpoints.

\textbf{Off-axis associations fail universally.} ``Metaphor'' (associated with both language and thought) appeared at 0.00 on the language→thought path. ``Energy'' appeared at 0.00 on hot→cold. These concepts are associated with both endpoints but do not name the primary axis of connection.

\textbf{Process-naming outperforms object-naming.} ``Germination'' (the process connecting seed to garden) outperformed ``plant'' (the object that inhabits both) in all four models: Claude 1.00 vs 0.00, GPT 0.95 vs 0.65, Grok 0.15 vs 0.00, Gemini 1.00 vs 0.00 (O4). The explanation parallels the spectrum/deposit pattern: ``germination'' names the \emph{transformation} from seed to garden, adding directional information; ``plant'' names something already implied by both endpoints, adding no navigational value. Navigational salience depends on directional information content, not solely on generic associative strength.

\textbf{Too-central concepts are skipped.} ``Fire'' appeared at near-zero frequency across all configurations in Phase 5B --- 0.00 on spark→ash (except GPT 0.15), 0.00 on ambition→motivation, 0.00 on spark→motivation (O6). Both ``spark'' and ``ash'' already imply fire; naming it explicitly adds no information. Phase 7C extended this finding: ``water'' for rain→ocean was 0.00 across all four models (O16). Too-central concepts are navigable \emph{to} but not \emph{through} --- a distinct failure mode from off-axis associations.

\textbf{Forced crossings create near-obligatory bottlenecks.} When a polysemous concept is the sole connection between two domains, it becomes a near-obligatory waypoint under unconstrained elicitation. Loan→bank→shore produced bridge frequency of 0.95--1.00 for Claude, GPT, and Grok (O5); Gemini showed 0.05, consistent with its apparent topological disconnection between senses. This is a bank-specific result --- the benchmark's strongest forced-crossing case --- and generality beyond ``bank'' remains untested. The forced-crossing mechanism is the inverse of the too-central phenomenon: instead of being redundantly implied, the bridge is the \emph{only} route between otherwise disconnected domains.

\textbf{A major failed prediction} was ``metaphor'' as the language→thought bridge (graveyard entry G5). Intuitive semantic importance had zero predictive power for navigational bridge behavior. This failure sharpened the operational definition: bridges must name the primary axis of connection between endpoints, not merely be associated with both.

\subsection{Cue-Strength Gradient}\label{cue-strength-gradient}

Phase 5A tested whether bridge frequency decreases monotonically with the bridge concept's cue strength --- its associative proximity to both endpoints. Four concept families were constructed, each with three bridge candidates at varying cue levels (very-high, medium, low), and bridge frequency was measured across 4 models (2,040 new runs).

\textbf{The gradient is real.} Twelve of 16 family/model combinations showed monotonic decrease in bridge frequency from highest to lowest cue level. All four failures occurred in a single anomalous family (biological-growth: seed→B→garden), where the ``medium'' cue bridge (``germination'') outperformed the ``very-high'' bridge (``plant'') --- the process-naming effect described in Section 5.4. The three well-behaved families --- physical-causation (sun→B→desert), compositional-hierarchy (word→B→paragraph), and abstract-affect (emotion→B→melancholy) --- showed perfect monotonic gradients across all four models.

Logistic regression revealed similar cue-strength thresholds across models (Claude 2.02, GPT 1.96, Grok 2.16, Gemini 1.79), with the difference CI {[}−1.00, 1.07{]} including zero. This falsified the Phase 4 hypothesis that Gemini's bridge failures stemmed from a higher cue-strength threshold. Gemini responds to cue strength at approximately the same threshold as other models; its failures must be driven by something other than associative sensitivity. R² values ranged from 0.311 to 0.640, indicating that cue strength is a real but incomplete predictor --- bridge selection is closer to a threshold phenomenon than a smooth gradient.

\subsection{Waypoint Distributions Are Non-Uniform}\label{waypoint-distributions-are-non-uniform}

Phase 6A collected 1,200 runs across 8 concept pairs and 4 models to characterize the full waypoint frequency distribution for each (pair, model) condition.

\textbf{Navigational traffic concentrates in a small number of high-frequency waypoints.} Of 32 pair-model conditions (8 pairs × 4 models), 31 rejected the null hypothesis of uniform waypoint frequency at Bonferroni-corrected significance (chi-squared goodness-of-fit test, threshold p \textless{} 0.0063). The sole non-rejection was Claude on seed→garden, where only 5 unique waypoints were observed --- near-deterministic navigation leaving insufficient diversity to distinguish from uniform.

Entropy varied systematically with gait: Claude showed the lowest mean entropy (2.59), reflecting near-deterministic waypoint selection from a small vocabulary; GPT showed the highest (3.44), reflecting broad exploration across a larger candidate set. This mirrors the gait spectrum from Section 4.1 --- high-gait models concentrate traffic on a few landmarks, while low-gait models distribute it across many.

The non-uniformity finding is important because it rules out a class of null models. If waypoints were drawn approximately uniformly from a large vocabulary, the consistency metrics throughout the benchmark could reflect sampling noise rather than navigational structure. The non-uniform, concentrated distributions confirm that models navigate through preferred waypoints, not through arbitrary word generation.

\subsection{Bridge Positioning: Early Anchoring, Not Midpoint}\label{bridge-positioning-early-anchoring-not-midpoint}

Phase 5C tested whether bridge concepts create a convergence peak at the midpoint of waypoint paths (the W-shape hypothesis). The aggregate result was null --- bridge-present pairs showed no significant midpoint convergence signal when averaged across models and pairs. Phase 6C revised the approach: instead of assuming the bridge occupies the midpoint, it scanned all positions for bridge frequency peaks.

\textbf{Bridges anchor early, not at the midpoint.} Across 10 pairs and 4 models (40 positional profiles), the cross-model modal bridge position was 1--2 (0-indexed; the 2nd--3rd waypoint out of 7) for 8 of 10 pairs. The peak-detection contrast --- bridge frequency at the modal position minus the mean of non-modal positions --- was 0.345 (95\% CI {[}0.224, 0.459{]}), significantly positive. The fixed-midpoint contrast was −0.080 (95\% CI {[}−0.141, −0.024{]}), confirming that the Phase 5C W-shape null was a methodological artifact of the rigid midpoint assumption, not evidence against bridge positional structure.

\textbf{The taxonomic exception.} The sole departure from early anchoring was the hierarchical pair animal→poodle, where ``dog'' anchored at positions 4--5 across models. This late positioning reflects the structure of the taxonomic chain: the path from animal to poodle narrows progressively (mammal → carnivore → canine → dog → breed), and ``dog'' is the penultimate step, not an early waypoint. Taxonomic bridges reflect their position in the hierarchy rather than a universal early-anchoring tendency.

\textbf{Forced-crossing bridges are positionally unstable.} The two forced-crossing pairs (loan→shore, deposit→river) showed positional standard deviation of 1.71, compared to 0.52 for the 8 non-forced pairs (O13). The polysemous pivot sits at a frame junction where the timing of sense-crossing is model-dependent: Claude places ``bank'' at position 1 on loan→shore; Gemini places it at position 4. The bridge is obligatory (it must appear somewhere) but its position is not fixed --- a contrast with non-forced bridges like ``spectrum'' or ``sadness,'' which reliably anchor at positions 1--2.

\begin{center}\rule{0.5\linewidth}{0.5pt}\end{center}

The seven findings in this section collectively support a characterization of LLM conceptual space as consistent with quasimetric geometry, exhibiting compositional structure and non-uniform navigational traffic. Distances are asymmetric but compose properly (triangle inequality \textasciitilde91\%); hierarchical relationships preserve transitivity; bridge concepts function as structural bottlenecks, not associative co-occurrences; and navigational traffic concentrates at preferred waypoints that anchor early in the path sequence. This geometric characterization provides the foundation for the causal intervention experiments in the next section, which probe whether the observed structure is merely statistical regularity or reflects genuine causal dependencies in the navigation process.
         % Section 5: Act II -- Topology
\input{06-mechanism}        % Section 6: Act III -- Mechanism
\input{07-limits}           % Section 7: Act IV -- Limits
\section{Act V --- Generality: Structure Scales Across Architectures}\label{act-v-generality-structure-scales-across-architectures}

The structural findings in Sections 4--6 were established on four models from four providers. A natural concern is that the properties might reflect idiosyncrasies of those specific models rather than general features of language model navigation. Phases 10A and 11A addressed this by testing 8 additional models from 7 independent training pipelines, bringing the total to 12 models from 11 independent training pipelines. The central result is that both navigational structure and navigational content generalize across architectures --- with one scale-dependent exception in content generalization, alongside separate measurement caveats discussed in Sections 5 and 9.

\subsection{Replication Across 12 Models}\label{replication-across-12-models}

Phase 10A tested 5 models (Qwen 3.5 397B-A17B, MiniMax M2.5, Kimi K2.5, GLM 5, Llama 3.1 8B Instruct) across the same 12 concept pairs used in earlier phases. Four passed the reliability gate; GLM 5 was blocked by upstream rate limiting and excluded. Phase 11A tested 4 additional models (DeepSeek V3.2, Mistral Large 3, Cohere Command A, Llama 4 Maverick), all of which passed reliability.

\textbf{Methodological note.} The initial Phase 10A run used 60-second timeouts and blocked 4 of 5 models on OpenRouter latency, not model capability --- all four parsed correctly with 100\% connectivity. After relaxing timeouts to 300 seconds via a \texttt{-\/-patient} CLI flag and re-running overnight, 4 of 5 models passed. This infrastructure episode is instructive: aggressive timeout thresholds can masquerade as model failures, and a consequential methodological change for the benchmark's generality results was a timeout relaxation, not a prompt revision.

\textbf{R1 (gait) replicates universally.} All 12 models show characteristic conceptual gaits --- stable, model-specific consistency levels measured by intra-model Jaccard similarity. The full spectrum:

\begin{longtable}[]{@{}lrl@{}}
\toprule\noalign{}
Model & Gait (Mean Jaccard) & Cohort \\
\midrule\noalign{}
\endhead
\bottomrule\noalign{}
\endlastfoot
Mistral Large 3 & 0.747 & Phase 11A \\
Claude Sonnet 4.6 & 0.578 & Original \\
DeepSeek V3.2 & 0.540 & Phase 11A \\
Llama 4 Maverick & 0.539 & Phase 11A \\
Qwen 3.5 397B-A17B & 0.508 & Phase 10A \\
Cohere Command A & 0.502 & Phase 11A \\
MiniMax M2.5 & 0.419 & Phase 10A \\
Kimi K2.5 & 0.414 & Phase 10A \\
Gemini 3 Flash & 0.372 & Original \\
Llama 3.1 8B Instruct & 0.298 & Phase 10A \\
Grok 4.1 Fast & 0.293 & Original \\
GPT-5.2 & 0.258 & Original \\
\end{longtable}

The range spans 2.9× from GPT (0.258) to Mistral (0.747). Gait is not bounded at Claude's level (\textasciitilde0.58); architectural and training differences can produce substantially more rigid navigation. The property is universal --- every model tested has a measurable, characteristic gait --- but the specific value is model-specific.

\textbf{R2 (asymmetry) replicates universally.} All 12 models exceed the 0.60 directional asymmetry threshold: Phase 10A models ranged from 0.638 (MiniMax) to 0.785 (Llama 8B); Phase 11A models from 0.673 (Llama 4 Maverick) to 0.729 (Mistral). The quasimetric property --- that forward and reverse paths share few waypoints --- is universal across models from 11 independent training pipelines. The resolution-dependence caveat from Section 5.1 (asymmetry falls below threshold at 5 waypoints in the Phase 11C robustness subset) applies to the measurement, not the phenomenon.

\subsection{The Structure/Content/Scale Hierarchy}\label{the-structurecontentscale-hierarchy}

The most informative result from the generality phases is not that structure replicates --- that was expected --- but that navigational \emph{content} (which specific bridge concepts models use) also mostly generalizes. The hierarchy:

\textbf{Structure generalizes across tested models under the primary elicitation protocol.} Gait and asymmetry are universal properties across all 12 models from 11 independent training pipelines. No model fails either property, though asymmetry detectability requires sufficient waypoint count (Section 5.1).

\textbf{Content mostly generalizes among large models.} Aggregate bridge frequency for the benchmark's pre-specified bridges is statistically indistinguishable between cohorts, with some bridges near-universal and others remaining model-dependent. Phase 10A: new cohort mean 0.721 vs original 0.817, diff −0.096, CI including zero. Phase 11A combined: new cohort 0.717 vs original 0.817, diff −0.100, CI including zero. The strongest bridges --- ``spectrum'' for light→color, ``dog'' for animal→poodle, ``warm'' for hot→cold --- appear consistently across providers, but others show substantial heterogeneity: ``sadness'' ranges from 0.000 to 1.000 across models, ``harmony'' from 0.133 to 0.600, and ``leather'' from 0.267 to 1.000. The aggregate pattern is convergence; the bridge-level pattern is a mixture of universal and model-dependent landmarks.

\textbf{Scale differentiates.} Llama 3.1 8B Instruct is the sole outlier in bridge frequency: mean 0.200, compared to 0.717--0.817 for all frontier-scale models. Phase 11A provides a within-family comparison: Llama 4 Maverick (frontier-scale) shows bridge frequency of 0.724, while Llama 3.1 8B shows 0.200 --- a 3.6× difference within the same model family. Both models have navigational structure (gait and asymmetry are present), but the small model navigates through different landmarks. Scale is the clearest observed differentiator in current data, though the threshold and the role of architecture remain open.

\textbf{Connection to \citet{karkada2025translation}.} Karkada et al.~proved that geometric structure in LLM representations arises inevitably from translation symmetry in co-occurrence statistics --- structure that all models share because they train on similar data distributions. Our empirical finding of structure/content generality is consistent with this prediction: if geometric structure is determined by data statistics, then navigational structure derived from that geometry should also generalize. The scale effect suggests a minimum capacity threshold for extracting the full navigational structure from shared data statistics.

\textbf{Connection to the Platonic Representation Hypothesis.} The Platonic Representation Hypothesis \citep{huh2024platonic} predicts that larger models converge toward shared representations, supported by measures like CKA alignment. Our behavioral test partially confirms this: navigational \emph{structure} converges (all models have gait and asymmetry), navigational \emph{landmarks} converge for large models (same bridge concepts), but navigational \emph{gaits} remain model-specific (the 2.9× gait range does not collapse with scale). The models share a map but prefer different routes --- partial confirmation of representational convergence, with a clear boundary where convergence stops.

\subsection{Gait Profiles of New Models}\label{gait-profiles-of-new-models}

Phase 11A extended the gait spectrum beyond the original four models, revealing that the range is wider than initially characterized.

\textbf{Mistral sets the record at 0.747.} Mistral Large 3 showed near-deterministic navigation on several pairs: 0.936 on music→mathematics (producing nearly identical waypoint sets across all repetitions) and 0.934 on hot→cold. This exceeds Claude's 0.578 by a substantial margin and demonstrates that rigid, near-deterministic navigation is an achievable endpoint of the gait spectrum.

\textbf{DeepSeek, Cohere, and Llama 4 Maverick cluster at moderate consistency.} DeepSeek V3.2 at 0.540, Llama 4 Maverick at 0.539, and Cohere Command A at 0.502 occupy the middle of the gait spectrum --- more consistent than GPT/Grok but less rigid than Claude/Mistral. The Llama within-family comparison is particularly informative: Maverick (0.539) vs 8B (0.298) shows that gait increases with scale within a model family, suggesting that scale may contribute to navigational consistency within the Llama family, though broader causal attribution remains open.

The 12-model gait spectrum reveals that architectural and training differences produce qualitatively different navigation strategies. This is not merely quantitative variation --- a model with gait 0.747 (Mistral) generates near-identical paths across runs, while a model with gait 0.258 (GPT) explores broadly across a large waypoint vocabulary. The gait metric captures a genuine behavioral dimension that spans the full range from near-deterministic to broadly exploratory.

\begin{center}\rule{0.5\linewidth}{0.5pt}\end{center}

The generality results establish that the benchmark's structural claims are not artifacts of the original four-model cohort. Gait and asymmetry are universal across 12 models from 11 independent training pipelines. Frontier-scale models share the benchmark's strongest bridges and similar aggregate bridge frequencies, but several bridges remain model-dependent. The scale effect provides a meaningful boundary condition: frontier-scale models converge on the same geometric structure and share the strongest navigational landmarks, while small models retain the structure but navigate through different waypoints. The next section tests whether these findings are also robust to protocol variation.
       % Section 8: Act V -- Generality
\input{09-robustness}       % Section 9: Act VI -- Robustness
\section{Discussion}\label{discussion}

The preceding six acts established that LLM conceptual navigation has measurable geometric structure (Section 4), consistent with quasimetric topology (Section 5), causally sensitive to early anchoring (Section 6), resistant to single-variable mechanistic explanation (Section 7), generalizable across 12 models from 11 independent training pipelines (Section 8), and robust to protocol variation (Section 9). This section synthesizes those findings, states their limitations directly, and identifies the directions they open. Across approximately 21,540 runs, 12 models, 11 phases, and 29 documented dead ends, the benchmark reveals a consistent empirical picture: conceptual navigation is structured, asymmetric, model-specific in dynamics, and opaque in mechanism.

\subsection{Summary of the Six-Act Narrative}\label{summary-of-the-six-act-narrative}

\textbf{Structure (Section 4).} Models have distinct, stable conceptual gaits spanning a 2.9x range from GPT (0.258) to Mistral (0.747) {[}robust{]} R1. Navigation is self-consistent across time and resolution --- cross-batch Jaccard falls within 0.05 of within-batch (O9), and 5-waypoint paths appear as proper subsequences of 10-waypoint paths 67.9\% of the time (O10). The dual-anchor effect confirms that both endpoints shape the trajectory, ruling out pure forward-chaining (O2).

\textbf{Topology (Section 5).} Conceptual space is consistent with quasimetric structure. The symmetry axiom fails comprehensively --- mean directional asymmetry is 0.811 in the Phase 2 cohort, exceeding 0.60 for all 12 models tested {[}robust{]} R2. The triangle inequality holds at approximately 91\% across three independent samples, tested on within-direction Jaccard distance. These two properties are measured on different distance operationalizations (cross-direction Jaccard for asymmetry, within-direction Jaccard for triangle inequality), so the quasimetric characterization is convergent rather than formally established on a single unified metric. Bridge concepts function as structural bottlenecks --- obligatory intermediaries rather than mere associations {[}robust{]} R6. Hierarchical paths are compositional, with 4.9x higher transitivity than random controls {[}robust{]} R4.

\textbf{Mechanism (Section 6).} Pre-filling the first waypoint causally displaces bridges (O15, mean displacement 0.515, CI {[}0.357, 0.664{]} excluding zero). The primary mechanism is associative primacy: whatever occupies position 1 anchors the trajectory regardless of its semantic content, with displacement magnitudes statistically indistinguishable across incongruent, congruent, and neutral conditions. Relation class matters as a secondary modulator (O26): related pre-fills (on-axis or same-domain) preserve the navigational context and allow bridge recovery, while unrelated pre-fills destroy it (Friedman p = 0.034). Bridge survival under perturbation is bimodal --- ``sadness'' and ``dog'' survive robustly; ``harmony'' and ``germination'' collapse to near-zero.

\textbf{Limits (Section 7).} Seven follow-up hypotheses were falsified --- spanning mechanistic, model-deficit, methodological, and effect-direction types; one initially falsified result was resurrected with additional data. Prediction accuracy descended from 81.3\% in the characterization phases to 20-24\% in the mechanism phases (O24). The pattern is systematic: qualitative structure is characterizable and replicable at approximately 80\%; quantitative mechanism is opaque, with single-variable models succeeding at approximately 0-15\%.

\textbf{Generality (Section 8).} Gait and asymmetry replicate universally across 12 models from 11 independent training pipelines {[}robust{]} R1, R2. Bridge frequency generalizes among frontier-scale models (aggregate difference CI includes zero) but not to Llama 3.1 8B, which navigates through different landmarks (bridge frequency 0.200 vs 0.717-0.817 for frontier models). The within-family Llama comparison is particularly informative: Maverick (frontier-scale) shows bridge frequency 0.724 while 8B shows 0.200, a 3.6x difference. The hierarchy is structure \textgreater{} content \textgreater{} scale: structural properties are universal, navigational content is shared among large models, and scale is the clearest differentiator {[}robust{]} R6 extended.

\textbf{Robustness (Section 9).} Model identity dominates protocol variation in an ANOVA decomposition (eta-squared = 0.242 for model vs approximately 0 for waypoint count and temperature) (O30). Bridge frequency is the most robust property, exceeding 0.97 across all conditions including the most challenging (5 waypoints, temperature 0.9) (O31). Gait rank ordering is largely stable (Kendall's W = 0.840), with Claude maintaining the top position in every condition. Asymmetry is resolution-dependent: 9-waypoint conditions clear the 0.60 threshold (0.669-0.684) while 5-waypoint conditions fall just short (0.593-0.594) (O32). Temperature effects on all metrics are negligible.

\subsection{What the Benchmark Reveals About Conceptual Space}\label{what-the-benchmark-reveals-about-conceptual-space}

Four interpretive points emerge from the empirical results.

\textbf{Consistent with quasimetric structure.} Symmetry fails comprehensively (mean 0.811). The triangle inequality holds at approximately 91\% across three independent samples. These two findings use different operationalizations of distance --- cross-direction Jaccard for asymmetry, within-direction Jaccard for triangle inequality --- so the formal quasimetric claim is ``consistent with'' rather than deductively established on a single metric. The asymmetry finding parallels two independent results from different literatures: \citeauthor{tversky1977features}'s (\citeyear{tversky1977features}) demonstration that human similarity judgments are inherently asymmetric (North Korea is judged more similar to China than China is to North Korea), and \citeauthor{berglund2024reversal}'s (\citeyear{berglund2024reversal}) reversal curse showing that autoregressive models cannot reverse learned associations. Our finding is consistent with both: navigational asymmetry parallels the reversal curse and Tversky's human similarity asymmetry. The connection is thematic --- all three concern directional asymmetry in conceptual structure --- rather than mechanistic; we do not demonstrate that the reversal curse causes the specific asymmetry we observe.

\textbf{Bottleneck topology.} The most reliable navigational landmarks are not the most strongly associated concepts but obligatory intermediaries: ``spectrum'' for light-to-color (frequency 1.00 across all four original models), ``dog'' for animal-to-poodle, ``warm'' for hot-to-cold. ``Metaphor'' --- intuitively central to the language-to-thought relationship --- appeared at 0.00 frequency (graveyard entry G5). ``Energy'' --- associated with both hot and cold --- appeared at 0.00 on hot-to-cold. Process-naming concepts outperformed object-naming concepts: ``germination'' outperformed ``plant'' for seed-to-garden in all four models. This echoes a finding from knowledge graph research: shortest-path intermediaries are more informative than high-degree hubs. Navigational salience depends on directional information content, not on generic associative strength.

\textbf{Model-specific gaits on shared geometry.} All frontier-scale models navigate the same basic geometry through the same strongest landmarks, but with different route preferences. Mistral (gait 0.747) produces near-deterministic paths --- 0.936 Jaccard on music-to-mathematics, producing nearly identical waypoint sets across all repetitions. GPT (0.258) explores broadly, sampling different landmarks on each traversal. The aggregate bridge frequency difference between original and expanded cohorts has a CI including zero, yet gait varies 2.9x. This is the behavioral analog of a refinement to the Platonic Representation Hypothesis: representations may converge \citep{huh2024platonic}, but navigation diversifies. Same map, different routes. Structure converges; gaits remain model-specific. Neither extreme is intrinsically superior --- they represent genuinely different navigational strategies that a model-agnostic benchmark must accommodate.

\textbf{Navigational structure is not compositional from simple variables.} The mechanism ceiling --- 7 of 8 hypotheses falsified at primary test --- suggests that navigation emerges from the interaction of multiple factors rather than any single structural variable. Competitor count (G20), dominance ratio (G23), gradient type (G21, G24), and gait normalization (G22) all failed individually. The prediction accuracy descent from 81.3\% to 20\% tracks a shift in prediction type, not a loss of signal: replication predictions succeed at approximately 80\% throughout, while mechanistic predictions hover near 0-15\%. This parallels findings in computational neuroscience, where spatial navigation in rodents and humans depends on multi-factorial interactions between grid cells, place cells, head-direction cells, and boundary cells. The analogy is structural, not mechanistic: in both cases, the emergent navigational behavior resists decomposition into single-variable explanations.

\textbf{Statistical caveat on the quasimetric characterization.} The two tests supporting the quasimetric claim --- asymmetry and triangle inequality --- use different distance operationalizations. Asymmetry uses cross-direction Jaccard (comparing A-to-B paths against B-to-A paths); triangle inequality uses within-direction Jaccard (consistency of repeated same-direction runs). Additionally, all metrics are lexical --- synonymous waypoints are treated as distinct. The magnitude of the resulting bias is bounded by the canonicalization pipeline but is otherwise unknown. These caveats do not undermine the qualitative finding (systematic asymmetry, triangle inequality holding at approximately 91\%), but the formal quasimetric claim should be understood as ``consistent with'' rather than ``formally established.''

\subsection{Limitations}\label{limitations}

The benchmark's limitations divide into two distinct categories: design issues with the benchmark itself, and findings about the phenomenon that constrain interpretation. Conflating these would obscure the distinction between things we should have done better and things that turned out to be harder than expected.

\subsubsection{Benchmark Limitations}\label{benchmark-limitations}

\textbf{1. Control pair failure {[}robust{]} R5 (qualified).} This is the benchmark's most serious design limitation. The original control pair stapler-monsoon fails strict criteria for all 12 models (top waypoint frequency 0.650-1.000). Phase 11B tested 4 replacement candidates (accordion-stalactite, turmeric-trigonometry, barnacle-sonnet, magnesium-ballet) across 6 models: 0 of 24 cells passed either the frequency gate (\textless{} 0.15) or the entropy gate (\textgreater{} 5.0). Each candidate revealed creative structured routing --- turmeric-trigonometry produced ``ancient India'' (Claude), ``spice'' (GPT), ``golden ratio'' (Mistral) --- yet each model still converged on a single dominant waypoint. LLMs find navigable routes between any concept pair at 7 waypoints, defeating single-pair strict-threshold control validation. The benchmark's validity rests on the predictability gap --- experimental bridges are predicted from the endpoint relationship; control bridges are ad hoc --- not on presence versus absence of structure. This limitation should be weighted prominently in any evaluation of the benchmark's claims.

\textbf{2. Self-grounding circularity.} Model outputs generate their own ground truth. The concern is that structural findings could be artifacts of the generation process rather than properties of conceptual organization. Three mitigations partially address this:

\begin{enumerate}
\def\labelenumi{(\alph{enumi})}
\tightlist
\item
  Testing against external mathematical properties that have model-independent definitions (metric axioms --- the triangle inequality either holds or it does not, independent of the generation process).
\item
  Cross-model comparison across 12 independent models (if all 12 show the same property, it cannot be attributed to idiosyncrasies of any single model's decoding).
\item
  Causal intervention via the pre-fill experiments (bridge displacement demonstrates that structure is causally sensitive, not merely a statistical regularity).
\end{enumerate}

These mitigations reduce the concern but never fully eliminate it. The benchmark cannot distinguish between ``models have navigational geometry'' and ``the generation process has statistical regularities that mimic navigational geometry'' --- though the causal sensitivity of bridge displacement (Section 6.1) makes the second interpretation less parsimonious.

\textbf{3. Lexical metrics only.} All metrics use Jaccard similarity after lemmatization. ``Melody'' and ``tune'' are treated as entirely distinct. The asymmetry mean of 0.811 is lexical asymmetry; true conceptual asymmetry is likely lower because some apparent non-overlap reflects synonymous waypoints scored as different. The magnitude of this bias is unknown but bounded by the canonicalization pipeline, which handles morphological variants (running/run, cities/city).

Embedding-based semantic similarity was implemented but deferred due to API cost. This is the benchmark's most pervasive methodological limitation, present in every metric and every phase. It affects all structural claims: gait underestimates true conceptual consistency, asymmetry overestimates true conceptual asymmetry, and bridge frequency may miss bridges expressed through synonyms rather than exact tokens.

\textbf{4. API-mediated access.} Internal model state cannot be controlled, temperature settings may be approximate for some providers, and silent model updates may occur during multi-day data collection. Partially addressed by the no-temporal-drift finding (O9, cross-batch Jaccard within 0.05 of within-batch) and the robustness to temperature variation (Phase 11C, eta-squared = 0.002 for temperature). These mitigations apply to the tested conditions; they do not rule out sensitivity to untested API-level changes.

\textbf{5. Concept pair selection bias.} Pairs were selected by the experimenters based on semantic interest and hypothesis-driven design, not sampled from a principled distribution over conceptual space. The benchmark tests navigational structure on the pairs it tests, not on a representative sample of all possible concept relationships. Partially addressed by the breadth of the pair inventory (100+ unique pairs across 11 phases) and the generality analysis (Section 8), but the selection is not random and the results should not be extrapolated to arbitrary concept pairs without further testing.

\textbf{6. Resolution dependence.} Asymmetry requires 9 or more waypoints to reliably exceed the 0.60 threshold in the robustness subset (O32): the 5-waypoint conditions fall just short (0.593-0.594), the 7-waypoint baseline also fell short (0.599) in the sparse 3-model subset, and only the 9-waypoint conditions clear it (0.669-0.684). The full Phase 2 cohort at 5 waypoints produced 0.811, but that used a larger pair set and model cohort. Other structural properties may have similar resolution thresholds that have not been identified --- the benchmark tested waypoint counts of 5, 7, and 9 but not the full range. This is a measurement sensitivity finding, not evidence against asymmetry, but it means that claims about asymmetry magnitude are conditioned on sufficient path length.

\subsubsection{Phenomenon Limitations}\label{phenomenon-limitations}

\textbf{7. Mechanism ceiling.} Seven of eight hypotheses were falsified at the primary-test level (Section 7.2). This is not a design flaw --- it is the benchmark's most informative negative result. The mechanism ceiling establishes that conceptual navigation resists single-variable explanation and that the resolving power of one-factor experimental designs has been exhausted for this phenomenon. The prediction accuracy trajectory --- 81.3\% (Phase 4) to 20\% (Phase 9) --- tracks a systematic shift in prediction type, not random degradation.

\textbf{8. Cross-model distance failure (O18, O19).} Models rank-order navigational distances differently. The cross-model distance correlation was r = 0.170, far below the 0.50 validity threshold. Gait normalization produced zero improvement (r = 0.212 before and after, with 0.000 change on every model pair, G22). Some model pairs actively anti-correlate: Grok-Gemini at r = -0.580. No universal distance function is recoverable from path-based measurements. Within-model distance analysis remains viable; cross-model distance comparison is blocked.

\textbf{9. Gemini characterization failure.} Gemini's approximately 30\% bridge frequency deficit relative to the non-Gemini mean (0.480 vs approximately 0.67) is real and stable across the tested phases. Three successive mechanistic characterizations --- frame-crossing threshold (G7), gradient blindness (G21), transformation-chain blindness (G24) --- all failed, with two producing results in the opposite direction from predicted. Whatever pair type is designated as ``hard'' for Gemini, Gemini does relatively better on it. The type-based approach to characterizing the deficit appears fundamentally misguided; the deficit is model-level, not type-level. Future approaches might require representational analysis rather than behavioral type-classification.

\subsection{The Graveyard as Methodology}\label{the-graveyard-as-methodology}

The benchmark documents 29 dead ends across 11 phases --- a transparent record of the benchmark's learning curve.

Four lessons emerge from the graveyard.

First, intuitive semantic importance has zero predictive power for navigational behavior. ``Metaphor'' --- the concept most intuitively connected to both language and thought --- appeared at 0.00 frequency as a bridge for language-to-thought (G5). ``Energy'' --- associated with both hot and cold --- appeared at 0.00 on hot-to-cold. Navigational salience and human intuition about conceptual centrality are orthogonal; bridges must name the primary axis of connection, not merely inhabit the shared associative neighborhood.

Second, qualitative models can be correct while their quantitative predictions fail. Gemini's deficit is real (approximately 30\% reduction in bridge frequency), but three quantitative characterizations of its mechanism failed (G7, G21, G24). The qualitative observation replicates across every phase in which Gemini was tested; the mechanism remains opaque. This pattern --- robust qualitative finding, intractable quantitative mechanism --- characterizes the benchmark as a whole.

Third, retrospective signals on small samples do not generalize. The dominance ratio showed a promising correlation (rho = 0.60) on the original 8-pair sample but collapsed to rho = 0.157 on an independent sample (G23). The competitor count predictor showed a similar failure (G20, rho = 0.116). Both variables were plausible on the original data; neither survived replication.

Fourth, infrastructure limitations can masquerade as substantive findings. The bridge generalization failure (G26) was initially attributed to the phenomenon based on data from a single small model (Llama 8B) under aggressive 60-second timeouts that excluded most new models on latency, not capability. After timeout relaxation to 300 seconds and testing of 3 additional frontier-scale models, the aggregate finding reversed --- bridge frequency CI shifted from excluding zero to including it. Llama 8B's low bridge frequency (0.200) remains real as a scale effect, but the original conclusion that bridge structure does not generalize was wrong. G26 is the sole clean resurrection in the graveyard and a caution against both premature interpretation of negative results and confounding infrastructure limitations with phenomenon limitations.

Transparent dead-end documentation should be standard practice in benchmarking work. The incentive structure in machine learning research favors reporting only successful experiments, but the graveyard entries in this benchmark --- particularly G5 (metaphor), G22 (gait normalization), and G26 (resurrected) --- are as methodologically informative as the successful findings. We recommend that future benchmark papers adopt a similar practice of documenting falsified hypotheses alongside confirmed ones.

\subsection{Connections to Theory}\label{connections-to-theory}

\textbf{\citet{karkada2025translation}.} They proved that static geometric structure in LLM representations arises inevitably from translation symmetry in co-occurrence statistics --- the manifold shape is determined by the data distribution, not by architectural choices. Our benchmark tests the dynamic complement: whether that geometry supports structured, consistent behavioral navigation. Together, the contributions are complementary.

Their result explains why static geometry exists; ours shows that it supports navigation that is structured and asymmetric at the universal level, but model-specific in its dynamics. Geometry comes from data statistics (their contribution); navigational behavior on that geometry is universal in structure but model-specific in gait (ours). The scale effect from Section 8 is consistent with their framework: if geometric structure requires sufficient capacity to extract translation symmetries from training data, smaller models (Llama 8B) may have the structure but lack the resolution to navigate through the same landmarks as frontier models.

\textbf{Gardenfors operationalized.} \citeauthor{gardenfors2000}'s (\citeyear{gardenfors2000}) conceptual spaces theory posits that concepts live in navigable geometric spaces with measurable distance. The benchmark provides a large-scale computational probe of this theory across 12 models and 100+ concept pairs. The partial confirmation is substantive: concepts do live in navigable spaces with measurable, compositional distance (triangle inequality approximately 91\%, hierarchical transitivity 4.9x over random). Polysemy sense differentiation (cross-pair Jaccard near zero for different senses of ``bank,'' ``bat,'' and ``crane'') confirms that paths reflect conceptual structure rather than surface word statistics. But the pervasive asymmetry (mean 0.811) challenges the standard metric assumption in conceptual spaces theory, which typically presumes symmetric distance functions. The data suggest that quasimetric structure --- metric-like but inherently directional --- may be a more appropriate formal framework for computational conceptual spaces.

\textbf{Platonic Representation Hypothesis refined.} \citet{huh2024platonic} predicted that larger models converge toward shared representations. Our behavioral test partially confirms this: navigational structure converges (all 12 models have gait and asymmetry), navigational landmarks converge for frontier-scale models (aggregate bridge frequency difference CI includes zero), but navigational gaits remain model-specific (the 2.9x gait range does not collapse with scale). The result establishes a boundary for the hypothesis: representational convergence does not imply behavioral convergence. Structure converges; gaits diversify. Partial confirmation with a clear boundary condition.

\textbf{\citet{wyss2025spectral}.} Spectral analysis of semantic attractors predicts discrete basins in LLM conceptual space. Our bridge bottlenecks and the ``deep basins'' observed in the word convergence game that inspired this benchmark \citep{hodge2025convergence} are plausibly analogous to those spectral attractors --- the behavioral manifestation of the basin structure their theory predicts. Bridge concepts like ``spectrum'' and ``dog'' behave as attractor basins --- navigational points that paths converge toward regardless of starting trajectory. The pre-fill displacement results (Section 6.1) demonstrate that these basins are causally active: they anchor paths when present and paths reorganize when they are displaced. The bridge frequency robustness finding (\textgreater{} 0.97 across all protocol conditions, O31) further supports the attractor interpretation --- these are deep structural features of the navigational landscape, not artifacts of a particular elicitation protocol.

\subsection{Future Work}\label{future-work}

The benchmark's findings and limitations suggest ten specific directions, ordered from the most directly motivated by current data to the most speculative.

\begin{enumerate}
\def\labelenumi{\arabic{enumi}.}
\item
  \textbf{Multi-variable bridge fragility model.} The mechanism ceiling (Section 7) was reached by testing single variables in isolation. A joint model --- bridge role x pre-fill content x dominance ratio x competitor count --- could capture the interaction effects that single-variable designs miss. The bimodal survival pattern (Section 6.1) suggests that at least two factors interact: ``sadness'' survives because it has no competitors and sits on an obligatory navigational axis, while ``harmony'' collapses because it competes with rhythm, pattern, and frequency in a crowded navigational space. The combinatorial cost is high (likely 5,000+ runs for a 4-factor design), but this is the only known path past the mechanism ceiling.
\item
  \textbf{Embedding-based distance validation.} Replace Jaccard with embedding cosine similarity to rescue the lexical limitation (Section 10.3, limitation 3). This would also test whether the cross-model distance failure (O18, O19) is a property of lexical metrics or of the navigational distances themselves.
\item
  \textbf{Multilingual conceptual topology.} Test whether the navigational geometry holds for non-English prompts and multilingual models. If the structure arises from data statistics \citep{karkada2025translation}, language-specific training distributions may produce different geometries.
\item
  \textbf{Scale threshold investigation.} The Llama family provides a natural within-family comparison: 8B (bridge frequency 0.200) vs Maverick at frontier scale (0.724). Intermediate checkpoints (Llama 70B, Llama 405B) would locate the scale threshold at which bridge content converges with frontier models and determine whether the transition is gradual or abrupt.
\item
  \textbf{Temporal stability under model updates.} The no-drift finding (O9) covers the data collection window. Longitudinal monitoring --- re-running the core pair set on the same model IDs at 3-month intervals --- would test whether navigational structure is preserved across training updates and provider-side model swaps.
\item
  \textbf{Human waypoint comparison.} A calibration subset of 10-15 concept pairs administered to human participants under the same prompt format would establish whether the structural properties (asymmetry, bridge bottlenecks, gait variation) are shared between human and LLM navigation or are specific to autoregressive generation. Given \citeauthor{tversky1977features}'s (\citeyear{tversky1977features}) finding that human similarity is asymmetric and \citeauthor{constantinescu2016grid}'s (\citeyear{constantinescu2016grid}) evidence for grid-cell-based conceptual navigation, partial overlap is plausible but untested.
\item
  \textbf{External anchoring.} Validate bridge predictions against ConceptNet relation strength, Small World of Words association norms, or Wikipedia link structure. This would partially address the self-grounding circularity (limitation 2) by grounding navigational landmarks in independently measured conceptual structure.
\item
  \textbf{Control pair redesign.} The single-pair strict-threshold approach has failed (limitation 1). Three alternatives: (a) relative consistency metrics comparing experimental vs control pairs on predictability rather than presence of structure; (b) pair batteries testing 20+ control pairs to characterize the distribution of spurious bridge frequencies; (c) difficulty-graded controls sampling from the full spectrum of conceptual distance.
\item
  \textbf{Chain-of-thought as intervention.} Eliciting waypoints with and without chain-of-thought reasoning would test whether explicit reasoning changes navigational structure or merely rationalizes the same implicit routes. This connects to the compositionality gap literature: if chain-of-thought improves multi-hop reasoning, it may also alter bridge selection, path composition, and the gait spectrum.
\item
  \textbf{Curvature estimation.} Within-model curvature is estimable from the triangle inequality data --- the distribution of triangle excess (d(A,B) + d(B,C) - d(A,C)) characterizes local geometry, with positive excess suggesting positive curvature and near-zero excess suggesting flatness. Cross-model curvature comparison is blocked by the distance failure (O18, O19) and should not be attempted until embedding-based distances are validated (item 2 above).
\end{enumerate}

\begin{center}\rule{0.5\linewidth}{0.5pt}\end{center}

The discussion synthesizes a benchmark that achieves high replication accuracy (approximately 80\%) for structural characterization while reaching a mechanism ceiling (approximately 15-20\%) for quantitative explanation. The core empirical contribution is a set of structural properties --- characteristic gaits spanning 2.9x across 12 models, quasimetric asymmetry at 0.811, bottleneck bridge topology with frequencies exceeding 0.97 under robustness testing, and compositional hierarchy with 4.9x transitivity over random controls --- that generalize across architectures and survive protocol variation. The core negative contribution is equally important: 7 falsified mechanistic hypotheses and 29 documented dead ends establish that the qualitative geometry of conceptual navigation is accessible to current methods while the quantitative mechanism is not. The benchmark's most serious limitation --- the control pair failure, with 0 of 24 replacement cells passing --- is acknowledged directly, and its validity is reframed as resting on predictability gaps rather than absolute baselines. The six theoretical connections --- to quasimetric topology, Tversky's asymmetry, Karkada et al.'s geometric inevitability, Gardenfors's conceptual spaces, the Platonic Representation Hypothesis, and spectral semantic attractors --- position the findings within a broader landscape where static geometry is well-characterized but dynamic navigation is only beginning to be understood. The most productive future directions involve multi-variable mechanism models, embedding-based distance validation, and human comparison --- each addressing a specific limitation identified across the six acts.
       % Section 10: Discussion
\section{Conclusion}\label{conclusion}

Across approximately 21,540 runs spanning 12 models from 11 independent training pipelines, the Conceptual Topology Mapping Benchmark establishes that LLMs have consistent, measurable geometric structure in how they navigate between concepts. This structure is consistent with quasimetric geometry --- systematically asymmetric (mean 0.811), with the triangle inequality holding at approximately 91\%. Navigation is organized around bottleneck bridges that function as structural landmarks, not mere associations, and expressed through model-specific gaits spanning a 2.9x range from GPT (0.258) to Mistral (0.747). These findings generalize across all 12 models and resist protocol variation.

The mechanism ceiling is itself the most important finding. Seven follow-up hypotheses --- spanning mechanistic, model-deficit, methodological, and effect-direction types --- were falsified across Phases 8--10; one was resurrected with additional data. Prediction accuracy descended from 81.3\% for characterization questions to 20--24\% for mechanistic ones --- not because signal disappeared, but because qualitative structure (characterizable at \textasciitilde80\%) resists quantitative mechanistic decomposition (opaque at 0--15\%). The dynamics generating navigational regularity do not reduce to any individual variable tested. This parallels the broader challenge in interpretability: structure is detectable before it is explainable.

Static embedding geometry is well-studied; navigation on that geometry is not. \citet{karkada2025translation} proved why the geometric structure exists --- translation symmetry in co-occurrence statistics determines manifold shape. We showed what that geometry supports: structured, asymmetric, model-specific behavioral navigation. The Conceptual Topology Mapping Benchmark provides a different evaluation paradigm --- not testing what models know, but how they navigate what they know. The 29 documented dead ends and the mechanism ceiling are not failures of the research program but honest accounting of where the field's explanatory tools fall short, and where the next generation of multi-factorial mechanistic models will need to reach.
       % Section 11: Conclusion

% =============================================================================
% Bibliography
% =============================================================================

\bibliography{references}

% =============================================================================
% Appendices
% =============================================================================

\appendix

\section{Concept Pair and Triple Inventory}\label{appendix-a-concept-pair-and-triple-inventory}

This appendix provides the complete inventory of concept pairs, triples, and phase-specific stimuli used across the benchmark's 11 experimental phases.

\subsection{A.1 Core Concept Pairs (Phase 1)}\label{a.1-core-concept-pairs-phase-1}

36 pairs across 9 categories, split into holdout (15, used only for prompt selection) and reporting (21, used for all published analyses).

\subsubsection{Anchor Pairs}\label{anchor-pairs}

Known basin structure from the word convergence game experiments.

\begin{longtable}[]{@{}
  >{\raggedright\arraybackslash}p{(\columnwidth - 12\tabcolsep) * \real{0.0870}}
  >{\raggedright\arraybackslash}p{(\columnwidth - 12\tabcolsep) * \real{0.1304}}
  >{\raggedright\arraybackslash}p{(\columnwidth - 12\tabcolsep) * \real{0.0870}}
  >{\raggedright\arraybackslash}p{(\columnwidth - 12\tabcolsep) * \real{0.2826}}
  >{\raggedright\arraybackslash}p{(\columnwidth - 12\tabcolsep) * \real{0.1522}}
  >{\raggedright\arraybackslash}p{(\columnwidth - 12\tabcolsep) * \real{0.1522}}
  >{\raggedright\arraybackslash}p{(\columnwidth - 12\tabcolsep) * \real{0.1087}}@{}}
\toprule\noalign{}
\begin{minipage}[b]{\linewidth}\raggedright
ID
\end{minipage} & \begin{minipage}[b]{\linewidth}\raggedright
From
\end{minipage} & \begin{minipage}[b]{\linewidth}\raggedright
To
\end{minipage} & \begin{minipage}[b]{\linewidth}\raggedright
Concreteness
\end{minipage} & \begin{minipage}[b]{\linewidth}\raggedright
Basin
\end{minipage} & \begin{minipage}[b]{\linewidth}\raggedright
Depth
\end{minipage} & \begin{minipage}[b]{\linewidth}\raggedright
Set
\end{minipage} \\
\midrule\noalign{}
\endhead
\bottomrule\noalign{}
\endlastfoot
anchor-beyonce-erosion & Beyonce & erosion & concrete--abstract & formation & deep & reporting \\
anchor-tesla-mycelium & Tesla & mycelium & concrete--concrete & network & deep & reporting \\
anchor-shadow-melody & shadow & melody & concrete--abstract & echo & moderate & holdout \\
anchor-skull-garden & skull & garden & concrete--concrete & --- & moderate & reporting \\
anchor-kanye-lattice & Kanye West & lattice & concrete--abstract & --- & flat & holdout \\
\end{longtable}

\subsubsection{Hierarchy Pairs}\label{hierarchy-pairs}

Hypernym--hyponym relationships at varying abstraction levels.

\begin{longtable}[]{@{}
  >{\raggedright\arraybackslash}p{(\columnwidth - 8\tabcolsep) * \real{0.1250}}
  >{\raggedright\arraybackslash}p{(\columnwidth - 8\tabcolsep) * \real{0.1875}}
  >{\raggedright\arraybackslash}p{(\columnwidth - 8\tabcolsep) * \real{0.1250}}
  >{\raggedright\arraybackslash}p{(\columnwidth - 8\tabcolsep) * \real{0.4062}}
  >{\raggedright\arraybackslash}p{(\columnwidth - 8\tabcolsep) * \real{0.1562}}@{}}
\toprule\noalign{}
\begin{minipage}[b]{\linewidth}\raggedright
ID
\end{minipage} & \begin{minipage}[b]{\linewidth}\raggedright
From
\end{minipage} & \begin{minipage}[b]{\linewidth}\raggedright
To
\end{minipage} & \begin{minipage}[b]{\linewidth}\raggedright
Concreteness
\end{minipage} & \begin{minipage}[b]{\linewidth}\raggedright
Set
\end{minipage} \\
\midrule\noalign{}
\endhead
\bottomrule\noalign{}
\endlastfoot
hierarchy-animal-poodle & animal & poodle & concrete--concrete & reporting \\
hierarchy-vehicle-skateboard & vehicle & skateboard & concrete--concrete & holdout \\
hierarchy-emotion-nostalgia & emotion & nostalgia & abstract--abstract & reporting \\
hierarchy-structure-lattice & structure & lattice & abstract--abstract & holdout \\
\end{longtable}

\subsubsection{Cross-Domain Pairs}\label{cross-domain-pairs}

Concepts from different semantic domains.

\begin{longtable}[]{@{}
  >{\raggedright\arraybackslash}p{(\columnwidth - 8\tabcolsep) * \real{0.1250}}
  >{\raggedright\arraybackslash}p{(\columnwidth - 8\tabcolsep) * \real{0.1875}}
  >{\raggedright\arraybackslash}p{(\columnwidth - 8\tabcolsep) * \real{0.1250}}
  >{\raggedright\arraybackslash}p{(\columnwidth - 8\tabcolsep) * \real{0.4062}}
  >{\raggedright\arraybackslash}p{(\columnwidth - 8\tabcolsep) * \real{0.1562}}@{}}
\toprule\noalign{}
\begin{minipage}[b]{\linewidth}\raggedright
ID
\end{minipage} & \begin{minipage}[b]{\linewidth}\raggedright
From
\end{minipage} & \begin{minipage}[b]{\linewidth}\raggedright
To
\end{minipage} & \begin{minipage}[b]{\linewidth}\raggedright
Concreteness
\end{minipage} & \begin{minipage}[b]{\linewidth}\raggedright
Set
\end{minipage} \\
\midrule\noalign{}
\endhead
\bottomrule\noalign{}
\endlastfoot
cross-music-mathematics & music & mathematics & abstract--abstract & reporting \\
cross-cooking-architecture & cooking & architecture & concrete--concrete & holdout \\
cross-justice-erosion & justice & erosion & abstract--concrete & reporting \\
cross-time-crystal & time & crystal & abstract--concrete & holdout \\
\end{longtable}

\subsubsection{Polysemy Pairs}\label{polysemy-pairs}

Same ambiguous word paired with different targets to steer sense activation.

\begin{longtable}[]{@{}
  >{\raggedright\arraybackslash}p{(\columnwidth - 10\tabcolsep) * \real{0.0833}}
  >{\raggedright\arraybackslash}p{(\columnwidth - 10\tabcolsep) * \real{0.1250}}
  >{\raggedright\arraybackslash}p{(\columnwidth - 10\tabcolsep) * \real{0.0833}}
  >{\raggedright\arraybackslash}p{(\columnwidth - 10\tabcolsep) * \real{0.3333}}
  >{\raggedright\arraybackslash}p{(\columnwidth - 10\tabcolsep) * \real{0.2708}}
  >{\raggedright\arraybackslash}p{(\columnwidth - 10\tabcolsep) * \real{0.1042}}@{}}
\toprule\noalign{}
\begin{minipage}[b]{\linewidth}\raggedright
ID
\end{minipage} & \begin{minipage}[b]{\linewidth}\raggedright
From
\end{minipage} & \begin{minipage}[b]{\linewidth}\raggedright
To
\end{minipage} & \begin{minipage}[b]{\linewidth}\raggedright
Polysemy Group
\end{minipage} & \begin{minipage}[b]{\linewidth}\raggedright
Target Sense
\end{minipage} & \begin{minipage}[b]{\linewidth}\raggedright
Set
\end{minipage} \\
\midrule\noalign{}
\endhead
\bottomrule\noalign{}
\endlastfoot
polysemy-bank-river & bank & river & bank & geographic & reporting \\
polysemy-bank-mortgage & bank & mortgage & bank & financial & holdout \\
polysemy-bat-cave & bat & cave & bat & animal & holdout \\
polysemy-bat-baseball & bat & baseball & bat & sports & reporting \\
polysemy-crane-construction & crane & construction & crane & machine & reporting \\
polysemy-crane-wetland & crane & wetland & crane & bird & holdout \\
\end{longtable}

\subsubsection{Near-Synonym Pairs}\label{near-synonym-pairs}

Dense semantic neighborhoods testing navigational friction.

\begin{longtable}[]{@{}
  >{\raggedright\arraybackslash}p{(\columnwidth - 8\tabcolsep) * \real{0.1250}}
  >{\raggedright\arraybackslash}p{(\columnwidth - 8\tabcolsep) * \real{0.1875}}
  >{\raggedright\arraybackslash}p{(\columnwidth - 8\tabcolsep) * \real{0.1250}}
  >{\raggedright\arraybackslash}p{(\columnwidth - 8\tabcolsep) * \real{0.4062}}
  >{\raggedright\arraybackslash}p{(\columnwidth - 8\tabcolsep) * \real{0.1562}}@{}}
\toprule\noalign{}
\begin{minipage}[b]{\linewidth}\raggedright
ID
\end{minipage} & \begin{minipage}[b]{\linewidth}\raggedright
From
\end{minipage} & \begin{minipage}[b]{\linewidth}\raggedright
To
\end{minipage} & \begin{minipage}[b]{\linewidth}\raggedright
Concreteness
\end{minipage} & \begin{minipage}[b]{\linewidth}\raggedright
Set
\end{minipage} \\
\midrule\noalign{}
\endhead
\bottomrule\noalign{}
\endlastfoot
synonym-cemetery-graveyard & cemetery & graveyard & concrete--concrete & reporting \\
synonym-forest-woods & forest & woods & concrete--concrete & holdout \\
synonym-happy-joyful & happy & joyful & abstract--abstract & reporting \\
synonym-anger-rage & anger & rage & abstract--abstract & holdout \\
\end{longtable}

\subsubsection{Antonym Pairs}\label{antonym-pairs}

Oppositional continua testing axis traversal.

\begin{longtable}[]{@{}lllll@{}}
\toprule\noalign{}
ID & From & To & Concreteness & Set \\
\midrule\noalign{}
\endhead
\bottomrule\noalign{}
\endlastfoot
antonym-hot-cold & hot & cold & concrete--concrete & reporting \\
antonym-order-chaos & order & chaos & abstract--abstract & holdout \\
\end{longtable}

\subsubsection{Control Pairs}\label{control-pairs}

\begin{longtable}[]{@{}
  >{\raggedright\arraybackslash}p{(\columnwidth - 8\tabcolsep) * \real{0.2308}}
  >{\raggedright\arraybackslash}p{(\columnwidth - 8\tabcolsep) * \real{0.1923}}
  >{\raggedright\arraybackslash}p{(\columnwidth - 8\tabcolsep) * \real{0.2308}}
  >{\raggedright\arraybackslash}p{(\columnwidth - 8\tabcolsep) * \real{0.1538}}
  >{\raggedright\arraybackslash}p{(\columnwidth - 8\tabcolsep) * \real{0.1923}}@{}}
\toprule\noalign{}
\begin{minipage}[b]{\linewidth}\raggedright
Type
\end{minipage} & \begin{minipage}[b]{\linewidth}\raggedright
ID
\end{minipage} & \begin{minipage}[b]{\linewidth}\raggedright
From
\end{minipage} & \begin{minipage}[b]{\linewidth}\raggedright
To
\end{minipage} & \begin{minipage}[b]{\linewidth}\raggedright
Set
\end{minipage} \\
\midrule\noalign{}
\endhead
\bottomrule\noalign{}
\endlastfoot
Identity & control-identity-apple & apple & apple & reporting \\
Identity & control-identity-justice & justice & justice & holdout \\
Random & control-random-stapler-monsoon & stapler & monsoon & reporting \\
Random & control-random-parliament-cinnamon & parliament & cinnamon & holdout \\
Random & control-random-telescope-jealousy & telescope & jealousy & reporting \\
Random & control-random-shoelace-democracy & shoelace & democracy & reporting \\
Random & control-random-flamingo-calculus & flamingo & calculus & reporting \\
Random & control-random-umbrella-photosynthesis & umbrella & photosynthesis & reporting \\
Random & control-random-origami-gravity & origami & gravity & reporting \\
Nonsense & control-nonsense-xkplm-qrvzt & xkplm & qrvzt & reporting \\
Nonsense & control-nonsense-bwfnj-tlmgk & bwfnj & tlmgk & holdout \\
\end{longtable}

\begin{center}\rule{0.5\linewidth}{0.5pt}\end{center}

\subsection{A.2 Phase 3B Triples (Transitive Path Structure)}\label{a.2-phase-3b-triples-transitive-path-structure}

8 triples testing transitivity, triangle inequality, and bridge frequency. Each triple (A, B, C) requires 6 directional legs; some legs were reused from Phase 1/2 data.

\begin{longtable}[]{@{}
  >{\raggedright\arraybackslash}p{(\columnwidth - 10\tabcolsep) * \real{0.1176}}
  >{\raggedright\arraybackslash}p{(\columnwidth - 10\tabcolsep) * \real{0.0882}}
  >{\raggedright\arraybackslash}p{(\columnwidth - 10\tabcolsep) * \real{0.0882}}
  >{\raggedright\arraybackslash}p{(\columnwidth - 10\tabcolsep) * \real{0.0882}}
  >{\raggedright\arraybackslash}p{(\columnwidth - 10\tabcolsep) * \real{0.1765}}
  >{\raggedright\arraybackslash}p{(\columnwidth - 10\tabcolsep) * \real{0.4412}}@{}}
\toprule\noalign{}
\begin{minipage}[b]{\linewidth}\raggedright
ID
\end{minipage} & \begin{minipage}[b]{\linewidth}\raggedright
A
\end{minipage} & \begin{minipage}[b]{\linewidth}\raggedright
B
\end{minipage} & \begin{minipage}[b]{\linewidth}\raggedright
C
\end{minipage} & \begin{minipage}[b]{\linewidth}\raggedright
Type
\end{minipage} & \begin{minipage}[b]{\linewidth}\raggedright
Target Bridge
\end{minipage} \\
\midrule\noalign{}
\endhead
\bottomrule\noalign{}
\endlastfoot
triple-animal-dog-poodle & animal & dog & poodle & hierarchical & dog \\
triple-emotion-nostalgia-melancholy & emotion & nostalgia & melancholy & hierarchical & nostalgia \\
triple-music-harmony-mathematics & music & harmony & mathematics & semantic-chain & harmony \\
triple-hot-energy-cold & hot & energy & cold & semantic-chain & energy \\
triple-beyonce-justice-erosion & Beyonce & justice & erosion & existing-pair & justice \\
triple-bank-river-ocean & bank & river & ocean & polysemy-extend & river \\
triple-music-stapler-mathematics & music & stapler & mathematics & random-control & stapler \\
triple-hot-flamingo-cold & hot & flamingo & cold & random-control & flamingo \\
\end{longtable}

\begin{center}\rule{0.5\linewidth}{0.5pt}\end{center}

\subsection{A.3 Phase 4 Triples (Cross-Model Bridge Topology)}\label{a.3-phase-4-triples-cross-model-bridge-topology}

8 targeted triples testing specific bridge predictions, with focus on Gemini fragmentation.

\begin{longtable}[]{@{}
  >{\raggedright\arraybackslash}p{(\columnwidth - 10\tabcolsep) * \real{0.1081}}
  >{\raggedright\arraybackslash}p{(\columnwidth - 10\tabcolsep) * \real{0.0811}}
  >{\raggedright\arraybackslash}p{(\columnwidth - 10\tabcolsep) * \real{0.0811}}
  >{\raggedright\arraybackslash}p{(\columnwidth - 10\tabcolsep) * \real{0.0811}}
  >{\raggedright\arraybackslash}p{(\columnwidth - 10\tabcolsep) * \real{0.4324}}
  >{\raggedright\arraybackslash}p{(\columnwidth - 10\tabcolsep) * \real{0.2162}}@{}}
\toprule\noalign{}
\begin{minipage}[b]{\linewidth}\raggedright
ID
\end{minipage} & \begin{minipage}[b]{\linewidth}\raggedright
A
\end{minipage} & \begin{minipage}[b]{\linewidth}\raggedright
B
\end{minipage} & \begin{minipage}[b]{\linewidth}\raggedright
C
\end{minipage} & \begin{minipage}[b]{\linewidth}\raggedright
Diagnostic Type
\end{minipage} & \begin{minipage}[b]{\linewidth}\raggedright
Bridge
\end{minipage} \\
\midrule\noalign{}
\endhead
\bottomrule\noalign{}
\endlastfoot
p4-bank-river-ocean & bank & river & ocean & polysemy-retest & river \\
p4-bank-deposit-savings & bank & deposit & savings & polysemy-financial & deposit \\
p4-light-spectrum-color & light & spectrum & color & cross-domain-concrete & spectrum \\
p4-emotion-nostalgia-melancholy & emotion & nostalgia & melancholy & abstract-retest & nostalgia \\
p4-language-metaphor-thought & language & metaphor & thought & abstract-bridge & metaphor \\
p4-tree-forest-ecosystem & tree & forest & ecosystem & concrete-hierarchical & forest \\
p4-light-chandelier-color & light & chandelier & color & random-control & chandelier \\
p4-emotion-calendar-melancholy & emotion & calendar & melancholy & random-control & calendar \\
\end{longtable}

\begin{center}\rule{0.5\linewidth}{0.5pt}\end{center}

\subsection{A.4 Phase 5 Stimuli (Cue-Strength, Dimensionality, Convergence)}\label{a.4-phase-5-stimuli-cue-strength-dimensionality-convergence}

\subsubsection{Part A: Cue-Strength Triples}\label{part-a-cue-strength-triples}

14 triples across 4 concept families at 3--4 cue levels, plus 2 random controls.

\begin{longtable}[]{@{}lllll@{}}
\toprule\noalign{}
Family & A & C & Cue Level & B (Bridge) \\
\midrule\noalign{}
\endhead
\bottomrule\noalign{}
\endlastfoot
Physical causation & sun & desert & very-high & heat \\
Physical causation & sun & desert & medium & radiation \\
Physical causation & sun & desert & low & solstice \\
Biological growth & seed & garden & very-high & plant \\
Biological growth & seed & garden & medium & germination \\
Biological growth & seed & garden & low & husk \\
Compositional hierarchy & word & paragraph & very-high & sentence \\
Compositional hierarchy & word & paragraph & medium & clause \\
Compositional hierarchy & word & paragraph & low & syllable \\
Abstract affect & emotion & melancholy & high & sadness \\
Abstract affect & emotion & melancholy & medium & nostalgia \\
Abstract affect & emotion & melancholy & low & apathy \\
Control & sun & desert & --- & umbrella \\
Control & seed & garden & --- & telescope \\
\end{longtable}

\subsubsection{Part B: Dimensionality Triples}\label{part-b-dimensionality-triples}

14 triples probing polysemous focal concepts across same-axis vs.~cross-axis configurations.

\begin{longtable}[]{@{}lllll@{}}
\toprule\noalign{}
Focal Concept & A & C & Axis Pattern & Bridge \\
\midrule\noalign{}
\endhead
\bottomrule\noalign{}
\endlastfoot
light & photon & color & same-axis (physics) & light \\
light & feather & heavy & same-axis (weight) & light \\
light & candle & darkness & same-axis (illumination) & light \\
light & photon & heavy & cross-axis & light \\
light & candle & color & partial-overlap & light \\
light & prayer & darkness & cross-axis & light \\
bank & loan & savings & same-axis (financial) & bank \\
bank & river & shore & same-axis (geographic) & bank \\
bank & loan & shore & cross-axis & bank \\
fire & spark & ash & same-axis (combustion) & fire \\
fire & ambition & motivation & same-axis (passion) & fire \\
fire & spark & motivation & cross-axis & fire \\
control & photon & color & --- & bicycle \\
control & loan & savings & --- & penguin \\
\end{longtable}

\subsubsection{Part C: Convergence Pairs}\label{part-c-convergence-pairs}

8 pairs for W-shape positional convergence testing.

\begin{longtable}[]{@{}llll@{}}
\toprule\noalign{}
From & To & Type & Expected Bridge \\
\midrule\noalign{}
\endhead
\bottomrule\noalign{}
\endlastfoot
light & color & bridge-present & spectrum \\
bank & savings & bridge-present & deposit \\
animal & poodle & bridge-present & dog \\
tree & ecosystem & bridge-present & forest \\
language & thought & bridge-absent & --- \\
hot & cold & bridge-absent & --- \\
music & mathematics & bridge-variable & harmony \\
telescope & jealousy & no-bridge-control & --- \\
\end{longtable}

\begin{center}\rule{0.5\linewidth}{0.5pt}\end{center}

\subsection{A.5 Phase 6 Pairs (Salience, Forced-Crossing, Position)}\label{a.5-phase-6-pairs-salience-forced-crossing-position}

\subsubsection{Part A: Salience Mapping (8 pairs, 30--40 reps each)}\label{part-a-salience-mapping-8-pairs-3040-reps-each}

\begin{longtable}[]{@{}llll@{}}
\toprule\noalign{}
From & To & Known Bridge & Category \\
\midrule\noalign{}
\endhead
\bottomrule\noalign{}
\endlastfoot
music & mathematics & harmony & bridge-present \\
sun & desert & heat & bridge-present \\
seed & garden & germination & bridge-present \\
light & color & spectrum & bridge-present \\
bank & ocean & river & bridge-present \\
emotion & melancholy & sadness & bridge-present \\
language & thought & --- & bridge-absent \\
hot & cold & --- & bridge-absent \\
\end{longtable}

\subsubsection{Part B: Forced-Crossing Asymmetry (8 pairs)}\label{part-b-forced-crossing-asymmetry-8-pairs}

\begin{longtable}[]{@{}llll@{}}
\toprule\noalign{}
From & To & Bridge & Type \\
\midrule\noalign{}
\endhead
\bottomrule\noalign{}
\endlastfoot
loan & shore & bank & forced-crossing \\
deposit & river & bank & forced-crossing \\
savings & cliff & bank & forced-crossing \\
photon & heavy & light & forced-crossing \\
loan & savings & --- & same-axis comparison \\
river & shore & --- & same-axis comparison \\
sun & desert & --- & same-axis comparison \\
seed & garden & --- & same-axis comparison \\
\end{longtable}

\subsubsection{Part C: Positional Bridge Scanning (10 pairs)}\label{part-c-positional-bridge-scanning-10-pairs}

\begin{longtable}[]{@{}llll@{}}
\toprule\noalign{}
From & To & Bridge & Expected Position \\
\midrule\noalign{}
\endhead
\bottomrule\noalign{}
\endlastfoot
light & color & spectrum & middle (3--5) \\
bank & savings & deposit & middle (3--5) \\
animal & poodle & dog & middle (3--5) \\
music & mathematics & harmony & middle (3--5) \\
tree & ecosystem & forest & middle (3--5) \\
loan & shore & bank & middle (3--5) \\
deposit & river & bank & middle (3--5) \\
sun & desert & heat & early (2--3) \\
seed & garden & germination & middle (3--5) \\
emotion & melancholy & sadness & late (4--6) \\
\end{longtable}

\begin{center}\rule{0.5\linewidth}{0.5pt}\end{center}

\subsection{A.6 Phase 7 Pairs and Triples (Anchoring, Curvature, Too-Central)}\label{a.6-phase-7-pairs-and-triples-anchoring-curvature-too-central}

\subsubsection{Part A: Anchoring-Effect Pairs (8 pairs, 3 pre-fill conditions each)}\label{part-a-anchoring-effect-pairs-8-pairs-3-pre-fill-conditions-each}

\begin{longtable}[]{@{}
  >{\raggedright\arraybackslash}p{(\columnwidth - 12\tabcolsep) * \real{0.1053}}
  >{\raggedright\arraybackslash}p{(\columnwidth - 12\tabcolsep) * \real{0.0702}}
  >{\raggedright\arraybackslash}p{(\columnwidth - 12\tabcolsep) * \real{0.1404}}
  >{\raggedright\arraybackslash}p{(\columnwidth - 12\tabcolsep) * \real{0.2281}}
  >{\raggedright\arraybackslash}p{(\columnwidth - 12\tabcolsep) * \real{0.1930}}
  >{\raggedright\arraybackslash}p{(\columnwidth - 12\tabcolsep) * \real{0.1579}}
  >{\raggedright\arraybackslash}p{(\columnwidth - 12\tabcolsep) * \real{0.1053}}@{}}
\toprule\noalign{}
\begin{minipage}[b]{\linewidth}\raggedright
From
\end{minipage} & \begin{minipage}[b]{\linewidth}\raggedright
To
\end{minipage} & \begin{minipage}[b]{\linewidth}\raggedright
Bridge
\end{minipage} & \begin{minipage}[b]{\linewidth}\raggedright
Incongruent
\end{minipage} & \begin{minipage}[b]{\linewidth}\raggedright
Congruent
\end{minipage} & \begin{minipage}[b]{\linewidth}\raggedright
Neutral
\end{minipage} & \begin{minipage}[b]{\linewidth}\raggedright
Role
\end{minipage} \\
\midrule\noalign{}
\endhead
\bottomrule\noalign{}
\endlastfoot
music & mathematics & harmony & symmetry & rhythm & element & heading-bridge \\
sun & desert & heat & drought & warmth & element & heading-bridge \\
seed & garden & germination & sprout & growth & element & heading-bridge \\
light & color & spectrum & wavelength & prism & element & heading-bridge \\
bank & ocean & river & estuary & stream & element & heading-bridge \\
emotion & melancholy & sadness & mood & grief & element & heading-bridge \\
loan & shore & bank & interest & credit & element & forced-crossing \\
animal & poodle & dog & mammal & canine & element & taxonomic \\
\end{longtable}

\subsubsection{Part B: Curvature Triangles (8 triangles)}\label{part-b-curvature-triangles-8-triangles}

\begin{longtable}[]{@{}
  >{\raggedright\arraybackslash}p{(\columnwidth - 8\tabcolsep) * \real{0.0882}}
  >{\raggedright\arraybackslash}p{(\columnwidth - 8\tabcolsep) * \real{0.0882}}
  >{\raggedright\arraybackslash}p{(\columnwidth - 8\tabcolsep) * \real{0.0882}}
  >{\raggedright\arraybackslash}p{(\columnwidth - 8\tabcolsep) * \real{0.3529}}
  >{\raggedright\arraybackslash}p{(\columnwidth - 8\tabcolsep) * \real{0.3824}}@{}}
\toprule\noalign{}
\begin{minipage}[b]{\linewidth}\raggedright
A
\end{minipage} & \begin{minipage}[b]{\linewidth}\raggedright
B
\end{minipage} & \begin{minipage}[b]{\linewidth}\raggedright
C
\end{minipage} & \begin{minipage}[b]{\linewidth}\raggedright
Vertex Type
\end{minipage} & \begin{minipage}[b]{\linewidth}\raggedright
Relationship
\end{minipage} \\
\midrule\noalign{}
\endhead
\bottomrule\noalign{}
\endlastfoot
loan & bank & river & polysemous (homonym) & financial/geographic \\
deposit & bank & shore & polysemous (homonym) & financial/geographic \\
photon & light & heavy & polysemous (homonym) & electromagnetic/weight \\
candle & light & feather & polysemous (homonym) & illumination/weight \\
sun & heat & desert & non-polysemous & causal-chain \\
seed & germination & garden & non-polysemous & process-chain \\
music & harmony & mathematics & non-polysemous & cross-domain \\
word & sentence & paragraph & non-polysemous & compositional \\
\end{longtable}

\subsubsection{Part C: Too-Central Pairs (10 pairs)}\label{part-c-too-central-pairs-10-pairs}

\begin{longtable}[]{@{}llll@{}}
\toprule\noalign{}
From & To & Candidate Bridge & Category \\
\midrule\noalign{}
\endhead
\bottomrule\noalign{}
\endlastfoot
spark & ash & fire & too-central \\
acorn & timber & tree & too-central \\
flour & bread & dough & too-central \\
hot & cold & warm & obvious-useful \\
infant & elderly & adolescent & obvious-useful \\
whisper & shout & speak & obvious-useful \\
rain & ocean & water & boundary \\
egg & chicken & embryo & boundary \\
ice & steam & water & boundary \\
dawn & dusk & noon & boundary \\
\end{longtable}

\begin{center}\rule{0.5\linewidth}{0.5pt}\end{center}

\subsection{A.7 Phase 8 Pairs (Fragility, Gradient, Distance)}\label{a.7-phase-8-pairs-fragility-gradient-distance}

\subsubsection{Part A: Bridge Fragility (8 new pairs)}\label{part-a-bridge-fragility-8-new-pairs}

\begin{longtable}[]{@{}lllll@{}}
\toprule\noalign{}
From & To & Predicted Bridge & Competitor Level & Pre-fill \\
\midrule\noalign{}
\endhead
\bottomrule\noalign{}
\endlastfoot
question & answer & reasoning & low (0--1) & inquiry \\
parent & child & birth & low (0--1) & family \\
cause & effect & mechanism & low (1--2) & reason \\
winter & summer & spring & medium (3--4) & snow \\
student & professor & research & medium (3--5) & university \\
science & art & creativity & high (6--10) & experiment \\
ocean & mountain & landscape & high (6--10) & wave \\
brain & computer & algorithm & high (8--12) & neuron \\
\end{longtable}

\subsubsection{Part B: Gradient vs.~Causal-Chain Pairs (20 pairs)}\label{part-b-gradient-vs.-causal-chain-pairs-20-pairs}

\textbf{Gradient-midpoint (10):}

\begin{longtable}[]{@{}llll@{}}
\toprule\noalign{}
From & To & Bridge & Dimension \\
\midrule\noalign{}
\endhead
\bottomrule\noalign{}
\endlastfoot
whisper & shout & speak & vocal intensity \\
hot & cold & warm & temperature \\
dawn & dusk & noon & time of day \\
infant & elderly & adolescent & life stage \\
crawl & sprint & walk & locomotion speed \\
pond & ocean & lake & water body scale \\
pebble & boulder & rock & stone size \\
drizzle & downpour & rain & precipitation \\
village & metropolis & city & settlement scale \\
murmur & scream & speech & vocal intensity \\
\end{longtable}

\textbf{Causal-chain (10):}

\begin{longtable}[]{@{}llll@{}}
\toprule\noalign{}
From & To & Bridge & Process \\
\midrule\noalign{}
\endhead
\bottomrule\noalign{}
\endlastfoot
grape & wine & fermentation & winemaking \\
ore & jewelry & metal & smelting-crafting \\
caterpillar & butterfly & cocoon & metamorphosis \\
clay & vase & pottery & crafting \\
wool & sweater & yarn & textile \\
milk & cheese & curd & dairy \\
seed & fruit & flower & reproduction \\
sand & glass & heat & manufacturing \\
acorn & timber & tree & growth-harvest \\
flour & bread & dough & baking \\
\end{longtable}

\subsubsection{Part C: Gait-Normalized Distance (16 pairs)}\label{part-c-gait-normalized-distance-16-pairs}

8 reference pairs (short to very-long expected distance) and 8 test pairs spanning the distance range. See Section 7.4 for the finding that cross-model distance normalization produces zero improvement.

\begin{center}\rule{0.5\linewidth}{0.5pt}\end{center}

\subsection{A.8 Phase 9--11 Pairs}\label{a.8-phase-911-pairs}

Phase 9 tested 6 dominance pairs (selected from Phase 8B data), 10 transformation/gradient pairs (retesting Gemini deficit), and 8 facilitation pairs (testing the pre-fill facilitation crossover).

Phase 10A used a 12-pair battery (8 forward + 4 reverse from the core pairs above) across 5 new models. Phase 10B tested 8 pairs under 3 pre-fill relation classes (on-axis, same-domain, unrelated).

Phase 11A reused the same 12-pair battery across 4 additional new models. Phase 11B screened 4 control candidates (accordion--stalactite, turmeric--trigonometry, barnacle--sonnet, magnesium--ballet) across 6 models. Phase 11C tested 6 pairs across a $2\times 2$ waypoint/temperature grid (3 forward pairs: light--color, hot--cold, emotion--melancholy; plus their reverses) on 3 models (Claude, GPT, DeepSeek).

\begin{center}\rule{0.5\linewidth}{0.5pt}\end{center}

\subsection{A.9 Model Inventory}\label{a.9-model-inventory}

\begin{longtable}[]{@{}
  >{\raggedright\arraybackslash}p{(\columnwidth - 8\tabcolsep) * \real{0.0816}}
  >{\raggedright\arraybackslash}p{(\columnwidth - 8\tabcolsep) * \real{0.2653}}
  >{\raggedright\arraybackslash}p{(\columnwidth - 8\tabcolsep) * \real{0.2041}}
  >{\raggedright\arraybackslash}p{(\columnwidth - 8\tabcolsep) * \real{0.2857}}
  >{\raggedright\arraybackslash}p{(\columnwidth - 8\tabcolsep) * \real{0.1633}}@{}}
\toprule\noalign{}
\begin{minipage}[b]{\linewidth}\raggedright
ID
\end{minipage} & \begin{minipage}[b]{\linewidth}\raggedright
Display Name
\end{minipage} & \begin{minipage}[b]{\linewidth}\raggedright
Provider
\end{minipage} & \begin{minipage}[b]{\linewidth}\raggedright
OpenRouter ID
\end{minipage} & \begin{minipage}[b]{\linewidth}\raggedright
Cohort
\end{minipage} \\
\midrule\noalign{}
\endhead
\bottomrule\noalign{}
\endlastfoot
claude & Claude Sonnet 4.6 & Anthropic & anthropic/claude-sonnet-4.6 & Original \\
gpt & GPT-5.2 & OpenAI & openai/gpt-5.2 & Original \\
grok & Grok 4.1 Fast & xAI & x-ai/grok-4.1-fast & Original \\
gemini & Gemini 3 Flash & Google & google/gemini-3-flash-preview & Original \\
minimax & MiniMax M2.5 & MiniMax & minimax/minimax-m2.5 & Phase 10A \\
kimi & Kimi K2.5 & Moonshot AI & moonshotai/kimi-k2.5 & Phase 10A \\
glm & GLM 5 & Zhipu AI & z-ai/glm-5 & Phase 10A (rate-limited) \\
qwen & Qwen 3.5 397B-A17B & Alibaba Cloud & qwen/qwen3.5-397b-a17b & Phase 10A \\
llama & Llama 3.1 8B Instruct & Meta & meta-llama/llama-3.1-8b-instruct & Phase 10A \\
deepseek & DeepSeek V3.2 & DeepSeek & deepseek/deepseek-v3.2 & Phase 11A \\
mistral & Mistral Large 3 & Mistral AI & mistralai/mistral-large-2512 & Phase 11A \\
cohere & Cohere Command A & Cohere & cohere/command-a & Phase 11A \\
llama4 & Llama 4 Maverick & Meta & meta-llama/llama-4-maverick & Phase 11A \\
\end{longtable}

All models accessed via OpenRouter API. GLM 5 was blocked by upstream rate limiting and excluded from analysis. Llama 3.1 8B Instruct is the sole small-scale model, included to probe scale effects.
        % Appendix A: Concept Pair Inventory
\section{Phase-by-Phase Results Summary}\label{appendix-b-phase-by-phase-results-summary}

This appendix provides a condensed summary of each phase's design, findings, predictions, and prediction accuracy. It is intended for readers who want to trace the exploration-first methodology in detail.

\subsection{Phase 1: Pilot and Prompt Selection}\label{phase-1-pilot-and-prompt-selection}

\textbf{Runs:} \textasciitilde2,480 (4 models, 36 pairs, 20 runs each)

\textbf{Key question:} Do models produce consistent waypoint paths?

\textbf{Design:} Evaluated two prompt formats (semantic, direct) on 15 holdout pairs. Collected 20-run batches on 21 reporting pairs across 4 models (Claude Sonnet 4.6, GPT-5.2, Grok 4.1 Fast, Gemini 3 Flash). Tested both 5-waypoint and 10-waypoint conditions.

\textbf{Primary findings:} - Models have distinct conceptual gaits: Claude 0.578 vs GPT 0.258 (2.2x gap) {[}R1{]}. - Control validation: nonsense pairs produce near-zero Jaccard (0.062). - Polysemy sense differentiation: cross-pair Jaccard 0.011--0.062 {[}R3{]}. - 5-waypoint paths are genuine coarse-grainings of 10-waypoint paths (70.5\% shared, 67.9\% subsequence rate) {[}O10{]}. - Semantic prompt format selected over direct format.

\textbf{Prediction accuracy:} N/A (exploratory phase).

\begin{center}\rule{0.5\linewidth}{0.5pt}\end{center}

\subsection{Phase 2: Reversals and Path Consistency}\label{phase-2-reversals-and-path-consistency}

\textbf{Runs:} \textasciitilde960 (840 reverse + 120 polysemy supplementary)

\textbf{Key question:} Is A-to-B navigation symmetric with B-to-A?

\textbf{Design:} Reversed all 21 reporting pairs across 4 models (20 runs each). Added 120 supplementary runs to correct the Phase 1 polysemy artifact (G1).

\textbf{Primary findings:} - Navigation is fundamentally asymmetric: mean directional asymmetry 0.811 {[}R2{]}. - 87\% of 84 pair/model combinations show significant asymmetry (permutation test, p \textless{} 0.05). - Dual-anchor effect: U-shaped convergence profile {[}O2{]}. - Category symmetry predictions failed for 4 of 8 categories {[}G3{]}.

\textbf{Prediction accuracy:} N/A (exploratory phase).

\begin{center}\rule{0.5\linewidth}{0.5pt}\end{center}

\subsection{Phase 3: Positional Convergence and Transitive Path Structure}\label{phase-3-positional-convergence-and-transitive-path-structure}

\textbf{Runs:} \textasciitilde1,260 (\textasciitilde660 Phase 3B new runs; Phase 3A was pure analysis)

\textbf{Key question:} Do paths compose? Does the triangle inequality hold?

\textbf{Design:} Part A: pure analysis of Phase 2 reversal data for positional convergence. Part B: 8 concept triples (hierarchical, semantic-chain, polysemy-extend, random-control) with 10 runs per leg, testing transitivity and triangle inequality.

\textbf{Primary findings:} - Hierarchical transitivity 0.175 vs random 0.036 (4.9x, non-overlapping CIs) {[}R4{]}. - Triangle inequality holds in 90.6\% of cases (29/32 triple/model combinations). - Bridge frequency model-dependent: ``dog'' 15--100\%, ``harmony'' 0--100\%. - ``Energy'' fails as hot--cold bridge {[}G4{]}.

\textbf{Prediction accuracy:} N/A (exploratory phase).

\begin{center}\rule{0.5\linewidth}{0.5pt}\end{center}

\subsection{Phase 4: Cross-Model Bridge Topology}\label{phase-4-cross-model-bridge-topology}

\textbf{Runs:} \textasciitilde1,600 (Part A: 0 new runs, pure analysis; Part B: \textasciitilde1,520 new runs)

\textbf{Key question:} Do models agree on bridge concepts?

\textbf{Design:} Part A: analyzed inter-model bridge agreement across 6 non-control triples from Phase 3B. Part B: 8 targeted triples with 20 runs per leg across 4 models, testing specific bridge predictions.

\textbf{Primary findings:} - Three models agree, Gemini diverges (isolation index 0.136) {[}O1{]}. - Bridges are structural bottlenecks: ``spectrum'' 1.00 all models {[}R6{]}. - ``Metaphor'' fails at 0.00 all models {[}G5{]}. - Universal concrete bridging (spectrum 100\%), universal abstract failure (metaphor 0\%). - Triangle inequality replicates at 93.8\%.

\textbf{Prediction accuracy:} 81.3\% (26/32). The highest in the benchmark --- characterization questions.

\begin{center}\rule{0.5\linewidth}{0.5pt}\end{center}

\subsection{Phase 5: Cue-Strength and Dimensionality}\label{phase-5-cue-strength-and-dimensionality}

\textbf{Runs:} \textasciitilde3,720 new + 200 reused

\textbf{Key question:} Is bridge frequency graded by semantic cue strength?

\textbf{Design:} Part A: 4 concept families x 3 cue levels x 4 models. Part B: polysemy dimensionality with same-axis vs cross-axis triples. Part C: triple-anchor convergence (W-shape test).

\textbf{Primary findings:} - Cue-strength gradient real: 12/16 monotonic {[}R7{]}. All 4 failures in biological-growth family. - ``Germination'' \textgreater{} ``plant'': process \textgreater{} object {[}O4, G8{]}. - Forced crossing discovery (loan--bank--shore) {[}O5{]}. - ``Fire'' dead as bridge (too-central) {[}O6, G10{]}. - W-shape aggregate null {[}G11{]}. Gemini threshold hypothesis fails {[}G7{]}.

\textbf{Prediction accuracy:} 42.9\% (Phase 5A: 64.6\%; Phase 5B: 31/48 but overall combined lower).

\begin{center}\rule{0.5\linewidth}{0.5pt}\end{center}

\subsection{Phase 6: Navigational Salience and Bridge Position}\label{phase-6-navigational-salience-and-bridge-position}

\textbf{Runs:} \textasciitilde2,080 new + 280 reused

\textbf{Key question:} Where do bridges appear in paths, and how does salience vary?

\textbf{Design:} Part A: salience mapping (40 runs/condition, 8 pairs, 4 models). Part B: forced-crossing asymmetry test (8 pairs). Part C: positional bridge scanning (10 pairs, 4 models).

\textbf{Primary findings:} - Salience distributions non-uniform: 31/32 conditions reject uniformity {[}O11{]}. - Forced-crossing asymmetry matches baseline (0.817 vs 0.811) {[}G12{]}. - Bridges anchor early (positions 1--2), not midpoint {[}O12{]}. - Peak-detection contrast 0.345 (CI excludes zero) {[}O13{]}. - Forced-crossing bridges positionally unstable (SD 1.71 vs 0.52) {[}G14{]}. - GPT highest entropy (3.44); Claude lowest (2.59).

\textbf{Prediction accuracy:} \textasciitilde50\% (structural predictions mixed).

\begin{center}\rule{0.5\linewidth}{0.5pt}\end{center}

\subsection{Phase 7: Early Anchoring and Navigational Mechanics}\label{phase-7-early-anchoring-and-navigational-mechanics}

\textbf{Runs:} \textasciitilde2,360 new + 920 reused

\textbf{Key question:} Does pre-filling a waypoint causally affect the path?

\textbf{Design:} Part A: 8 pairs x 3 pre-fill conditions (incongruent, congruent, neutral) x 4 models. Part B: 8 curvature triangles. Part C: 10 too-central pairs.

\textbf{Primary findings:} - Pre-filling causally displaces bridges: mean 0.515, CI excludes zero {[}O15{]}. - Mechanism is associative primacy, not directional heading {[}G15{]}. - Bridge survival is bimodal: sadness/dog survive; harmony/germination collapse. - Taxonomic bridges resist displacement (0.140). - Triangle inequality replicates at 90.6\%. - Cross-model distance fails (r = 0.170) {[}O18, G17{]}. - Gradient \textgreater{} causal-chain bridges (0.730 vs 0.496).

\textbf{Prediction accuracy:} \textasciitilde50\% (structural replications succeed, mechanistic novel predictions fail).

\begin{center}\rule{0.5\linewidth}{0.5pt}\end{center}

\subsection{Phase 8: Bridge Fragility and Gemini Deficit}\label{phase-8-bridge-fragility-and-gemini-deficit}

\textbf{Runs:} \textasciitilde2,690 new + 2,960 reused

\textbf{Key question:} Can bridge fragility be predicted from single variables?

\textbf{Design:} Part A: 14 fragility pairs testing route exclusivity (competitor count). Part B: 20 gradient/causal pairs testing Gemini gradient blindness. Part C: 16 distance pairs testing gait normalization.

\textbf{Primary findings:} - Route exclusivity fails: \(\rho\) = 0.116, opposite direction {[}G20{]}. - Gemini gradient blindness fails: interaction 0.046, null {[}G21{]}. - Gait normalization: zero improvement (r = 0.212 before and after) {[}G22{]}. - O17 replicates: gradient 0.770 vs causal 0.578. - New hypotheses proposed: H9 (dominance ratio), H10 (transformation blindness), H11 (facilitation).

\textbf{Prediction accuracy:} 24\% (6/25). The benchmark's lowest --- all novel mechanistic predictions fail.

\begin{center}\rule{0.5\linewidth}{0.5pt}\end{center}

\subsection{Phase 9: Dominance, Transformation, and Facilitation}\label{phase-9-dominance-transformation-and-facilitation}

\textbf{Runs:} \textasciitilde3,037 new + \textasciitilde5,270 reused

\textbf{Key question:} Can bridge survival be explained by dominance ratio, transformation type, or facilitation crossover?

\textbf{Design:} Part A: 14 pairs testing dominance ratio as fragility predictor. Part B: 20 transformation/gradient pairs retesting Gemini. Part C: 14 facilitation pairs testing the crossover regression.

\textbf{Primary findings:} - Dominance ratio fails: \(\rho\) = 0.157 {[}G23{]}. - Transformation blindness fails with reversal: Gemini transformation mean 0.667 vs gradient 0.293 {[}G24{]}. - Crossover regression: slope correct direction but CI includes zero {[}G25{]}. - Pre-fill content is secondary modulator; 5/8 Phase 7A pairs replicate. - Marginal bridges show massive facilitation: mean 3.761x {[}O22{]}.

\textbf{Prediction accuracy:} 20\% (5/25). Mechanism ceiling confirmed.

\begin{center}\rule{0.5\linewidth}{0.5pt}\end{center}

\subsection{Phase 10: Model Generality and Relation Classes}\label{phase-10-model-generality-and-relation-classes}

\textbf{Runs:} \textasciitilde1,680 new + 778 reused

\textbf{Key question:} Do structural findings generalize to new models?

\textbf{Design:} Part A: 5 new models (Qwen, MiniMax, Kimi, GLM, Llama 8B) on 12 pairs with reliability probing. Part B: 4 core models on 8 pairs with 3 pre-fill relation classes.

\textbf{Primary findings:} - R1 replicates: new model gaits 0.298--0.508 {[}R1 extended{]}. - R2 replicates: all new models exceed 0.60 asymmetry threshold {[}R2 extended{]}. - Bridge frequency CI includes zero (0.721 vs 0.817): bridge structure generalizes. - Llama 8B sole outlier (bridge freq 0.200): scale effect {[}O25{]}. - Relation class matters (Friedman p = 0.034): unrelated \textless{} related {[}O26{]}. - Predicted relation class ordering wrong {[}G27{]}. G26 resurrected.

\textbf{Prediction accuracy:} 50\% (9/18). Recovery from mechanism floor --- more replication/structural predictions.

\begin{center}\rule{0.5\linewidth}{0.5pt}\end{center}

\subsection{Phase 11: Expanded Generality, Control Revision, and Robustness}\label{phase-11-expanded-generality-control-revision-and-robustness}

\textbf{Runs:} \textasciitilde2,040 new + \textasciitilde340 reused

\textbf{Key question:} Do findings survive more models and protocol variation?

\textbf{Design:} Part A: 4 new models (DeepSeek, Mistral, Cohere, Llama 4 Maverick) on 12 pairs. Part B: 4 control pair candidates across 6 models. Part C: 3 models x 2x2 waypoint/temperature grid.

\textbf{Primary findings:} - R1/R2 universal across all 12 models. Mistral record gait (0.747) {[}O28{]}. - Llama scale confirmed: Maverick 0.724 vs 8B 0.200. - All 4 control candidates fail (0/24 cells pass). Stapler-monsoon fails for all 12 models {[}G28{]}. - ANOVA: model identity \(\eta^2\) = 0.242; protocol \(\approx\) 0 {[}O30{]}. - Bridge frequency most robust (\textgreater{} 0.97 all conditions) {[}O31{]}. - Asymmetry resolution-dependent: 5wp 0.594, 9wp 0.669 {[}O32, G29{]}.

\textbf{Prediction accuracy:} 44\% (8/18). Intermediate --- mix of replication and structural predictions.

\begin{center}\rule{0.5\linewidth}{0.5pt}\end{center}

\subsection{Cumulative Statistics}\label{cumulative-statistics}

\begin{longtable}[]{@{}llll@{}}
\toprule\noalign{}
Phase & New Runs & Prediction Accuracy & Dominant Prediction Type \\
\midrule\noalign{}
\endhead
\bottomrule\noalign{}
\endlastfoot
1--3 & \textasciitilde4,700 & N/A (exploratory) & --- \\
4 & \textasciitilde1,600 & 81.3\% & Characterization \\
5 & \textasciitilde3,720 & 42.9\% & Characterization + structure \\
6 & \textasciitilde2,080 & \textasciitilde50\% & Structure \\
7 & \textasciitilde2,360 & \textasciitilde50\% & Structure + mechanism \\
8 & \textasciitilde2,690 & 24\% & Mechanism \\
9 & \textasciitilde3,037 & 20\% & Mechanism \\
10 & \textasciitilde1,680 & 50\% & Replication + structure \\
11 & \textasciitilde2,040 & 44\% & Replication + robustness \\
\textbf{Total} & \textbf{\textasciitilde21,540} & & \\
\end{longtable}

The prediction accuracy trajectory --- 81.3\% \(\to\) 42.9\% \(\to\) \textasciitilde50\% \(\to\) \textasciitilde50\% \(\to\) 24\% \(\to\) 20\% \(\to\) 50\% \(\to\) 44\% --- tracks the shift between prediction types. Replication predictions succeed at \textasciitilde80\%; structural predictions at \textasciitilde50\%; mechanistic predictions at 0--15\%.
        % Appendix B: Phase-by-Phase Results
\section{The Graveyard --- Dead Ends, Failed Hypotheses, and Deprecated Approaches}\label{appendix-c-the-graveyard-dead-ends-failed-hypotheses-and-deprecated-approaches}

This appendix documents 29 dead ends across 11 experimental phases. These entries are ordered chronologically and represent an honest accounting of the benchmark's learning curve. One entry (G26) was subsequently resurrected with additional data.

\subsection{Summary Table}\label{summary-table}

\begin{longtable}[]{@{}
  >{\raggedright\arraybackslash}p{(\columnwidth - 10\tabcolsep) * \real{0.0816}}
  >{\raggedright\arraybackslash}p{(\columnwidth - 10\tabcolsep) * \real{0.1429}}
  >{\raggedright\arraybackslash}p{(\columnwidth - 10\tabcolsep) * \real{0.2245}}
  >{\raggedright\arraybackslash}p{(\columnwidth - 10\tabcolsep) * \real{0.2245}}
  >{\raggedright\arraybackslash}p{(\columnwidth - 10\tabcolsep) * \real{0.2041}}
  >{\raggedright\arraybackslash}p{(\columnwidth - 10\tabcolsep) * \real{0.1224}}@{}}
\toprule\noalign{}
\begin{minipage}[b]{\linewidth}\raggedright
ID
\end{minipage} & \begin{minipage}[b]{\linewidth}\raggedright
Phase
\end{minipage} & \begin{minipage}[b]{\linewidth}\raggedright
Hypothesis
\end{minipage} & \begin{minipage}[b]{\linewidth}\raggedright
Predicted
\end{minipage} & \begin{minipage}[b]{\linewidth}\raggedright
Observed
\end{minipage} & \begin{minipage}[b]{\linewidth}\raggedright
Type
\end{minipage} \\
\midrule\noalign{}
\endhead
\bottomrule\noalign{}
\endlastfoot
G1 & 1 & Perfect polysemy differentiation & Jaccard = 0.000 & Artifact (empty data) & Data error \\
G2 & 2 & Monotonic convergence & Increasing overlap & U-shape (dual-anchor) & Model too simple \\
G3 & 2 & Category symmetry (4 of 8) & High symmetry for antonyms, near-synonyms, controls & Asymmetry 0.596--0.986 & Asymmetry is universal \\
G4 & 3 & ``Energy'' as hot--cold bridge & Bridge frequency \textgreater{} 0 & Near-zero (0.036 transitivity) & Association \(\neq\) navigation \\
G5 & 4 & ``Metaphor'' as language--thought bridge & Frequency 0.50--1.00 & 0.00 all models & Intuition \(\neq\) navigation \\
G6 & 4 & Concrete/abstract predicts Gemini & tree--forest--ecosystem \textgreater{} 0.40 & Gemini 0.10 & Category too coarse \\
G7 & 5 & Gemini cue-strength threshold & Threshold highest & Threshold lowest (1.79) & Quantitative model fails \\
G8 & 5 & Plant \textgreater{} germination (intuitive cue) & Plant near 100\% & Germination wins (process \textgreater{} object) & Human intuition fails \\
G9 & 5 & Dimensional independence of polysemy & Same-axis \textgreater{} cross-axis by 0.40 & Cross-axis higher (forced crossings) & Constraint \textgreater{} association \\
G10 & 5 & Fire as reliable control bridge & Reliable same-axis bridging & Near-zero all models & Too-central failure mode \\
G11 & 5 & W-shape aggregate prediction & Bridge-present \textgreater{} bridge-absent & No aggregate signal & Aggregation kills signal \\
G12 & 6 & Forced-crossing reduces asymmetry & 0.50--0.70 asymmetry & 0.817 (= baseline) & Asymmetry is structural \\
G13 & 6 & Semantic distance predicts position & r significant & r = 0.239, p = 0.486 & Early-anchoring overrides geometry \\
G14 & 6 & Forced-crossing positional stability & SD \textless{} 0.8 & SD = 1.71 & Obligatory \(\neq\) stable \\
G15 & 7 & Directional-heading theory (strong form) & Incongruent \textgreater{} congruent displacement & Indistinguishable (0.515 vs 0.436) & Associative primacy dominates \\
G16 & 7 & Polysemous curvature & Higher triangle excess & Difference 0.053, CI includes zero & Distance metric invalid \\
G17 & 7 & Cross-model distance validity & r \textgreater{} 0.50 & r = 0.170 & Distances are model-dependent \\
G18 & 7 & Too-central as binary category & Clean separation & CI includes zero; only fire/water qualify & Category based on intuition \\
G19 & 7 & Forced-crossing resists pre-fill & Survival \textgreater{} 0.90 & Survival 0.267 & Obligatory \(\neq\) robust \\
G20 & 8 & Route exclusivity (competitor count) & \(\rho\) \textless{} --0.70 & \(\rho\) = 0.116 & Single variable fails \\
G21 & 8 & Gemini gradient blindness & Gradient zeros \(\geq\) 5/10 & Causal zeros 6/10 (reversed) & Type-based approach wrong \\
G22 & 8 & Gait normalization rescues distance & Normalized r \textgreater{} 0.50 & r = 0.212 (0.000 improvement) & Ordinal disagreement \\
G23 & 9 & Dominance ratio predicts fragility & \(\rho\) \textgreater{} 0.50 & \(\rho\) = 0.157 & Retrospective overfitting \\
G24 & 9 & Gemini transformation blindness & Transformation deficit & Transformation advantage (reversed) & Third type-based failure \\
G25 & 9 & Pre-fill crossover regression & Significant slope & CI includes zero (partial) & Too heterogeneous \\
G26 & 10 & Bridge generalization fails & CI excludes zero & \textbf{Resurrected}: CI includes zero & Infrastructure artifact \\
G27 & 10 & Relation class ordering & On-axis \textless{} unrelated \textless{} same-domain & Unrelated \textless{} on-axis \(\approx\) same-domain & Related vs.~unrelated is binary \\
G28 & 11 & Control pair revision & \(\geq\) 2/4 candidates pass & 0/24 cells pass & LLMs navigate between anything \\
G29 & 11 & Asymmetry universality across protocol & Mean \textgreater{} 0.60 all conditions & 5wp: 0.594; 9wp: 0.669 & Resolution-dependent \\
\end{longtable}

\subsection{Detailed Entries}\label{detailed-entries}

\subsubsection{G1 --- Phase 1: Polysemy Cross-Pair Jaccard of 0.000 (Artifact)}\label{g1-phase-1-polysemy-cross-pair-jaccard-of-0.000-artifact}

\textbf{Hypothesis:} Perfect sense differentiation with cross-pair Jaccard of 0.000 for all three polysemy groups.

\textbf{What happened:} Holdout pairs had no pilot data --- the zero Jaccard was comparing real data against empty sets.

\textbf{Resolution:} Phase 2 supplementary runs (120 forward runs) corrected values to 0.011--0.062. Sense differentiation is genuine and strong, but not perfect.

\textbf{Lesson:} Always verify the denominator. Zero is suspicious.

\subsubsection{G2 --- Phase 2: Monotonic Convergence Prediction}\label{g2-phase-2-monotonic-convergence-prediction}

\textbf{Hypothesis:} Mirror-matching forward position \emph{i} against reverse position \emph{(5-i)} would show monotonically increasing overlap, based on the starting-point hypothesis.

\textbf{What happened:} U-shape instead --- positions 1 and 5 both elevated (0.102, 0.129), valley in the middle (0.057--0.085). Both endpoints act as anchors.

\textbf{Resolution:} Replaced starting-point hypothesis with dual-anchor hypothesis (O2).

\textbf{Lesson:} Simple unidirectional models underestimate the influence of the target concept.

\subsubsection{G3 --- Phase 2: Category Symmetry Predictions}\label{g3-phase-2-category-symmetry-predictions}

\textbf{Hypothesis:} Several categories would show high path symmetry: antonyms, near-synonyms, random controls, nonsense controls.

\textbf{What happened:} Observed asymmetry of 0.596 (antonyms), 0.665 (near-synonyms), 0.908 (random), 0.986 (nonsense). Four of eight category symmetry predictions failed.

\textbf{Resolution:} Established that asymmetry is pervasive and structural (R2). The 0.811 mean was far higher than predicted.

\textbf{Lesson:} Asymmetry is the default. Symmetry is the exception, produced only by strong semantic bridges.

\subsubsection{G4 --- Phase 3: ``Energy'' as Hot--Cold Bridge}\label{g4-phase-3-energy-as-hotcold-bridge}

\textbf{Hypothesis:} ``Energy,'' associated with both ``hot'' and ``cold,'' would appear as a bridge concept.

\textbf{What happened:} Near-zero transitivity (0.036). Hot--cold paths traverse the temperature gradient (warm, tepid, cool); hot--energy paths traverse physics terminology.

\textbf{Resolution:} Established that the temperature axis and physics-of-energy axis are parallel non-intersecting dimensions. Association is not navigation.

\textbf{Lesson:} Semantic association does not guarantee navigational intermediacy.

\subsubsection{G5 --- Phase 4: ``Metaphor'' as Language--Thought Bridge}\label{g5-phase-4-metaphor-as-languagethought-bridge}

\textbf{Hypothesis:} Bridge frequency of 0.50--1.00 for metaphor on language--thought. Metaphor is arguably the quintessential bridge between language and thought.

\textbf{What happened:} 0.00 for all four models. Zero appearances out of 80 runs. Yet transitivity was high (Claude 0.437) through ``communication,'' ``meaning,'' ``cognition.''

\textbf{Resolution:} Sharpened the bridge definition: bridges must name the primary axis of connection, not merely be associated with both endpoints.

\textbf{Lesson:} The biggest single prediction miss in the benchmark. Intuitive semantic importance has zero predictive power for navigational bridge behavior.

\subsubsection{G6 --- Phase 4: Concrete/Abstract as Gemini's Fragmentation Boundary}\label{g6-phase-4-concreteabstract-as-geminis-fragmentation-boundary}

\textbf{Hypothesis:} The concrete/abstract axis predicts Gemini's bridge failures. Predicted Gemini success on tree--forest--ecosystem (concrete, hierarchical).

\textbf{What happened:} Gemini observed at 0.10. Bank--deposit--savings (abstract/financial) succeeds at 1.00.

\textbf{Resolution:} Replaced with frame-crossing hypothesis. Forest-to-ecosystem requires a scale transition (botanical to ecological), not an abstractness gradient.

\textbf{Lesson:} Concrete/abstract is too coarse. Frame membership is the relevant variable.

\subsubsection{G7 --- Phase 5: Gemini Cue-Strength Threshold}\label{g7-phase-5-gemini-cue-strength-threshold}

\textbf{Hypothesis:} Gemini has a measurably higher cue-strength threshold than other models.

\textbf{What happened:} Gemini's logistic threshold (1.79) is the lowest, not highest. CI on difference includes zero.

\textbf{Resolution:} The quantitative threshold model is falsified. The qualitative frame-crossing model survives.

\textbf{Lesson:} A correct qualitative model can generate incorrect quantitative predictions.

\subsubsection{G8 --- Phase 5: Plant \textgreater{} Germination (Intuitive Cue Strength)}\label{g8-phase-5-plant-germination-intuitive-cue-strength}

\textbf{Hypothesis:} ``Plant'' (rated very-high cue strength) would bridge seed--garden at near-100\%.

\textbf{What happened:} Complete inversion. Germination outperforms plant in all 4 models. Claude: 1.00 vs 0.00.

\textbf{Resolution:} Discovered the distinction between associative strength and navigational salience. Process-naming (``germination'') outperforms object-naming (``plant'').

\textbf{Lesson:} Human intuition about ``most related'' does not predict navigational routing. Calibrate empirically.

\subsubsection{G9 --- Phase 5: Dimensional Independence of Polysemy}\label{g9-phase-5-dimensional-independence-of-polysemy}

\textbf{Hypothesis:} Same-axis bridge frequency would exceed cross-axis by at least 0.40.

\textbf{What happened:} Cross-axis mean (0.291) \emph{higher} than same-axis (0.220). Driven by loan--bank--shore at 0.95--1.00.

\textbf{Resolution:} Polysemy creates forced crossing points, not dimensional walls. Cross-axis paths have no alternative to the polysemous bridge.

\textbf{Lesson:} The most reliable bridges are bottlenecks, not associations. Constraint produces reliability.

\subsubsection{G10 --- Phase 5: ``Fire'' as Control Bridge}\label{g10-phase-5-fire-as-control-bridge}

\textbf{Hypothesis:} ``Fire'' would reliably bridge spark--ash as a non-polysemous control.

\textbf{What happened:} Near-zero frequency across all models. Fire is too central --- both endpoints already imply it.

\textbf{Resolution:} Established ``too-central'' as a new bridge failure mode (O6), distinct from off-axis (G5) and fragmentation.

\textbf{Lesson:} When both endpoints directly imply the candidate bridge, naming it adds no navigational information.

\subsubsection{G11 --- Phase 5: W-Shape Aggregate Prediction}\label{g11-phase-5-w-shape-aggregate-prediction}

\textbf{Hypothesis:} Bridge-present pairs show higher W-shape contrast than bridge-absent pairs in aggregate.

\textbf{What happened:} Bridge-present mean 0.0027, bridge-absent 0.0050. No aggregate signal.

\textbf{Resolution:} Effect is model-dependent and position-variable. Phase 6C redesigned the test as peak-detection.

\textbf{Lesson:} Aggregation can kill real signals when the effect is interaction-dependent.

\subsubsection{G12 --- Phase 6: Forced-Crossing Reduces Asymmetry}\label{g12-phase-6-forced-crossing-reduces-asymmetry}

\textbf{Hypothesis:} Forced crossings would reduce asymmetry below the 0.811 baseline to 0.50--0.70.

\textbf{What happened:} Forced-crossing asymmetry 0.817, indistinguishable from baseline.

\textbf{Resolution:} Asymmetry is a global structural property, not modulated by local bottlenecks.

\textbf{Lesson:} A shared waypoint does not produce shared paths.

\subsubsection{G13 --- Phase 6: Semantic Distance Predicts Bridge Position}\label{g13-phase-6-semantic-distance-predicts-bridge-position}

\textbf{Hypothesis:} Bridge position correlates with semantic distance ratio d(A,bridge)/d(A,C).

\textbf{What happened:} r = 0.239, p = 0.486. Irrelevant because bridges overwhelmingly anchor at positions 1--2.

\textbf{Resolution:} Position is determined by heading selection, not geometry.

\textbf{Lesson:} Navigational mechanics override geometric intuition.

\subsubsection{G14 --- Phase 6: Forced-Crossing Positional Stability}\label{g14-phase-6-forced-crossing-positional-stability}

\textbf{Hypothesis:} Forced-crossing bridges have low positional variance (SD \textless{} 0.8).

\textbf{What happened:} SD = 1.71 vs 0.52 for non-forced. Forced crossings are the most positionally unstable.

\textbf{Resolution:} Frame-crossing timing is model-dependent. Obligatory bridges are mandatory but not positionally fixed.

\textbf{Lesson:} Obligatory \(\neq\) stable.

\subsubsection{G15 --- Phase 7: Directional-Heading Theory (Strong Form)}\label{g15-phase-7-directional-heading-theory-strong-form}

\textbf{Hypothesis:} Incongruent pre-fills would produce more displacement than congruent pre-fills.

\textbf{What happened:} Displacement magnitudes indistinguishable: incongruent 0.515, congruent 0.436, neutral 0.536. CIs overlap.

\textbf{Resolution:} The mechanism is primarily associative primacy (whatever occupies position 1 anchors the trajectory) with secondary content modulation.

\textbf{Lesson:} Position matters more than content for displacement magnitude.

\subsubsection{G16 --- Phase 7: Polysemous Curvature}\label{g16-phase-7-polysemous-curvature}

\textbf{Hypothesis:} Polysemous-vertex triangles show higher curvature (triangle excess) than non-polysemous triangles.

\textbf{What happened:} Difference 0.053, CI includes zero. Categories indistinguishable.

\textbf{Resolution:} Polysemy affects \emph{which} concepts appear, not \emph{how far apart} they are. Also, the cross-model distance metric failed validity (G17).

\textbf{Lesson:} Always validate the instrument before interpreting measurements.

\subsubsection{G17 --- Phase 7: Cross-Model Distance Validity}\label{g17-phase-7-cross-model-distance-validity}

\textbf{Hypothesis:} Cross-model navigational distance correlation would exceed r = 0.50.

\textbf{What happened:} r = 0.170. Models assign wildly different distances to the same pairs. Some anti-correlate (Grok--Gemini r = --0.580).

\textbf{Resolution:} Navigational distances are fundamentally model-dependent. Blocks cross-model curvature comparison.

\textbf{Lesson:} Gait differences (R1) contaminate all attempts at shared distance metrics.

\subsubsection{G18 --- Phase 7: Too-Central as Binary Category}\label{g18-phase-7-too-central-as-binary-category}

\textbf{Hypothesis:} ``Too-central'' bridges (fire, tree, dough) would show frequency \textless{} 0.15, cleanly separated from ``obvious-useful'' bridges.

\textbf{What happened:} Too-central mean 0.496. ``Tree'' has frequency 1.000 for 3/4 models.

\textbf{Resolution:} Only fire/water meet the strict operational definition. The original categorization was based on intuition.

\textbf{Lesson:} Human intuition about ``too obvious'' does not match model navigation.

\subsubsection{G19 --- Phase 7: Forced-Crossing Resists Pre-Fill}\label{g19-phase-7-forced-crossing-resists-pre-fill}

\textbf{Hypothesis:} ``Bank'' on loan--shore would survive pre-filling at \textgreater{} 0.90 frequency.

\textbf{What happened:} Survival 0.267, positional shift 0.791.

\textbf{Resolution:} Forced crossings are obligatory only under unconstrained navigation. Pre-filling displaces even ``obligatory'' bottlenecks.

\textbf{Lesson:} Obligatory \(\neq\) robust under perturbation.

\subsubsection{G20 --- Phase 8: Route Exclusivity (Competitor Count)}\label{g20-phase-8-route-exclusivity-competitor-count}

\textbf{Hypothesis:} Bridges with more competitors show lower pre-fill survival (\(\rho\) \textless{} --0.70).

\textbf{What happened:} \(\rho\) = 0.116 (opposite direction). Sadness has 8 competitors yet highest survival.

\textbf{Resolution:} Competitor count ignores the distribution of traffic. A bridge with 8 weak competitors is more robust than one with 3 strong competitors.

\textbf{Lesson:} Quantity \(\neq\) quality in the competitive landscape.

\subsubsection{G21 --- Phase 8: Gemini Gradient Blindness}\label{g21-phase-8-gemini-gradient-blindness}

\textbf{Hypothesis:} Gemini selectively fails on gradient-midpoint bridges (interaction \(\geq\) 0.20).

\textbf{What happened:} Interaction 0.046 (null). Gemini's zeros concentrate on causal-chain pairs (6/10), not gradient pairs (1/10).

\textbf{Resolution:} Proposed transformation-chain blindness (G24).

\textbf{Lesson:} When the result is opposite to prediction, the categorization itself may be wrong.

\subsubsection{G22 --- Phase 8: Gait Normalization Rescues Distance}\label{g22-phase-8-gait-normalization-rescues-distance}

\textbf{Hypothesis:} Normalizing by model-specific baselines would raise cross-model distance correlation to r \textgreater{} 0.50.

\textbf{What happened:} Raw r = 0.212, normalized r = 0.212. Zero improvement. Pearson correlation is invariant to scalar division.

\textbf{Resolution:} Cross-model disagreement is ordinal, not scalar. Model-independent geometry is blocked.

\textbf{Lesson:} Always check whether a proposed fix is mathematically capable of addressing the pattern.

\subsubsection{G23 --- Phase 9: Dominance Ratio Predicts Fragility}\label{g23-phase-9-dominance-ratio-predicts-fragility}

\textbf{Hypothesis:} Bridge-to-competitor frequency ratio predicts survival (\(\rho\) \textgreater{} 0.50).

\textbf{What happened:} \(\rho\) = 0.157. Warm (ratio 1.00) is destroyed; fermentation (ratio 1.07) is bulletproof.

\textbf{Resolution:} Pre-fill content matters more than dominance ratio. Both single-variable structural predictors (G20, G23) have failed.

\textbf{Lesson:} Retrospective signals on small samples (N = 8) do not generalize.

\subsubsection{G24 --- Phase 9: Gemini Transformation-Chain Blindness}\label{g24-phase-9-gemini-transformation-chain-blindness}

\textbf{Hypothesis:} Gemini shows a transformation deficit (mean \textless{} 0.30).

\textbf{What happened:} Gemini transformation mean 0.667 (vs predicted \textless{} 0.30). Result reversed from prediction.

\textbf{Resolution:} Third type-based Gemini characterization fails. Whatever pair type is designated as ``hard,'' Gemini does relatively better on it.

\textbf{Lesson:} When a hypothesis fails backward in two successive tests, the classification itself is wrong. Abandon the type-based approach.

\subsubsection{G25 (partial) --- Phase 9: Pre-Fill Crossover Regression}\label{g25-partial-phase-9-pre-fill-crossover-regression}

\textbf{Hypothesis:} Significant negative slope predicting survival from unconstrained frequency; crossover at 0.40--0.50.

\textbf{What happened:} Slope --3.355 (correct direction) but CI includes zero. Crossover at 0.790 (far from predicted).

\textbf{Resolution:} Qualitative threshold survives (marginal facilitation, dominant displacement). Quantitative regression fails.

\textbf{Lesson:} Qualitative predictions succeed; quantitative predictions fail. The benchmark's defining pattern.

\subsubsection{G26 --- Phase 10: Bridge Generalization Fails (RESURRECTED)}\label{g26-phase-10-bridge-generalization-fails-resurrected}

\textbf{Hypothesis:} Bridge frequency CI would include zero for new models.

\textbf{What happened initially:} With Llama 8B only (1 model), CI excluded zero. Entered graveyard.

\textbf{What happened after:} Timeout relaxation (60s \(\to\) 300s) recovered 3 additional models. With 4-model cohort, CI includes zero. Bridge structure generalizes.

\textbf{Resolution:} The original finding was an infrastructure artifact. Llama 8B's low bridge frequency (0.200) is a scale effect, not a generalization failure.

\textbf{Lesson:} Infrastructure limitations can masquerade as substantive findings. The sole resurrection in the graveyard.

\subsubsection{G27 --- Phase 10: Relation Class Ordering}\label{g27-phase-10-relation-class-ordering}

\textbf{Hypothesis:} On-axis \textless{} unrelated \textless{} same-domain in \(\geq\) 5/8 pairs.

\textbf{What happened:} Unrelated (0.388) \textless{} on-axis (0.643) \(\approx\) same-domain (0.708). Only 1/8 pairs showed predicted ordering.

\textbf{Resolution:} The distinction is binary (related vs.~unrelated), not three-way. Related pre-fills maintain navigational context.

\textbf{Lesson:} Coarse categorical distinctions capture phenomena; fine-grained orderings fail.

\subsubsection{G28 --- Phase 11: Control Pair Revision}\label{g28-phase-11-control-pair-revision}

\textbf{Hypothesis:} At least 2/4 new control candidates would pass screening.

\textbf{What happened:} 0/24 cells pass. Top frequencies 0.60--1.00. Retrospective: stapler--monsoon fails for all 12 models.

\textbf{Resolution:} R5 single-pair strict-threshold validation is fundamentally inadequate. LLMs navigate between any two concepts.

\textbf{Lesson:} LLMs have richer associative connectivity than human intuition predicts.

\subsubsection{G29 --- Phase 11: Asymmetry Universality Across Protocol}\label{g29-phase-11-asymmetry-universality-across-protocol}

\textbf{Hypothesis:} Mean asymmetry \textgreater{} 0.60 under all protocol conditions.

\textbf{What happened:} 5-waypoint conditions: 0.594 (below threshold). 9-waypoint conditions: 0.669--0.684 (above).

\textbf{Resolution:} Asymmetry is resolution-dependent. R2 qualified: robust at \(\geq\) 7 waypoints.

\textbf{Lesson:} Path length affects measurement sensitivity. Threshold calibration is resolution-dependent.
            % Appendix C: Dead Ends and Failed Hypotheses
\section{Prompt Templates}\label{appendix-d-prompt-templates}

This appendix documents the exact prompt templates used for waypoint elicitation across the benchmark's 11 phases.

\subsection{D.1 Standard Elicitation Prompts}\label{d.1-standard-elicitation-prompts}

Two prompt formats were developed and evaluated during Phase 1 prompt selection. The semantic format was selected as the primary template based on higher extraction accuracy and more interpretable waypoints on the holdout set.

\subsubsection{Semantic Format (Primary)}\label{semantic-format-primary}

\begin{verbatim}
Imagine walking through conceptual space from "{A}" to "{B}".
What {N} landmarks do you pass along the way?
List only the landmark concepts, one per line, numbered 1 through {N}.
Do not include the starting or ending concepts.
\end{verbatim}

\subsubsection{Direct Format (Fallback)}\label{direct-format-fallback}

\begin{verbatim}
List exactly {N} intermediate concepts that form a path from "{A}" to "{B}".
Respond with only the concepts, one per line, numbered 1 through {N}.
Do not include the starting or ending concepts.
\end{verbatim}

The direct format was used as a fallback for models that produced unparseable outputs under the semantic prompt, notably during Phase 10A reliability probing of new models.

\textbf{Default parameters.} Unless otherwise noted: N = 7 waypoints, temperature = 0.7, 10--20 independent runs per (model, pair, direction) condition. Phases 1--2 used 20 runs; later phases standardized to 10 runs.

\subsection{D.2 Pre-Fill Elicitation Prompt}\label{d.2-pre-fill-elicitation-prompt}

Introduced in Phase 7A for causal intervention experiments. The first waypoint is pre-specified, and the model generates the remaining N-1 waypoints.

\begin{verbatim}
The first intermediate concept between "{A}" and "{B}" is "{pre-fill}".
List exactly {N-1} more intermediate concepts that continue the path
from "{pre-fill}" to "{B}".
Respond with only the concepts, one per line, numbered 1 through {N-1}.
Do not include "{A}", "{pre-fill}", or "{B}" in your list.
\end{verbatim}

With default parameters, N-1 = 6 (the pre-fill occupies position 1, and the model generates positions 2--7).

\subsection{D.3 System Prompt}\label{d.3-system-prompt}

All elicitations used the following system prompt:

\begin{verbatim}
You are a helpful assistant that provides single-word or short-phrase
responses representing conceptual landmarks.
\end{verbatim}

This system prompt was consistent across all models and phases. No chain-of-thought or reasoning instructions were included; the waypoint task is presented as a direct generation task.

\subsection{D.4 API Configuration}\label{d.4-api-configuration}

All models were accessed via the OpenRouter API (\texttt{https://openrouter.ai/api/v1/chat/completions}) with the following parameters:

\begin{itemize}
\tightlist
\item
  \texttt{max\_tokens}: 256
\item
  \texttt{temperature}: 0.7 (default; varied in Phase 11C robustness analysis)
\item
  \texttt{top\_p}: 1.0
\item
  \texttt{frequency\_penalty}: 0
\item
  \texttt{presence\_penalty}: 0
\end{itemize}

Phase 10A introduced a \texttt{-\/-patient} CLI flag that extended the request timeout from 60 seconds to 300 seconds, recovering 3 of 5 new models that failed the default timeout on OpenRouter latency.

\subsection{D.5 Response Parsing}\label{d.5-response-parsing}

Model responses were parsed using a multi-strategy extraction pipeline attempting, in order:

\begin{enumerate}
\def\labelenumi{\arabic{enumi}.}
\tightlist
\item
  JSON array extraction
\item
  Numbered list parsing (e.g., ``1. concept'')
\item
  Bullet list parsing (e.g., ``- concept'')
\item
  Comma-separated value extraction
\item
  Arrow-separated extraction (recognizing \texttt{-\textgreater{}}, \texttt{=\textgreater{}}, \texttt{-\/-\textgreater{}}, \texttt{==\textgreater{}})
\item
  Line-by-line fallback
\end{enumerate}

After parsing, responses passed through the canonicalization pipeline (Section 2.3 of the main text): lowercase, article stripping, lemmatization via the compromise NLP library, and phrase normalization.

\subsection{D.6 Prompt Variations Across Phases}\label{d.6-prompt-variations-across-phases}

\begin{longtable}[]{@{}
  >{\raggedright\arraybackslash}p{(\columnwidth - 4\tabcolsep) * \real{0.2692}}
  >{\raggedright\arraybackslash}p{(\columnwidth - 4\tabcolsep) * \real{0.4231}}
  >{\raggedright\arraybackslash}p{(\columnwidth - 4\tabcolsep) * \real{0.3077}}@{}}
\toprule\noalign{}
\begin{minipage}[b]{\linewidth}\raggedright
Phase
\end{minipage} & \begin{minipage}[b]{\linewidth}\raggedright
Variation
\end{minipage} & \begin{minipage}[b]{\linewidth}\raggedright
Reason
\end{minipage} \\
\midrule\noalign{}
\endhead
\bottomrule\noalign{}
\endlastfoot
1 & Both formats tested on holdout pairs & Prompt format selection \\
1--6 & Semantic format only & Standard elicitation \\
7A--10B & Pre-fill variant added & Causal intervention \\
10A & Direct format fallback for some models & Parse reliability \\
11C & N varied (5, 7, 9); temperature varied (0.5, 0.7, 0.9) & Robustness analysis \\
\end{longtable}

No other prompt variations were introduced. The consistency of the prompt template across 11 phases is a deliberate design choice supporting cross-phase comparability.
     % Appendix D: Prompt Templates
\section{Statistical Methods}\label{appendix-e-statistical-methods}

This appendix documents the statistical procedures used throughout the benchmark. All analyses were implemented in TypeScript (Bun runtime) unless otherwise noted.

\subsection{E.1 Bootstrap Confidence Intervals}\label{e.1-bootstrap-confidence-intervals}

Confidence intervals for Jaccard-based metrics (gait consistency, path asymmetry, bridge frequency, transitivity) are 95\% percentile bootstrap intervals computed from 1,000 resamples with a seeded PRNG (seed = 42) for reproducibility.

\textbf{Procedure.} For a metric computed over \emph{n} independent runs:

\begin{enumerate}
\def\labelenumi{\arabic{enumi}.}
\tightlist
\item
  Resample \emph{n} runs with replacement (1,000 times).
\item
  Recompute the metric on each resample.
\item
  Report the 2.5th and 97.5th percentiles as the 95\% CI.
\end{enumerate}

\textbf{Limitation.} \citet{efron1993bootstrap} recommend 5,000--10,000 resamples for publication-quality confidence intervals. The benchmark uses 1,000 for computational efficiency during the exploration-first workflow. CIs should be interpreted as approximate; the qualitative conclusions (whether a CI includes zero) are robust to resample count, but exact CI bounds may shift slightly with more resamples.

\subsection{E.2 Permutation Tests for Asymmetry}\label{e.2-permutation-tests-for-asymmetry}

Statistical significance of path asymmetry is assessed via permutation test. For a given (model, pair) combination with \emph{n\_f} forward runs and \emph{n\_r} reverse runs:

\begin{enumerate}
\def\labelenumi{\arabic{enumi}.}
\tightlist
\item
  Pool all \emph{n\_f + n\_r} runs.
\item
  Randomly assign \emph{n\_f} to ``forward'' and \emph{n\_r} to ``reverse'' (1,000 permutations).
\item
  Compute the mean cross-direction Jaccard for each permutation.
\item
  The \emph{p}-value is the fraction of permutations producing a cross-direction Jaccard as low or lower than the observed value, using the \citet{phipson2010permutation} correction to avoid exact-zero \emph{p}-values:
\end{enumerate}

\[p = \frac{b + 1}{m + 1}\]

where \emph{b} is the count of permutations with test statistic as extreme as observed and \emph{m} is the total number of permutations.

\textbf{Limitation.} With 1,000 permutations, the minimum achievable \emph{p}-value is approximately 0.001 (after the Phipson--Smyth correction). This limits resolution for Bonferroni-corrected thresholds. The 87\% significance rate reported in Section 5.1 (73 of 84 combinations at p \textless{} 0.05) was computed without multiple-comparison correction. With Bonferroni correction (threshold 0.05/84 \(\approx\) 0.0006), the percentage would be lower, and the permutation test's resolution is insufficient to distinguish near this corrected threshold. The qualitative conclusion --- that the vast majority of combinations show genuine asymmetry --- is robust, but the specific percentage should be interpreted conservatively.

\subsection{E.3 Friedman Test and Post-Hoc Comparisons}\label{e.3-friedman-test-and-post-hoc-comparisons}

The effect of pre-fill relation class on bridge survival (Section 6.3) was tested using the Friedman test, a non-parametric repeated-measures alternative to one-way ANOVA. The 8 concept pairs serve as blocks, and the 3 relation classes (on-axis, same-domain, unrelated) serve as treatments.

\textbf{Procedure:}

\begin{enumerate}
\def\labelenumi{\arabic{enumi}.}
\tightlist
\item
  For each pair, rank the 3 survival rates.
\item
  Compute the Friedman chi-squared statistic from the rank sums.
\item
  Assess significance against the chi-squared distribution with \emph{k - 1} = 2 degrees of freedom.
\end{enumerate}

\textbf{Result:} \(\chi^2\) = 6.750, \emph{df} = 2, \emph{p} = 0.034.

\textbf{Post-hoc comparisons.} Pairwise Wilcoxon signed-rank tests:

\begin{longtable}[]{@{}lll@{}}
\toprule\noalign{}
Comparison & \emph{p}-value & Survives Bonferroni (0.05/3 = 0.017)? \\
\midrule\noalign{}
\endhead
\bottomrule\noalign{}
\endlastfoot
On-axis vs.~unrelated & 0.025 & No \\
Same-domain vs.~unrelated & 0.036 & No \\
On-axis vs.~same-domain & 0.889 & N/A \\
\end{longtable}

Neither pairwise comparison survives Bonferroni correction for 3 comparisons. The qualitative conclusion --- unrelated pre-fills are more disruptive than related ones --- is supported by the omnibus Friedman test but the specific pairwise \emph{p}-values are suggestive rather than confirmatory.

\subsection{E.4 ANOVA for Multiverse Robustness}\label{e.4-anova-for-multiverse-robustness}

The Phase 11C robustness analysis used a factorial ANOVA decomposing gait variance across model identity, waypoint count, and temperature.

\textbf{Design:} 3 models (Claude, GPT, DeepSeek) \(\times\) 2 waypoint counts (5, 9) \(\times\) 2 temperatures (0.5, 0.9), plus the 7-waypoint / 0.7-temperature baseline assembled from prior-phase data.

\textbf{Procedure:}

\begin{enumerate}
\def\labelenumi{\arabic{enumi}.}
\tightlist
\item
  Compute cell means: mean Jaccard for each (model, waypoint count, temperature, pair) combination.
\item
  Compute sum of squares for each factor and interaction.
\item
  Express as eta-squared (\(\eta^2\) = SS\(_{\text{factor}}\) / SS\(_{\text{total}}\)).
\item
  Approximate \emph{p}-values using the chi-squared distribution (ignoring denominator degrees of freedom).
\end{enumerate}

\textbf{Results:}

\begin{longtable}[]{@{}lll@{}}
\toprule\noalign{}
Factor & \(\eta^2\) & Approx. \emph{p} \\
\midrule\noalign{}
\endhead
\bottomrule\noalign{}
\endlastfoot
Model identity & 0.242 & \(\approx\) 0.001 \\
Waypoint count & 0.008 & \(\approx\) 0.520 \\
Temperature & 0.002 & \(\approx\) 0.743 \\
Waypoint \(\times\) Temperature & \(\approx\) 0.000 & \(\approx\) 0.886 \\
\end{longtable}

\textbf{Limitations:}

\begin{enumerate}
\def\labelenumi{\arabic{enumi}.}
\tightlist
\item
  \emph{P}-value approximation.* The chi-squared approximation ignores denominator degrees of freedom. Exact F-test \emph{p}-values would require a proper mean squares decomposition. The eta-squared effect sizes are more reliable than the exact \emph{p}-values.
\item
  \emph{Independence assumption.} Per-pair cells are treated as independent observations without repeated-measures modeling. A mixed-effects model with pair as a random effect would be more appropriate for confirmatory inference.
\item
  \emph{Coverage.} Only 3 models were tested. The conclusion that model identity dominates protocol variation is supported within this subset but should be extended with caution.
\item
  \emph{Baseline sparsity.} DeepSeek has no prior-phase data; its 7wp-t0.7 baseline cells are empty. This affects the baseline condition but not the four new conditions.
\end{enumerate}

The analysis should be understood as descriptive/exploratory rather than confirmatory. The qualitative conclusion --- that model identity accounts for far more variance than protocol parameters --- is robust to these approximations, as the effect size gap (0.242 vs.~0.008 and 0.002) is large.

\subsection{E.5 Chi-Squared Goodness-of-Fit Test}\label{e.5-chi-squared-goodness-of-fit-test}

Departure from uniform waypoint distributions (Section 5.6) is assessed via chi-squared goodness-of-fit test.

\textbf{Procedure.} For each (model, pair) condition:

\begin{enumerate}
\def\labelenumi{\arabic{enumi}.}
\tightlist
\item
  Count unique waypoints across all runs, treating each waypoint as appearing once per run.
\item
  Compute expected frequencies under the null hypothesis that all waypoints are equally likely.
\item
  Compute the chi-squared statistic: \(\chi^2 = \sum (O_i - E_i)^2 / E_i\).
\item
  Compute \emph{p}-value using the Wilson--Hilferty normal approximation.
\item
  Apply Bonferroni correction at the pair level (dividing the significance threshold by the number of pairs tested).
\end{enumerate}

\textbf{Note on naming.} The codebase function implementing this test is named \texttt{ksTestUniform} (after Kolmogorov--Smirnov), but the implementation is chi-squared goodness-of-fit, which is more appropriate for categorical frequency data. The paper uses the correct name (chi-squared).

\subsection{E.6 Kendall's W (Coefficient of Concordance)}\label{e.6-kendalls-w-coefficient-of-concordance}

Gait rank stability across protocol conditions (Section 9.3) is assessed via Kendall's W:

\[W = \frac{12 \sum_{i=1}^{n} (R_i - \bar{R})^2}{k^2(n^3 - n)}\]

where \(R_i\) is the rank sum for model \emph{i} across \emph{k} conditions, and \emph{n} is the number of models. W = 1 indicates perfect agreement across conditions; W = 0 indicates no agreement. The observed W = 0.840 indicates strong but imperfect rank stability.

\subsection{E.7 Bridge Frequency Comparison (Confidence Interval on Difference)}\label{e.7-bridge-frequency-comparison-confidence-interval-on-difference}

The structure/content/scale hierarchy (Section 8.2) compares aggregate bridge frequency between the original cohort and expanded cohorts using bootstrap CIs on the difference:

\begin{enumerate}
\def\labelenumi{\arabic{enumi}.}
\tightlist
\item
  For each cohort, compute mean bridge frequency across all (model, pair) conditions.
\item
  Bootstrap the difference (original mean \(-\) new mean) with 1,000 resamples.
\item
  If the 95\% CI includes zero, the cohorts are statistically indistinguishable.
\end{enumerate}

Phase 10A: diff = \(-0.096\), CI {[}\(-0.241\), 0.064{]} (includes zero). Phase 11A combined: diff = \(-0.100\), CI {[}\(-0.286\), 0.089{]} (includes zero).

\subsection{E.8 Canonicalization and Its Statistical Implications}\label{e.8-canonicalization-and-its-statistical-implications}

All metrics operate on canonicalized waypoint tokens. The canonicalization pipeline (Section 2.3) handles morphological variants but not synonyms. This has systematic statistical implications:

\begin{itemize}
\tightlist
\item
  \textbf{Jaccard similarity} is computed on token identity after canonicalization. Semantically equivalent but lexically distinct waypoints (``melody'' vs.~``tune,'' ``cemetery'' vs.~``graveyard'') are treated as completely different. All Jaccard-based metrics therefore \emph{underestimate} true conceptual overlap.
\item
  \textbf{Gait consistency} is a lower bound on true navigational consistency.
\item
  \textbf{Path asymmetry} is an upper bound on true conceptual asymmetry.
\item
  \textbf{Bridge frequency} may miss bridges expressed through synonyms rather than exact (or morphologically related) tokens.
\end{itemize}

The magnitude of this bias is bounded by the canonicalization (which collapses morphological variants) but is otherwise unknown. Embedding-based semantic similarity would help bound this effect but was deferred due to API cost.

\subsection{E.9 Multiple Comparison Corrections}\label{e.9-multiple-comparison-corrections}

The benchmark applies Bonferroni correction in two contexts:

\begin{enumerate}
\def\labelenumi{\arabic{enumi}.}
\tightlist
\item
  \textbf{Uniformity tests} (Section 5.6): 8 pairs tested per model; threshold adjusted to 0.05/8 = 0.00625.
\item
  \textbf{Asymmetry significance} (Section 5.1): 84 pair/model combinations; the 87\% significance rate is reported without correction. With Bonferroni (0.05/84 \(\approx\) 0.0006), the percentage would be lower.
\end{enumerate}

No correction is applied to the phase-level prediction accuracy tracking (Sections 7.1, 7.2), which is descriptive rather than inferential.

For the Friedman post-hoc comparisons (Section 6.3, Appendix E.3), Bonferroni correction for 3 pairwise comparisons (threshold 0.017) is documented but results in neither surviving comparison reaching significance. The paper reports both corrected and uncorrected results for transparency.
  % Appendix E: Statistical Methods
\section{Robustness Analysis Details}\label{appendix-f-robustness-analysis-details}

This appendix provides the full per-model, per-pair, per-condition breakdown from the Phase 11C multiverse robustness experiment. Section 9 of the main text presents aggregated results; this appendix provides the underlying data.

\subsection{F.1 Experimental Design}\label{f.1-experimental-design}

\textbf{Models:} Claude Sonnet 4.6, GPT-5.2, DeepSeek V3.2

\textbf{Pairs (forward + reverse):} - light \(\to\) color / color \(\to\) light (bridge: spectrum) - hot \(\to\) cold / cold \(\to\) hot (bridge: warm) - emotion \(\to\) melancholy / melancholy \(\to\) emotion (bridge: sadness)

\textbf{Conditions (2\$\times\$2 grid + baseline):}

\begin{longtable}[]{@{}lll@{}}
\toprule\noalign{}
Condition & Waypoints & Temperature \\
\midrule\noalign{}
\endhead
\bottomrule\noalign{}
\endlastfoot
Baseline (7wp-t0.7) & 7 & 0.7 \\
5wp-t0.5 & 5 & 0.5 \\
5wp-t0.9 & 5 & 0.9 \\
9wp-t0.5 & 9 & 0.5 \\
9wp-t0.9 & 9 & 0.9 \\
\end{longtable}

\textbf{Runs:} 1,080 new (4 conditions \(\times\) 3 models \(\times\) 6 pairs \(\times\) 15 reps) + \$\sim\$340 reused baseline runs from prior phases.

\textbf{Baseline data limitation:} The 7wp-t0.7 baseline is assembled from prior-phase data (Phases 1, 2, 6A). DeepSeek, added in Phase 11A, has no prior-phase data --- its baseline cells are empty. Additionally, reverse-direction data for some pairs is missing from prior phases. Baseline values for DeepSeek should be treated as missing, not zero.

\begin{center}\rule{0.5\linewidth}{0.5pt}\end{center}

\subsection{F.2 Gait (Intra-Model Jaccard) by Model and Condition}\label{f.2-gait-intra-model-jaccard-by-model-and-condition}

\begin{longtable}[]{@{}llllll@{}}
\toprule\noalign{}
Model & 7wp-t0.7 & 5wp-t0.5 & 5wp-t0.9 & 9wp-t0.5 & 9wp-t0.9 \\
\midrule\noalign{}
\endhead
\bottomrule\noalign{}
\endlastfoot
Claude & 0.676 & 0.794 & 0.772 & 0.699 & 0.624 \\
GPT & 0.378 & 0.413 & 0.477 & 0.382 & 0.365 \\
DeepSeek & 0.000* & 0.582 & 0.492 & 0.592 & 0.597 \\
\end{longtable}

*DeepSeek baseline is empty (no prior-phase data); reported as 0.000 by the analysis script.

\textbf{Observations:} - Claude shows the highest gait in every condition (0.624--0.794). - GPT shows the lowest gait in all new conditions (0.365--0.477). - DeepSeek occupies the middle position in all new conditions (0.492--0.597). - 5wp-t0.5 produces the highest gait for both Claude (0.794) and DeepSeek (0.582) --- fewer waypoints and lower temperature reduce navigational variability. - 9wp-t0.9 produces the lowest gait for Claude (0.624) --- more waypoints and higher temperature maximize variability.

\subsubsection{Gait Rank-Order Stability}\label{gait-rank-order-stability}

\begin{longtable}[]{@{}
  >{\raggedright\arraybackslash}p{(\columnwidth - 6\tabcolsep) * \real{0.3143}}
  >{\raggedright\arraybackslash}p{(\columnwidth - 6\tabcolsep) * \real{0.2286}}
  >{\raggedright\arraybackslash}p{(\columnwidth - 6\tabcolsep) * \real{0.2286}}
  >{\raggedright\arraybackslash}p{(\columnwidth - 6\tabcolsep) * \real{0.2286}}@{}}
\toprule\noalign{}
\begin{minipage}[b]{\linewidth}\raggedright
Condition
\end{minipage} & \begin{minipage}[b]{\linewidth}\raggedright
Rank 1
\end{minipage} & \begin{minipage}[b]{\linewidth}\raggedright
Rank 2
\end{minipage} & \begin{minipage}[b]{\linewidth}\raggedright
Rank 3
\end{minipage} \\
\midrule\noalign{}
\endhead
\bottomrule\noalign{}
\endlastfoot
7wp-t0.7 (baseline) & Claude (0.676) & GPT (0.378) & DeepSeek (0.000*) \\
5wp-t0.5 & Claude (0.794) & DeepSeek (0.582) & GPT (0.413) \\
5wp-t0.9 & Claude (0.772) & DeepSeek (0.492) & GPT (0.477) \\
9wp-t0.5 & Claude (0.699) & DeepSeek (0.592) & GPT (0.382) \\
9wp-t0.9 & Claude (0.624) & DeepSeek (0.597) & GPT (0.365) \\
\end{longtable}

Kendall's W = 0.840. Claude is always first. The baseline anomaly (GPT \textgreater{} DeepSeek) reflects DeepSeek's missing baseline data, not genuinely lower consistency.

\begin{center}\rule{0.5\linewidth}{0.5pt}\end{center}

\subsection{F.3 Asymmetry by Model, Pair, and Condition}\label{f.3-asymmetry-by-model-pair-and-condition}

\subsubsection{Full Breakdown}\label{full-breakdown}

\begin{longtable}[]{@{}lllll@{}}
\toprule\noalign{}
Model & Condition & Pair & Asymmetry & 95\% CI \\
\midrule\noalign{}
\endhead
\bottomrule\noalign{}
\endlastfoot
Claude & 7wp-t0.7 & hot--cold & 0.726 & {[}0.697, 0.750{]} \\
Claude & 7wp-t0.7 & emotion--melancholy & 1.000 & {[}1.000, 1.000{]} \\
GPT & 7wp-t0.7 & hot--cold & 0.471 & {[}0.414, 0.522{]} \\
GPT & 7wp-t0.7 & emotion--melancholy & 1.000 & {[}1.000, 1.000{]} \\
DeepSeek & 7wp-t0.7 & hot--cold & 1.000 & {[}1.000, 1.000{]} \\
DeepSeek & 7wp-t0.7 & emotion--melancholy & 1.000 & {[}1.000, 1.000{]} \\
Claude & 5wp-t0.5 & hot--cold & 0.750 & {[}0.750, 0.750{]} \\
Claude & 5wp-t0.5 & emotion--melancholy & 0.512 & {[}0.466, 0.550{]} \\
GPT & 5wp-t0.5 & hot--cold & 0.434 & {[}0.372, 0.503{]} \\
GPT & 5wp-t0.5 & emotion--melancholy & 0.520 & {[}0.461, 0.568{]} \\
DeepSeek & 5wp-t0.5 & hot--cold & 0.665 & {[}0.641, 0.689{]} \\
DeepSeek & 5wp-t0.5 & emotion--melancholy & 0.684 & {[}0.644, 0.719{]} \\
Claude & 5wp-t0.9 & hot--cold & 0.726 & {[}0.690, 0.750{]} \\
Claude & 5wp-t0.9 & emotion--melancholy & 0.519 & {[}0.473, 0.562{]} \\
GPT & 5wp-t0.9 & hot--cold & 0.392 & {[}0.354, 0.435{]} \\
GPT & 5wp-t0.9 & emotion--melancholy & 0.555 & {[}0.445, 0.654{]} \\
DeepSeek & 5wp-t0.9 & hot--cold & 0.653 & {[}0.621, 0.684{]} \\
DeepSeek & 5wp-t0.9 & emotion--melancholy & 0.712 & {[}0.667, 0.760{]} \\
Claude & 9wp-t0.5 & hot--cold & 0.671 & {[}0.647, 0.696{]} \\
Claude & 9wp-t0.5 & emotion--melancholy & 0.625 & {[}0.583, 0.662{]} \\
GPT & 9wp-t0.5 & hot--cold & 0.514 & {[}0.460, 0.562{]} \\
GPT & 9wp-t0.5 & emotion--melancholy & 0.682 & {[}0.645, 0.717{]} \\
DeepSeek & 9wp-t0.5 & hot--cold & 0.783 & {[}0.760, 0.804{]} \\
DeepSeek & 9wp-t0.5 & emotion--melancholy & 0.740 & {[}0.708, 0.766{]} \\
Claude & 9wp-t0.9 & hot--cold & 0.655 & {[}0.620, 0.690{]} \\
Claude & 9wp-t0.9 & emotion--melancholy & 0.612 & {[}0.565, 0.656{]} \\
GPT & 9wp-t0.9 & hot--cold & 0.596 & {[}0.555, 0.638{]} \\
GPT & 9wp-t0.9 & emotion--melancholy & 0.725 & {[}0.684, 0.769{]} \\
DeepSeek & 9wp-t0.9 & hot--cold & 0.760 & {[}0.732, 0.785{]} \\
DeepSeek & 9wp-t0.9 & emotion--melancholy & 0.754 & {[}0.720, 0.784{]} \\
\end{longtable}

\subsubsection{Condition Means}\label{condition-means}

\begin{longtable}[]{@{}llll@{}}
\toprule\noalign{}
Condition & Mean Asymmetry & 95\% CI & \(>\) 0.60? \\
\midrule\noalign{}
\endhead
\bottomrule\noalign{}
\endlastfoot
7wp-t0.7 (baseline)* & 0.599 & {[}0.471, 0.726{]} & no \\
5wp-t0.5 & 0.594 & {[}0.512, 0.675{]} & no \\
5wp-t0.9 & 0.593 & {[}0.502, 0.688{]} & no \\
9wp-t0.5 & 0.669 & {[}0.597, 0.733{]} & \textbf{yes} \\
9wp-t0.9 & 0.684 & {[}0.631, 0.736{]} & \textbf{yes} \\
\end{longtable}

*Baseline from 2 valid cells only (Claude and GPT on hot--cold).

\textbf{Key pattern:} The waypoint-count effect is consistent across all three models. At 9 waypoints, every model\(\times\)pair combination shows asymmetry \(\geq\) 0.514; at 5 waypoints, GPT hot--cold drops to 0.392--0.434. Temperature has negligible effect: within each waypoint count, 5wp-t0.5 (0.594) \(\approx\) 5wp-t0.9 (0.593) and 9wp-t0.5 (0.669) \(\approx\) 9wp-t0.9 (0.684).

\begin{center}\rule{0.5\linewidth}{0.5pt}\end{center}

\subsection{F.4 Bridge Frequency by Pair, Model, and Condition}\label{f.4-bridge-frequency-by-pair-model-and-condition}

\subsubsection{Full Breakdown}\label{full-breakdown-1}

\textbf{light \(\to\) color (bridge: spectrum)}

\begin{longtable}[]{@{}llllll@{}}
\toprule\noalign{}
Model & 7wp-t0.7 & 5wp-t0.5 & 5wp-t0.9 & 9wp-t0.5 & 9wp-t0.9 \\
\midrule\noalign{}
\endhead
\bottomrule\noalign{}
\endlastfoot
Claude & 1.000 & 1.000 & 1.000 & 1.000 & 1.000 \\
GPT & 0.975 & 1.000 & 1.000 & 1.000 & 1.000 \\
DeepSeek & 0.000* & 0.933 & 0.800 & 1.000 & 1.000 \\
\end{longtable}

\textbf{hot \(\to\) cold (bridge: warm)}

\begin{longtable}[]{@{}llllll@{}}
\toprule\noalign{}
Model & 7wp-t0.7 & 5wp-t0.5 & 5wp-t0.9 & 9wp-t0.5 & 9wp-t0.9 \\
\midrule\noalign{}
\endhead
\bottomrule\noalign{}
\endlastfoot
Claude & 1.000 & 1.000 & 1.000 & 1.000 & 1.000 \\
GPT & 1.000 & 1.000 & 1.000 & 1.000 & 1.000 \\
DeepSeek & 0.000* & 1.000 & 1.000 & 1.000 & 1.000 \\
\end{longtable}

\textbf{emotion \(\to\) melancholy (bridge: sadness)}

\begin{longtable}[]{@{}llllll@{}}
\toprule\noalign{}
Model & 7wp-t0.7 & 5wp-t0.5 & 5wp-t0.9 & 9wp-t0.5 & 9wp-t0.9 \\
\midrule\noalign{}
\endhead
\bottomrule\noalign{}
\endlastfoot
Claude & 1.000 & 1.000 & 1.000 & 0.933 & 1.000 \\
GPT & 0.975 & 1.000 & 1.000 & 1.000 & 0.867 \\
DeepSeek & 0.000* & 1.000 & 1.000 & 1.000 & 1.000 \\
\end{longtable}

*DeepSeek 7wp-t0.7 entries are 0.000 due to missing baseline data, not actual zero bridge frequency.

\subsubsection{Condition Means (Excluding DeepSeek Baseline)}\label{condition-means-excluding-deepseek-baseline}

\begin{longtable}[]{@{}lll@{}}
\toprule\noalign{}
Condition & Mean Bridge Freq & 95\% CI \\
\midrule\noalign{}
\endhead
\bottomrule\noalign{}
\endlastfoot
7wp-t0.7 (baseline) & 0.992 & {[}0.983, 1.000{]} \\
5wp-t0.5 & 0.993 & {[}0.978, 1.000{]} \\
5wp-t0.9 & 0.978 & {[}0.933, 1.000{]} \\
9wp-t0.5 & 0.993 & {[}0.978, 1.000{]} \\
9wp-t0.9 & 0.985 & {[}0.956, 1.000{]} \\
\end{longtable}

All conditions exceed 0.97. The two lowest individual cells are DeepSeek's spectrum at 5wp-t0.9 (0.800) and GPT's sadness at 9wp-t0.9 (0.867). Even these represent only minor dips from ceiling.

\begin{center}\rule{0.5\linewidth}{0.5pt}\end{center}

\subsection{F.5 ANOVA Interaction Test}\label{f.5-anova-interaction-test}

Factorial analysis testing gait (intra-model Jaccard) against model identity, waypoint count, temperature, and their interaction. Data from the four new conditions only (baseline excluded due to DeepSeek sparsity).

\begin{longtable}[]{@{}lrr@{}}
\toprule\noalign{}
Factor & \(\eta^2\) (SS Proportion) & p-value (approx.) \\
\midrule\noalign{}
\endhead
\bottomrule\noalign{}
\endlastfoot
Model identity & 0.242 & \(\approx\) 0.001 \\
Waypoint count & 0.008 & \(\approx\) 0.520 \\
Temperature & 0.002 & \(\approx\) 0.743 \\
Waypoint \(\times\) Temperature & \(\approx\) 0.000 & \(\approx\) 0.886 \\
\end{longtable}

\textbf{Methodological note:} p-values use a chi-squared approximation to the F-distribution (ignoring denominator degrees of freedom) and treat per-pair cells as independent without repeated-measures modeling. Only 3 models were tested. The effect sizes (\(\eta^2\)) are more reliable than the exact p-values. The qualitative conclusion --- that model identity dominates while protocol parameters are null --- is robust to the approximation, but the analysis should be understood as descriptive/exploratory rather than confirmatory.

\begin{center}\rule{0.5\linewidth}{0.5pt}\end{center}

\subsection{F.6 Waypoint Scaling}\label{f.6-waypoint-scaling}

Fraction of waypoints shared across different path lengths, computed as mean Jaccard similarity between the waypoint sets of matched pairs at different waypoint counts.

\begin{longtable}[]{@{}lc@{}}
\toprule\noalign{}
Comparison & Shared Fraction \\
\midrule\noalign{}
\endhead
\bottomrule\noalign{}
\endlastfoot
5wp vs 7wp & 0.761 \\
7wp vs 9wp & 0.504 \\
5wp vs 9wp & 0.481 \\
Reference: O10 (5wp vs 10wp, Phase 1) & 0.705 \\
\end{longtable}

The 5-to-7 scaling (0.761) is consistent with O10 (Phase 1's 5-to-10 at 0.705): shorter paths are genuine coarse-grainings of longer paths. The 7-to-9 scaling (0.504) is lower, suggesting that additional waypoints from 7 to 9 introduce genuinely new content rather than interpolating between existing waypoints.

\begin{center}\rule{0.5\linewidth}{0.5pt}\end{center}

\subsection{F.7 Control Pair Screening (Phase 11B)}\label{f.7-control-pair-screening-phase-11b}

\subsubsection{Candidate Results}\label{candidate-results}

4 candidates screened across 6 models (Claude, GPT, Grok, Gemini, DeepSeek, Mistral), 10 runs each. Frequency gate: top-waypoint frequency \(<\) 0.15. Entropy gate: waypoint entropy \(>\) 5.0.

\begin{longtable}[]{@{}
  >{\raggedright\arraybackslash}p{(\columnwidth - 8\tabcolsep) * \real{0.1692}}
  >{\centering\arraybackslash}p{(\columnwidth - 8\tabcolsep) * \real{0.2000}}
  >{\centering\arraybackslash}p{(\columnwidth - 8\tabcolsep) * \real{0.2000}}
  >{\centering\arraybackslash}p{(\columnwidth - 8\tabcolsep) * \real{0.2923}}
  >{\centering\arraybackslash}p{(\columnwidth - 8\tabcolsep) * \real{0.1385}}@{}}
\toprule\noalign{}
\begin{minipage}[b]{\linewidth}\raggedright
Candidate
\end{minipage} & \begin{minipage}[b]{\linewidth}\centering
Max Top Freq
\end{minipage} & \begin{minipage}[b]{\linewidth}\centering
Min Entropy
\end{minipage} & \begin{minipage}[b]{\linewidth}\centering
Best Model Entropy
\end{minipage} & \begin{minipage}[b]{\linewidth}\centering
Passes?
\end{minipage} \\
\midrule\noalign{}
\endhead
\bottomrule\noalign{}
\endlastfoot
accordion--stalactite & 1.000 & 3.29 & Gemini (4.92) & fail \\
turmeric--trigonometry & 1.000 & 3.44 & DeepSeek (4.77) & fail \\
barnacle--sonnet & 1.000 & 3.44 & Gemini (4.57) & fail \\
magnesium--ballet & 1.000 & 3.65 & Gemini (4.92) & fail \\
\end{longtable}

0 of 24 model\(\times\)candidate cells pass either gate.

\subsubsection{Retrospective: Stapler--Monsoon Across 12 Models}\label{retrospective-staplermonsoon-across-12-models}

\begin{longtable}[]{@{}lllc@{}}
\toprule\noalign{}
Model & Cohort & Top Waypoint & Top Frequency \\
\midrule\noalign{}
\endhead
\bottomrule\noalign{}
\endlastfoot
Claude & original & office & 0.775 \\
GPT & original & paperwork & 0.650 \\
Grok & original & flood & 0.700 \\
Gemini & original & humidity & 0.775 \\
MiniMax & phase10a & paper & 0.933 \\
Kimi & phase10a & paper & 0.933 \\
Qwen & phase10a & cloud & 0.867 \\
Llama 8B & phase10a & office & 0.800 \\
DeepSeek & phase11a & paperclip & 0.600 \\
Mistral & phase11a & workplace & 1.000 \\
Cohere & phase11a & office supply & 1.000 \\
Llama 4 & phase11a & paperweight & 1.000 \\
\end{longtable}

All 12 models fail R5 under strict criteria (top frequency \(<\) 0.15). The original 4 models show frequencies of 0.650--0.775; the Phase 10A and 11A models show 0.600--1.000. Under these criteria, stapler--monsoon has never been a valid control for any model.

\begin{center}\rule{0.5\linewidth}{0.5pt}\end{center}

\subsection{F.8 Phase 11C Predictions Summary}\label{f.8-phase-11c-predictions-summary}

\begin{longtable}[]{@{}
  >{\raggedright\arraybackslash}p{(\columnwidth - 8\tabcolsep) * \real{0.1026}}
  >{\raggedright\arraybackslash}p{(\columnwidth - 8\tabcolsep) * \real{0.2821}}
  >{\raggedright\arraybackslash}p{(\columnwidth - 8\tabcolsep) * \real{0.2051}}
  >{\raggedright\arraybackslash}p{(\columnwidth - 8\tabcolsep) * \real{0.1795}}
  >{\raggedright\arraybackslash}p{(\columnwidth - 8\tabcolsep) * \real{0.2308}}@{}}
\toprule\noalign{}
\begin{minipage}[b]{\linewidth}\raggedright
\#
\end{minipage} & \begin{minipage}[b]{\linewidth}\raggedright
Prediction
\end{minipage} & \begin{minipage}[b]{\linewidth}\raggedright
Result
\end{minipage} & \begin{minipage}[b]{\linewidth}\raggedright
Value
\end{minipage} & \begin{minipage}[b]{\linewidth}\raggedright
Verdict
\end{minipage} \\
\midrule\noalign{}
\endhead
\bottomrule\noalign{}
\endlastfoot
1 & Gait rank-order survives (W \(>\) 0.70) & W = 0.840 & above threshold & confirmed \\
2 & Asymmetry \(>\) 0.60 all conditions & 5wp: 0.594, 0.593 & below threshold & not confirmed \\
3 & Bridge freq survives (spectrum \(>\) 0.50, warm \(>\) 0.30) & all \(>\) 0.97 & far above threshold & confirmed \\
4 & Temperature 0.5 increases Jaccard by 0.05--0.15 & \(\Delta\) = 0.226 & exceeds range & not confirmed \\
5 & Temperature 0.9 decreases Jaccard, all \(>\) 0.10 & \(\Delta\) = 0.203 & magnitude wrong & not confirmed \\
6 & 5wp bridge freq \(>\) 9wp & 5wp = 0.985, 9wp = 0.989 & no difference & not confirmed \\
7 & Control pair unstructured all conditions & structured & --- & not confirmed \\
\end{longtable}
   % Appendix F: Robustness Details

\end{document}
